\documentclass{amsart}
%\documentclass{article}
\usepackage[russian]{babel}
%\usepackage[koi8-r]{inputenc}
\usepackage[utf8]{inputenc}
\usepackage[all]{xy}

\usepackage{fixmath}

\usepackage[usenames]{color}
\usepackage{colortbl}


\newtheorem{lemma}{Лемма}
\newtheorem{theorem}{Теорема}
\newtheorem{proposition}{Предложение}
\newtheorem{corollary}{Следствие}
\theoremstyle{definition}
\newtheorem{definition}{Определение}
\newtheorem{notation}{Обозначение}
\newtheorem{remark}{Замечание}
\newtheorem{example}{Пример}

\newcommand{\courbe}[7]{\qbezier(#1,#3)(#4,#7)(#5,#7)
\qbezier(#5,#7)(#6,#7)(#2,#3)}

\newcommand{\courbebis}[3]{\qbezier(#1,#3)(#1,2*#3)(#1/2+#2/2,2*#3)
\qbezier(#1/2+#2/2,2*#3)(#2,2*#3)(#2,#3)}

\newcommand{\NN}{\mathbb{N}^*}

\newcommand{\bs}[1]{\boldsymbol{#1}}

\begin{document}

\makeatletter
\edef\@@and{и}
\makeatother




%\let\uppercase\MakeUppercase

%{ \bf ДВОЙСТВЕННОСТЬ ОРБИТ В ИНД-МНОГООБРАЗИЯХ МАКСИМАЛЬНЫХ ОБОБЩЁННЫХ ФЛАГОВ}
%\title{Двойственность орбит в инд-многообразиях максимальных обобщённых флагов}
\title[\CYRD\CYRV\CYRO\CYRISHRT\CYRS\CYRT\CYRV\CYRE\CYRN\CYRN\CYRO\CYRS\CYRT\CYRSFTSN\
\CYRO\CYRR\CYRB\CYRI\CYRT\ \CYRV\
\CYRI\CYRN\CYRD-\CYRM\CYRN\CYRO\CYRG\CYRO\CYRO\CYRB\CYRR\CYRA\CYRZ\CYRI\CYRYA\CYRH...]
{\CYRD\CYRV\CYRO\CYRISHRT\CYRS\CYRT\CYRV\CYRE\CYRN\CYRN\CYRO\CYRS\CYRT\CYRSFTSN\
\CYRO\CYRR\CYRB\CYRI\CYRT\ \CYRV\
\CYRI\CYRN\CYRD-\CYRM\CYRN\CYRO\CYRG\CYRO\CYRO\CYRB\CYRR\CYRA\CYRZ\CYRI\CYRYA\CYRH\
\CYRM\CYRA\CYRK\CYRS\CYRI\CYRM\CYRA\CYRL\CYRSFTSN\CYRN\CYRERY\CYRH\
\CYRO\CYRB\CYRO\CYRB\CYRSHCH\CYRYO\CYRN\CYRN\CYRERY\CYRH\
\CYRF\CYRL\CYRA\CYRG\CYRO\CYRV}
\author{Иван Пенков}
\address{Университет Лотарингии, Нац. Центр Научн. Иссл., Лотарингский Институт им. Эли Картана, Объед. Иссл. Подраздел. 7502, Вандёвр в Нанси, F-54506 Франция}
\email{lucas.fresse@univ-lorraine.fr}
\author{Лука Фресс}
\address{Университет Якобса в Бремене, Кампус Ринг 1, 28759 Бремен, Германия}
\email{i.penkov@jacobs-university.de}
%\begin{center}
%\footnotesize{Лука Фресс и Иван Пенков}
%\end{center}
\dedicatory{Посвящается Эрнесту Борисовичу Винбергу  \\ к восьмидесятилетию }

УДК {512.816.2, 512.816.4, 512.818}

\footnotetext[0]{Ключевые слова. Классические инд-группы, обобщённые флаги, симметрические пары, вещетвенные формы, Двойственность Мацуки.}
\keywords{}



\begin{abstract}
%We show an orbit duality for ind-varieties of maximal generalized flags associated to classical ind-groups, which extends the duality shown by Matsuki for finite-dimensional flag varieties. In addition we provide explicit parametrizations of orbits.
Мы распространяем двойственность Мацуки на произвольные инд-многообразия максимальных обобщённых флагов, другими словами, на любое однородное инд-многообразие $\mathbf{G}/\mathbf{B}$ для классической инд-группы $\mathbf{G}$  и расщепляющей Борелевской инд-подгруппы $\mathbf{B}\subset\mathbf{G}$.  Сначала мы приводим явную комбинаторную версию двойственности Мацуки в конечномерном случае, включающую явную параметризацию  $K$- и $G^0$-орбит на $G/B$.  После того, как мы доказываем двойственность Мацуки в бесконечномерном случае,  мы даём необходимые и достаточные условия на Борелевскую инд-подгруппу $\mathbf{B}\subset\mathbf{G}$ для существования открытых и замкнутых $\mathbf{K}$- и $\mathbf{G}^0$-орбит на $\mathbf{G}/\mathbf{B}$, где $\left(\mathbf{K},\mathbf{G}^0\right)$ --- согласованная пара, состоящая из симметрической инд-подгруппы $\mathbf{K}$ и вещественной формы $\mathbf{G}^0$  инд-группы $\mathbf{G}$.
\end{abstract}

%\keywords{Classical ind-groups; generalized flags; symmetric pairs; real forms; Matsuki duality}

\begin{footnotesize}

\end{footnotesize}
\maketitle


\section{Введение}

\label{S.intro}

В настоящей статье мы распространяем двойственность Мацуки на
инд-много\-образия максимальных обобщённых флагов, то есть, на однородные пространства вида $\mathbf{G}/\mathbf{B}$ для $\mathbf{G}=\mathrm{GL}(\infty)$, $\mathrm{SL}(\infty)$, $\mathrm{SO}(\infty)$, $\mathrm{Sp}(\infty)$.  В случае  конечномерной редуктивной алгебраической группы  $G$, двойственность Мацуки \cite{Gindikin-Matsuki, Matsuki, Matsuki3} это биекция между (конечным) множеством $K$-орбит на $G/B$ и множеством $G^0$-орбит на $G/B$, где $K$ является симметрической подгруппой $G$ и $G^0$ является вещественной формой $G$.  Более того, эта биекция обращает отношения включения на замыканиях орбит.  В частности, замечательная теорема об единственности замкнутой $G^0$-орбиты на $G/B$, см. \cite{W1}, следует  через двойственность  Мацуки из единственности  открытой (по Зарисскому) $K$-орбиты на $G/B$.
В монографии \cite{FHW} двойственность Мацуки была использована в качестве отправной точки для изучения  пространств циклов.

Если  $\mathbf{G}=\mathrm{GL}(\infty)$, $\mathrm{SL}(\infty)$, $\mathrm{SO}(\infty)$, $\mathrm{Sp}(\infty)$ --- классическая инд-группа, то её Борелевские инд-подгруппы не являются ни $\mathbf{G}$-сопряжёнными, ни $\mathrm{Aut}(\mathbf{G})$-сопряжёнными, поэтому существет много инд-многообразий вида $\mathbf{G}/\mathbf{B}$.  Мы показываем, что двойственность Мацуки продолжается на произвольное инд-много\-образие $\mathbf{G}/\mathbf{B}$, где $\mathbf{B}$ --- расщепляющая Борелевская инд-подгруппа  инд-группы $\mathbf{G}$, а $\mathbf{G}=\mathrm{GL}(\infty)$, $\mathrm{SL}(\infty)$, $\mathrm{SO}(\infty)$, $\mathrm{Sp}(\infty)$.  В бесконечномерном случае структура $\mathbf{G}^0$-орбит и $\mathbf{K}$-орбит на $\mathbf{G}/\mathbf{B}$ сложнее, чем в конечномерном случае, и  всегда существует бесконечно много орбит.

 Изучение $\mathbf{G}^0$-орбит на $\mathbf{G}/\mathbf{B}$ для $\mathbf{G}=\mathrm{GL}(\infty),\mathrm{SL}(\infty)$ было начато в работе~\cite{Ignatyev-Penkov-Wolf} и  продолжено в работе~\cite{W}.  В частности, в статье \cite{Ignatyev-Penkov-Wolf} было показано, что для некоторых веществнных форм $\mathbf{G}^0$ существуют расщепляющие Борелевские инд-подгруппы $\mathbf{B}\subset \mathbf{G}$ такие, что $\mathbf{G}/\mathbf{B}$ не имеет ни открытых, ни замкнутых  $\mathbf{G}^0$-орбит.  %On the other hand, Wolf's Theorem about the uniqueness of a closed $G^0$-orbit on $G/B$ in the finite-dimensional case extends also to the infinite-dimensional case: here the correct statement is that a closed $\mathbf{G}^0$-orbit is unique whenever it exists.
Нам не известно о каких-либо предшествующих исследованиях    $\mathbf{K}$-орбит на $\mathbf{G}/\mathbf{B}$ для $\mathbf{G}=\mathrm{GL}(\infty),\mathrm{SL}(\infty),\mathrm{SO}(\infty),\mathrm{Sp}(\infty)$.  Двойственность, которую мы устанавливаем в этой статье, показывает, что структура $\mathbf{K}$-орбит на $\mathbf{G}/\mathbf{B}$ является ``зеркальным отражением'' структуры $\mathbf{G}^0$-орбит на $\mathbf{G}/\mathbf{B}$.  В частности, факт, что $\mathbf{G}/\mathbf{B}$ допускает не более одной замкнутой $\mathbf{G}^0$-орбиты теперь следствие очевидного утверждения, что $\mathbf{G}/\mathbf{B}$ допускает не более одной открытой по Зарисскому $\mathbf{K}$-орбиты.  %In the case of classical ind-groups $\mathrm{SO}(\infty)$ and $\mathrm{Sp}(\infty)$, it is natural to undertake a study of the structure of $\mathbf{K}$-orbits first, and then transfer the results of $\mathbf{G}^0$-orbits via Matsuki duality...

%MENTION ALSO $G/P$!

Наш основной результат можно сформулировать следующим образом.
Пусть $(\mathbf{G},\mathbf{K},\mathbf{G}^0)$ --- одна из троек
перечисленных в разделе \ref{S2.1},
состоящих из классической (комплексной)
инд-группы $\mathbf{G}$, симметрической инд-подгруппы $\mathbf{K}\subset\mathbf{G}$
и соответствующей вещественной формы $\mathbf{G}^0\subset\mathbf{G}$.
Пусть $\mathbf{B}\subset\mathbf{G}$ --- расщепляющая Борелевская инд-подгруппа
такая, что $\mathbf{X}:=\mathbf{G}/\mathbf{B}$ является инд-многообразием максимальных обобщённых флагов (изотропных для типов B, C, D)
слабо согласованных с некоторым базисом в $\textbf{V} $, соответствующим
выбору $\mathbf{K}$ и $\mathbf{G}^0$ в смысле разделов \ref{S2.1}, \ref{S2.3}.
Имеется естественное разложение $\mathbf{G}=\bigcup_{n\geq 1}G_n$
и $\mathbf{X}=\bigcup_{n\geq 1}X_n$.
Здесь $G_n$ --- конечномерная алгебраическая группа,
$X_n$ --- полное многообразие флагов  в $G_n$,
и вложение $X_n\subset\mathbf{X}$, в частности, $G_n$-эквивариантно.
Подгруппы $K_n:=\mathbf{K}\cap G_n$ и $G^0_n:=\mathbf{G}^0\cap G_n$
являются, соответственно, симметрической подгруппой и вещественной формой группы $G_n$.
Подробности см. в разделе  \ref{S4.4}  .

\begin{theorem}
\label{theorem-1}
\begin{itemize}
\item[\rm (a)] Для всех $n\geq 1$ включение $X_n\subset \mathbf{X}$
индуцирует вложение множеств орбит $X_n/K_n\hookrightarrow \mathbf{X}/\mathbf{K}$
и $X_n/G^0_n\hookrightarrow \mathbf{X}/\mathbf{G}^0$.
\item[\rm (b)] Существет биекция $\Xi:\mathbf{X}/\mathbf{K}\to\mathbf{X}/\mathbf{G}^0$
такая, что диаграмма
\begin{equation}
\label{diag}
\xymatrix{
X_n/K_n \ar@{->}[d]^{\Xi_n} \ar@{^{(}->}[r] & \mathbf{X}/\mathbf{K} \ar@{->}[d]^{\Xi} \\
X_n/G^0_n \ar@{^{(}->}[r] & \mathbf{X}/\mathbf{G}^0
}
\end{equation}
комутативна. Здесь $\Xi_n$ означает двойственность Мацуки.
\item[\rm (c)] Для каждой $\mathbf{K}$-орбиты $\bs{\mathcal{O}}\subset\mathbf{X}$
пересечение $\bs{\mathcal{O}}\cap\Xi(\bs{\mathcal{O}})$ состоит из одной $\mathbf{K}\cap\mathbf{G}^0$-орбиты.
\item[\rm (d)] Биекция $\Xi$ обращает отношение включения на замыканиях орбит.
В частности, $\Xi$ отображает открытые (соответственно, замкнутые) $\mathbf{K}$-орбиты в
замкнутые (соответственно, открытые) $\mathbf{G}^0$-орбиты.
\end{itemize}
\end{theorem}

На самом деле наши результаты намного точнее: в предложениях \ref{P4-1.1}, \ref{P4.2-2}, \ref{P4.3} мы показывем, что $\mathbf{X}/\mathbf{K}$ и $\mathbf{X}/\mathbf{G}^0$ допускают одну и ту же явную параметризацию индуктивным  пределом подходящих одинаковых параметризаций  $X_n/K_n$ и $X_n/G_n^0$.  Это влечёт биекцию $\Xi$ из теоремы  \ref{theorem-1}\,{\rm (b)}. Утверждения {\rm (a)} и {\rm (b)} из теоремы \ref{theorem-1} удовлетворяют нашим требованиям (\ref{T-proof.1}), (\ref{T-proof.2}), (\ref{T-proof.3}) ниже. Теорема \ref{theorem-1}\,{\rm (c)} следует из соответсвующих утверждений в предложениях \ref{P4-1.1}, \ref{P4.2-2}, \ref{P4.3}.
Наконец, теорема \ref{theorem-1}\,{\rm (d)} получается из
теоремы \ref{theorem-1}\,{\rm (a)} --- {\rm (b)}, определения инд-топологии,
и из факта, что двойственность $\Xi_n$ обращает отношение включения на замыканиях орбит.

Например, в случае, когда $\mathbf{G}=\mathrm{GL}(\infty)$ и $\mathbf{K}\subset \mathbf{G}$ --- инд-подгруппа, сохраняющая ортогональную форму $\omega$, мы показываем, что оба множества орбит $\mathbf{X}/\mathbf{K}$ и $\mathbf{X}/\mathbf{G}^0$ параметризуются инволюциями $ \omega: \mathbb{N}^* \to \mathbb{N}^*$ такими, что $\omega(\ell)=i(\ell)$ для почти всех $\ell \in \mathbb{N}^*$, где $i$ --- инволюция, индуцированная матрицей формы $\omega$ в подходящем базисе натурального представления  $\mathbf{G}$ (см. раздел \ref{S4.1}).

Наши методы основаны на классификации симметрических подгрупп и вещественных форм классических простых алгебраических групп. Возможно, будущие исследования позволят привести доказательство наших результатов, не зависящее от классификации.

\subsection*{Организация статьи}

В разделе \ref{S2} мы вводим обозначения для классических инд-групп, симметрических инд-подгрупп и вещественнх форм.
Мы напоминаем несколько основных фактов о конечномерных многообразиях  флагов, а также понятие инд-многообразия обобщённых флагов \cite{Dimitrov-Penkov,Ignatyev-Penkov}.
В разделе \ref{S3} мы строим
совместную параметризацию  $K$- и $G^0$-орбит для конечномерных многообразий  флагов.
Возможно, эта параметризация  известна (см. \cite{MO,Y}),
однако мы не нашли ссылки, в которой она была бы представлена в точности, как у нас. Ради полноты статьи мы приводим явное доказательство этих результатов.
В разделе \ref{S4} мы излагаем наши основные результаты о параметризации $\mathbf{K}$- и $\mathbf{G}^0$-орбит в инд-многообразиях обобщённых флагов. Сформулированная выше теорема \ref{theorem-1}  является основным следствием наших результатов.
В секции \ref{S5} мы указываем некоторые дальнейшие следствия из  основных результатов.

Через $\NN$ мы обозначаем множество положительных целых чисел.
$|A|$  обозначает мощность множества $A$.  Симметрическая группа на $n$ символах будет обозначаться $\mathfrak{S}_n$, а $\mathfrak{S}_\infty=\lim\limits_{\longrightarrow}\,\mathfrak{S}_n$ ---  бесконечная симметрическая группа.
Часто мы будем писать $w_k$ для образа $w(k)$ элемента $k$ перестановки $w$.
Через $(k;\ell)$ мы обозначим перестановку, которая меняет местами $k$ и $\ell$.
Мы будем использовать полужирные символы для обозначения инд-многообразий.
В конце статьи дан список обозначений.

\subsection*{Благодарности}


Мы благодарим Алана Хаклберри и Михаила Игнатьева за поддержку идеи изучения двойственности Мацуки.  Мы также признательны рецензенту за вдумчивые замечания. Первый автор получил частичную поддержку от Немецкого фонда научных исследований DFG через грант PE $980/6-1$. Второй автор получил частичную поддержку от Израильского научного фонда ISF через грант 797/14, а также от проекта ANR-15-CE40-0012. 
%\thanks{L. Fresse: supported in part by ISF Grant Nr. 797/14 and by ANR project GeoLie (ANR-15-CE40-0012)}
%\thanks{I. Penkov: supported in part by DFG Grant PE 980/6-1}

\section{Обозначения и известные факты}

\label{S2}



\subsection{Классические группы и классические инд-группы}

\label{S2.1}

Пусть $\mathbf{V}$ - комлексное векторное пространство счётной размерности, с упорядоченным базисом
$E=  (e_1,e_2,\ldots) \\ =  (e_\ell)_{\ell\in \NN}$.
Каждый вектор $x\in\mathbf{V}$ отождествляется со столбцом своих координат в базисе $E$,
и $x\mapsto \overline{x}$ это отображаение комплексного сопряжения  по отношению к базису $E$.
Мы рассматриваем также конечномерные подпространства $V=V_n:=\langle e_1,\ldots,e_n\rangle_\mathbb{C}$ в $\mathbf{V}$.

Классическая инд-группа $\mathrm{GL}(\infty)$ определяется как
\[\mathrm{GL}(\infty)=\mathbf{G}(E):=\{g\in\mathrm{Aut}(\mathbf{V}):g(e_\ell)=e_\ell\mbox{ для всех $\ell\gg 1$}\}=\bigcup_{n\geq 1}\mathrm{GL}(V_n).\]
%
%(resp. $(K,G^0)$)
%of a symmetric subgroup $\mathbf{K}\subset \mathbf{G}(E)$ and a real form $\mathbf{G}^0\subset \mathbf{G}(E)$.
%of $\mathbf{G}(E)$ (resp. $\mathrm{GL}(V)$).
Вещественные формы инд-группы $\mathrm{GL}(\infty)$ хорошо известны и восходят к работе Баранова~\cite{Baranov}.  Ниже мы перечисляем согласованные пары $(\mathbf{K},\mathbf{G}^0)$, где $\textbf{G}^0$ --- вещественная форма инд-группы $\mathbf{G}$, а $\mathbf{K}$  --- симметрическая инд-подгруппа  $\mathbf{G}$. Пары $\left(\mathbf{K},\mathbf{G}^0\right)$, которые мы рассматриваем, согласованы следующим образом: для разложения $\mathbf{G}$ в качестве объединения $\bigcup_n\,\mathrm{GL}\left(V_n\right)$, подгруппа $K_n:=\mathbf{K}\cap\mathrm{GL}\left(V_n\right)$ является симметрической подгруппой группы $\mathrm{GL}\left(V_n\right)$, $G_0^n:=\mathbf{G}^0\cap\mathrm{GL}\left(V_n\right)$ является вещественной формой группы $\mathrm{GL}\left(V_n\right)$, и $K_n\cap G_n^0$ --- максимальная компактная подгруппа $G_n^0$.

\subsubsection{Типы A1 и A2}
\label{S:A1A2}

Через $\Omega$ обозначим $\NN\times\NN$-матрицу вида
\begin{equation}
\label{omega}
\Omega=\left(
\begin{array}{cccc}
J_1 & & (0) \\
 & J_2 \\
(0) & & \ddots
\end{array}
\right),\
\mbox{где}\
\left\{
\begin{array}{ll}
J_k\in\left\{\left(\begin{matrix} 0 & 1 \\ 1 & 0 \end{matrix}\right),\left(1\right)\right\}
& \parbox{2.6cm}{(ортогональный случай, тип A1)} \\[4mm]
J_k=\left(\begin{matrix} 0 & 1 \\ -1 & 0 \end{matrix}\right) & \parbox{2.6cm}{(симплектический случай,  тип A2)}
\end{array}
\right.
\end{equation}
Билинейная форма
\[\omega(x,y):={}^tx\Omega y\quad (x,y\in \mathbf{V})\]
симметрична для типа A1 и симплектична для типа A2, в то время как отображение
\[\gamma(x):=\Omega\overline{x}\quad (x\in \mathbf{V})\]
это инволюция пространства $\mathbf{V}$ для типа A1 и антиинволюция для типа A2.
Пусть

$$ \mathbf{K}=\mathbf{G}(E,\omega):=\{g\in \mathbf{G}(E):\omega(gx,gy)=\omega(x,y)\ \forall x,y\in\mathbf{V}\} $$
и
  $$ \mathbf{G}^0:=\{g\in \mathbf{G}(E):\gamma(gx)=g\gamma(x)\ \forall x\in\mathbf{V}\}.$$

%In other words we get $(\mathbf{K},\mathbf{G}^0)=(\mathrm{SO}(\infty),\mathrm{GL}(\infty,\mathbb{R}))$ in type A1 and $(\mathbf{K},\mathbf{G}^0)=(\mathrm{Sp}(\infty),\mathrm{GL}(\infty,\mathbb{H}))$ in type A2, where $\mathbb{H}$ stands for the field of quaternions.


\subsubsection{Тип A3}
\label{S:A3}
Зафиксируем (нетривиальное) разложение $\NN=N_+\sqcup N_-$ и положим
\begin{equation}
\label{phi}
\Phi=\left(
\begin{array}{cccc}
\epsilon_1 & & (0) \\
 & \epsilon_2 \\
(0) & & \ddots
\end{array}
\right),
\end{equation}
где $\epsilon_\ell=1$ for $\ell\in N_+$ и $\epsilon_\ell=-1$ для $\ell\in N_-$.
Таким образом
\[\phi(x,y):={}^t\overline{x} \Phi y\quad(x,y\in\mathbf{V})\]
--- Эрмитова форма сигнатуры $(|N_+|,|N_-|)$ и
\[\delta(x):=\Phi x\quad(x\in\mathbf{V})\]
--- инволюция. Наконец положим
\[
\mathbf{K}:=\{g\in \mathbf{G}(E):\delta(gx)=g\delta(x)\ \forall
x\in\mathbf{V}\}
\]
и
\[
\mathbf{G}^0:=\{g\in \mathbf{G}(E):\phi(gx,gy)=\phi(x,y)\ \forall
 x,y\in\mathbf{V}\}.
\]
%In other words the pair $(\mathbf{K},\mathbf{G}^0)$ is of the form $(\mathrm{GL}(\infty)\times \mathrm{GL}(\infty),\mathrm{U}(\infty,\infty))$ (if $N_+,N_-$ are both infinite) or $(\mathrm{GL}_p(\mathbb{C})\times \mathrm{GL}(\infty),\mathrm{U}(p,\infty))$ for some finite $p\geq 0$ (otherwise).

%\medskip

%In types A1--A3 $\mathbf{K}$ and $\mathbf{G}^0$ are respectively a symmetric subgroup and a real form of $\mathbf{G}=\mathrm{GL}(\infty)$.


\subsection*{Типы B, C, D}
%\label{S:BCD}
Далее, мы опишем пары $(\mathbf{K},\mathbf{G}^0)$, связанные
с другими классическими инд-группами $\mathrm{SO}(\infty)$ и $\mathrm{Sp}(\infty)$.  Пусть $\mathbf{G}=\mathbf{G}(E,\omega)$,
где $\omega$ --- (симметрическая или симплектическая) билинейная форма, задающаяся матрицей $\Omega$ как и в  (\ref{omega}).
С учётом  (\ref{omega}), для каждого $\ell\in\NN$ существует единственное число $\ell^*\in\NN$ такое, что
\[\omega(e_\ell,e_{\ell^*})\not=0.\]
Далее, $\ell^*\in\{\ell-1,\ell,\ell+1\}$. Отображение $\ell\mapsto\ell^*$ это инволюция множества $\NN$.

%\medskip

\subsubsection{Типы BD1 и C2}
\label{S:BD1C2}

Предположим, что форма $\omega$ симметрична для типа BD1 и симплектична для типа C2.
Зафиксируем (нетривиальное) разложение $\NN=N_+\sqcup N_-$ такое, что
\[\forall \ell\in\NN,\ \ell\in N_+\Leftrightarrow \ell^*\in N_+\]
и ограничение формы $\omega$ на каждое подпространство $\mathbf{V}_+:=\langle e_\ell:\ell\in N_+\rangle_\mathbb{C}$ и
$\mathbf{V}_-:=\langle e_\ell:\ell\in N_-\rangle_\mathbb{C}$ невырождено.
Пусть  $\Phi,\phi,\delta$ --- такие же как и в разделе ~\ref{S:A3}.  Тогда  положим
\begin{eqnarray}
\label{N1}
 \mathbf{K}:=\{g\in \mathbf{G}(E,\omega):\delta(gx)=g\delta(x)\ \forall
x\in\mathbf{V}\}
\end{eqnarray}
и
\begin{eqnarray}
\label{N2}
 \mathbf{G}^0:=\{g\in \mathbf{G}(E,\omega):\phi(gx,gy)=\phi(x,y)\ \forall
 x,y\in\mathbf{V}\}.
\end{eqnarray}
%The triple $(\mathbf{G},\mathbf{K},\mathbf{G}^0)$ is thereby of the form $(\mathrm{SO}(\infty),\mathrm{O}(p)\times \mathrm{O}(q),\mathrm{O}(p,q))$ (in type BD1) or $(\mathrm{Sp}(\infty),\mathrm{Sp}(p)\times \mathrm{Sp}(q),\mathrm{Sp}(p,q))$ (in type C2), where $(p,q)=(|N_+|,|N_-|)$.

%\medskip

\subsubsection{Типы C1 и D3}
\label{S:C1D3}
Предположим, что форма $\omega$ симметрична для типа D3 и симплектична для типа C1.
Зафиксируем разложение $\NN=N_+\sqcup N_-$, удовлетворяющее
\[
\forall \ell\in\NN,\ \ell\in N_+\Leftrightarrow \ell^*\in N_-.
\]
Заметим, что из-за этого каждый блок  $J_k$ в (\ref{omega}) должен иметь размер $2$.
В этой ситуации $\mathbf{V}_+:=\langle e_\ell:\ell\in N_+\rangle_\mathbb{C}$ и $\mathbf{V}_-:=\langle e_\ell:\ell\in N_-\rangle_\mathbb{C}$
являются максимальными изотропными подпространствами для формы $\omega$.
Определим $\Phi,\phi,\delta$ так же как и в разделе~\ref{S:A3}. Наконец,  определим инд-подгруппы $\mathbf{K},\mathbf{G}^0\subset\mathbf{G}$  как  в (\ref{N1}), (\ref{N2}).
%Hence $(\mathbf{G},\mathbf{K},\mathbf{G}^0)$ is of the form $(\mathrm{SO}(\infty),\mathrm{GL}(\infty),\mathrm{O}(p,q))$ (type D1) or $(\mathrm{Sp}(\infty),\mathrm{GL}(\infty),\mathrm{O}(p,q))$ (type C1).

%\medskip

\subsection*{Конечномерный случай}
В следующей таблице представлены пересечения $G=\textbf{G}\cap\mathrm{GL}\left(V_n\right)$, $K=\textbf{K}\cap\mathrm{GL}\left(V_n\right)$, $G^0=\textbf{G}^0\cap\mathrm{GL}\left(V_n\right)$, где число $n=2m$ чётное для типов A2, C1, C2 и D3.
Для типов A3, BD1 и C2, полагаем $(p,q)=(|N_+\cap\{1,\ldots,n\}|,|N_-\cap\{1,\ldots,n\}|)$.
Через $\mathbb{H}$  обозначаем тело кватернионов.
Таким образом мы получаем классические конечномерные симметрические пары и вещественные формы
(см., например, \cite{Berger,OV,Otha1}).

\begin{center}
\renewcommand{\arraystretch}{1.25}
\begin{tabular}{|c|c|c|c|}
\hline
type & $G:=\mathbf{G}\cap\mathrm{GL}\left(V_n\right)$ & $K:=\mathbf{K}\cap\mathrm{GL}\left(V_n\right)$ & $G^0:=\mathbf{G}^0\cap\mathrm{GL}\left(V_n\right)$ \\
\hline
A1 & & $\mathrm{O}_n(\mathbb{C})$ & $\mathrm{GL}_n(\mathbb{R})$ \\
A2 & $\mathrm{GL}_n(\mathbb{C})$ & $\mathrm{Sp}_n(\mathbb{C})$ & $\mathrm{GL}_m(\mathbb{H})$ \\
A3 & & $\mathrm{GL}_p(\mathbb{C})\times\mathrm{GL}_q(\mathbb{C})$ & $\mathrm{U}_{p,q}(\mathbb{C})$ \\
\hline
BD1 & $\mathrm{O}_n(\mathbb{C})$ & $\mathrm{O}_p(\mathbb{C})\times\mathrm{O}_q(\mathbb{C})$ & $\mathrm{O}_{p,q}(\mathbb{C})$ \\
\hline
C1 & ${}_{{}_{\mbox{$\mathrm{Sp}_n(\mathbb{C})$}}}$ & $\mathrm{GL}_m(\mathbb{C})$ & $\mathrm{Sp}_n(\mathbb{R})$ \\[-1.5mm]
C2 & & $\mathrm{Sp}_p(\mathbb{C})\times\mathrm{Sp}_q(\mathbb{C})$ & $\mathrm{Sp}_{p,q}(\mathbb{C})$ \\
\hline
D3 & $\mathrm{O}_n(\mathbb{C})=\mathrm{O}_{2m}(\mathbb{C})$ & $\mathrm{GL}_m(\mathbb{C})$ & $\mathrm{O}^*_n(\mathbb{C})$ \\
\hline
\end{tabular}
\end{center}
В каждом случае  вещественная форма $G^0$ получается из $K$ таким образом, чтобы пересечение $K\cap G^0$ было максимальной компактной подгруппой $G^0$.
Наоборот, $K$ получается из $G^0$ как комплексификация максимальной компактной подгруппы.

\subsection{Конечномерные многообразия флагов}

\label{S2.2}

Напомним, что $V=V_n$.
Многообразие флагов $X:=\mathrm{GL}(V)/B=\{gB:g\in \mathrm{GL}(V)\}$ (для Борелевской подгруппы $B\subset \mathrm{GL}(V)$)
может  рассматриваться как   множество Борелевских подгрупп $\{gBg^{-1}:g\in \mathrm{GL}(V)\}$
или как множество полных флагов
\begin{equation}
\label{2.2.X}
\big\{\mathcal{F}=(F_0\subset F_1\subset\ldots\subset F_n=V): \dim F_k=k\ \mbox{ для всех $k$}\big\}.
\end{equation}
Для каждого полного флага $\mathcal{F}$ обозначим через $B_\mathcal{F}:=\{g\in \mathrm{GL}(V):g\mathcal{F}=\mathcal{F}\}$  сответствующую Борелевскую подгруппу.
Если $(v_1,\ldots,v_n)$ --- базис пространства $V$, то мы будем писать
\[\mathcal{F}(v_1,\ldots,v_n):=\big(0\subset\langle v_1\rangle_\mathbb{C}\subset\langle v_1,v_2\rangle_\mathbb{C}\subset\ldots\subset\langle v_1,\ldots,v_n\rangle_\mathbb{C}\big)\in X.\]

\medskip

\paragraph{{\bf Разложение Брюа.}}
Двойное многообразие флагов $X\times X$ имеет конечное число $\mathrm{GL}(V)$-орбит, параметризуемых перестановками $w\in\mathfrak{S}_n$.
В самом деле, если даны два флага $\mathcal{F}=(F_k)_{k=0}^n$ и
$\mathcal{F}'=(F'_\ell)_{\ell=0}^n$, то существет единственная перестановка
$w=:w(\mathcal{F},\mathcal{F}')$ такая, что
\[\dim F_k\cap F'_\ell=\big|\big\{j\in\{1,\ldots,\ell\}:w_j\in\{1,\ldots,k\}\big\}\big|.\]
Перестановка $w(\mathcal{F},\mathcal{F}')$ называется \textit{относительной позицией} пары $(\mathcal{F},\mathcal{F}')\in X\times X$.
%Set
%\[\mathbb{O}_w:=\big\{(\mathcal{F},\mathcal{F}')\in X\times X:w(\mathcal{F},\mathcal{F}')=w\big\}.\]
При этом
\[X\times X=\bigsqcup_{w\in\mathfrak{S}_n}\mathbb{O}_w,
\quad\mbox{где $\mathbb{O}_w:=\big\{(\mathcal{F},\mathcal{F}')\in X\times X:w(\mathcal{F},\mathcal{F}')=w\big\}$},\]
есть разложение $X\times X$ в объединение $\mathrm{GL}(V)$-орбит.
Единственная замкнутая орбита это $\mathbb{O}_{\mathrm{id}}$, и единственная открытая орбита это $\mathbb{O}_{w_0}$, где $w_0$ --- инволюция, заданная равенством
$w_0(k)=n-k+1$ для всех $k$. Отображение
$\mathbb{O}_w\mapsto\mathbb{O}_{w_0w}$ это инволюция на множестве орбит, и она обращает отношение включения на замыканиях орбит. Представители $\mathbb{O}_w$ описываются следующим образом:
для каждого базиса $(v_1,\ldots,v_n)$ of $V$ мы имеем
\[\big(\mathcal{F}(v_1,\ldots,v_n),\mathcal{F}(v_{w_1},\ldots,v_{w_n})\big)\in\mathbb{O}_w.\]

\medskip

\paragraph{{\bf Многообразие изотропных флагов}} Пусть $V$ обладает невырожденной симметрической или симплектической  билинейной формой $\omega$.
Для подпространства $F\subset V$, положим $F^\perp=\{x\in V:\omega(x,y)=0\ \forall y\in F\}$.
Многообразие  изотропных флагов $X_\omega$ является подмногообразием  в $X$, где
\begin{equation}
\label{2.2.Xomega}
X_\omega=\{\mathcal{F}=(F_k)_{k=0}^n\in X:F_k^\perp=F_{n-k}\ \forall k=0,\ldots,n\}.
\end{equation}
Оно обладает транзитивно действующей подгруппой
$G(V,\omega)\subset\mathrm{GL}(V)$ автоморфизмов, сохраняющих форму $\omega$.

\begin{lemma}
\label{lemma-2.2.1}
{\rm (a)}
Для каждого эндоморфизма  $f\in\mathrm{End}(V)$, обозначим через $f^*\in\mathrm{End}(V)$  эндоморфизм,  сопряжённый к $f$ относительно формы $\omega$.
Пусть $H\subset \mathrm{GL}(V)$ подгруппа, удовлетворяющая условию
\begin{equation}
\label{2.2.1}
\mathbb{C}[g^*g]\cap\mathrm{GL}(V)\subset H\ \mbox{ для всех $g\in H$}.
\end{equation}
Предположим, что $\mathcal{F}\in X_\omega$ и $\mathcal{F}'\in X_\omega$ принадлежат одной $H$-орбите в $X$.  Тогда они принадлежат одной $H\cap G(V,\omega)$-орбите в $X_\omega$. \\
{\rm (b)} Пусть $H=\{g\in \mathrm{GL}(V):g(V_+)=V_+,\ g(V_-)=V_-\}$, где $V=V_+\oplus V_-$ это разложение такое, что
$(V_+^\perp,V_-^\perp)=(V_+,V_-)$ или $(V_-,V_+)$.
Тогда (\ref{2.2.1}) выполняется.
\end{lemma}

\begin{proof}
{\rm (a)}
Отметим, что $G(V,\omega)=\{g\in\mathrm{GL}(V):g^*=g^{-1}\}$.
Рассмотрим элемент $g\in H$ такой, что $\mathcal{F}'=g\mathcal{F}$.
Равенство $(gF)^\perp=(g^*)^{-1}F^\perp$ выполняется для всех подпространств $F\subset V$.
Так как $\mathcal{F}$ и $\mathcal{F}'$ принадлежат $X_\omega$, мы получаем $\mathcal{F}'=(g^*)^{-1}\mathcal{F}$,
поэтому $g^*g\mathcal{F}=\mathcal{F}$.
Пусть $g_1=g^*g$.
Согласно \cite[Лемма 1.5]{Jantzen} существет полином $P(t)\in\mathbb{C}[t]$ такой, что
$P(g_1)^2=g_1$. Положим $h=P(g_1)$.
Тогда $h\in\mathrm{GL}(V)$ (так как $h^2=g_1\in\mathrm{GL}(V)$), и из (\ref{2.2.1}) вытекает, что на самом деле $h\in H$.
Далее, $h^*=h$ (поскольку $h\in\mathbb{C}[g_1]$ и $g_1^*=g_1$)
и $h\mathcal{F}=\mathcal{F}$ ( так как и каждое подпространство в $\mathcal{F}$ является $g_1$-подмодулем и, следовательно, также $h$-подмодулем).
Положим $h_1:=gh^{-1}\in H$.
Тогда, с одной стороны,
\[h_1^*=(h^*)^{-1}g^*=h^{-1}g_1g^{-1}=h^{-1}h^2g^{-1}=hg^{-1}=h_1^{-1}\,.\]
Поэтому $h_1\in H\cap G(V,\omega)$.
С другой стороны,
$h_1\mathcal{F}=gh^{-1}\mathcal{F}=g\mathcal{F}=\mathcal{F}'$, и часть (a) доказана. \\
{\rm (b)}
Равенство  $g^*(gF)^\perp=F^\perp$ (которое уже упоминалось), применяемое к $F=V_\pm$, даёт $g^*\in H$, и таким образом $g^*g\in H$, если $g\in H$. Поэтому выполняется (\ref{2.2.1}).
\end{proof}

\begin{remark}
Идея доказательства леммы  \ref{lemma-2.2.1}\,{\rm (a)} навеяна нам \cite[\S1.4]{Jantzen}. В  работах \cite{Nishiyama, Otha} читатель найдёт аналогичные результаты и обобщения.
\end{remark}

\subsection{Инд-многообразия обобщённых флагов}

\label{S2.3}

Напомним, что $\mathbf{V}$ обозначает комплексное векторное пространство счётной размерности,
с упорядоченным базисом $E= \ (e_\ell)_{\ell\in\NN}$.

\begin{definition}[\cite{Dimitrov-Penkov}]
Пусть $\mathcal{F}$ будет цепью подпространств в $\mathbf{V}$, то есть, множеством подпространств из $\mathbf{V}$, полностью упорядоченным по  включению . Пусть $\mathcal{F}'$ (соответственно $\mathcal{F}''$) --- подцепь, состоящая из всех $F\in\mathcal{F}$ с непосредственным последующим (соответственно предыдущим) элементом.
Через $s(F)\in\mathcal{F}''$  мы обозначим подпространство непосредственно следующее за  $F\in\mathcal{F}'$.

{\em Обобщённый флаг} в $\mathbf{V}$ это цепь подпространств $\mathcal{F}$ такая что:
\begin{itemize}
\item[\rm (i)] каждый $F\in\mathcal{F}$ обладает непосредственно последующим или предыдущим элементом,
то есть, $\mathcal{F}=\mathcal{F}'\cup\mathcal{F}''$;
\item[\rm (ii)] для каждого $v\in\mathbf{V}\setminus\{0\}$ существует единственное подпространство $F_v\in\mathcal{F}'$ такое, что $v\in s(F_v)\setminus F_v$, то есть,
$\mathbf{V}\setminus\{0\}=\bigcup_{F\in\mathcal{F}'}(s(F)\setminus F)$.
\end{itemize}

Обобщённый флаг назывется  {\em максимальным}, если он не содержится ни в каком другом обобщённом  флаге.
В частности, $\mathcal{F}$ максимален тогда и только тогда, когда $\dim s(F)/F=1$ для всех $F\in\mathcal{F}'$.
\end{definition}

\begin{notation}
Пусть $\sigma:\NN\to (A,\prec)$-- сюръективное отображение на полностью упорядоченное множество.
Пусть $\underline{v}=(v_1,v_2,\ldots)$ --- базис в $\mathbf{V}$.
Для каждого $a\in A$ положим
\[F'_a:=\langle v_\ell:\sigma(\ell)\prec a\rangle_\mathbb{C},\quad
F''_a:=\langle v_\ell:\sigma(\ell)\preceq a\rangle_\mathbb{C}.\]
Тогда $\mathcal{F}=\mathcal{F}_\sigma(\underline{v}):=\{F'_a,F''_a:a\in A\}$ --- обобщённый флаг, такой что
$\mathcal{F}'=\{F'_a:a\in A\}$, $\mathcal{F}''=\{F''_a:a\in A\}$, и $s(F'_a)=F''_a$ для всех $a$.
Мы называем такой обобщённый флаг {\em согласованным с базисом $\underline{v}$}.

Более того, $\mathcal{F}_\sigma(\underline{v})$ максимален, тогда и только тогда, когда отображение $\sigma$ биективно.

В дальнейшем будем использовать сокращённое обозначение $\mathcal{F}_\sigma:=\mathcal{F}_\sigma(E)$.

Отметим, что для каждого обобщённого флага существует согласованный с ним базис \cite[предложение 4.1]{Dimitrov-Penkov}.
Обобщённый флаг {\em слабо согласован с $E$}, если он согласован с некоторым базисом  $\underline{v}$ таким, что $E\setminus(E\cap \underline{v})$ --- конечное множество (аналогично,  $\dim\mathbf{V}/\langle E\cap \underline{v}\rangle_\mathbb{C}<\infty$).
\end{notation}

Группа $\mathbf{G}(E)$ (так же как и $\mathrm{Aut}(\mathbf{V})$) действует на множестве обобщённых флагов естественным образом.
Пусть $\mathbf{P}_\mathcal{F}\subset\mathbf{G}(E)$ обозначает инд-подгруппу, состоящую из элементов сохраняющих  $\mathcal{F}$.
Это замкнутая инд-подгруппа группы $\mathbf{G}(E)$.
Если $\mathcal{F}$ согласован с  $E$, то $\mathbf{P}_\mathcal{F}$ --- расщепляющая параболическая инд-подгруппа в  $\mathbf{G}(E)$ в том смысле, что она локально параболическая (то есть, существует исчерпание $\mathbf{G}(E)$  конечномерными редуктивными алгебраическими подгруппами $G_n$ такими, что все пересечения $\mathbf{P}_\mathcal{F}\cap G_n$ ---  параболические подгруппы в $G_n$) и содержит Картановскую инд-подгруппу  $\mathbf{H}(E)\subset \mathbf{G}(E)$, состоящую из диагональных элементов относительно базиса  $E$.
Более того, если $\mathcal{F}$ максимален, тогда $\mathbf{B}_\mathcal{F}:=\mathbf{P}_\mathcal{F}$ является расщепляющей Борелевской инд-подгруппой
(то есть все вышеупомянутые пересечения $\mathbf{B}_\mathcal{F}\cap G_n$ суть Борелевские подгруппы в $G_n$).

\begin{definition}[\cite{Dimitrov-Penkov}]
Два обобщённых флага $\mathcal{F},\mathcal{G}$ называются {\em $E$-соизмеримыми},
если $\mathcal{F},\mathcal{G}$ слабо согласованы с $E$, и существует изоморфизм $\phi:\mathcal{F}\to\mathcal{G}$ упорядоченных множеств и конечномерное подпространство $U\subset\mathbf{V}$ такие, что
\begin{itemize}
\item[\rm (i)] $\phi(F)+U=F+U$ для всех $F\in\mathcal{F}$;
\item[\rm (ii)] $\dim \phi(F)\cap U=\dim F\cap U$ для всех $F\in\mathcal{F}$.
\end{itemize}
\end{definition}

$E$-соизмеримость --- отношение эквивалентности на множестве обобщённых флагов слабо согласованных с  $E$. На самом деле, согласно следующему предложению, каждый класс эквивалентности состоит из одной $\mathbf{G}(E)$-орбиты.
Если $\mathcal{F}$ обобщённый флаг слабо согласован с  $E$, мы обозначим через $\mathbf{X}(\mathcal{F},E)$ множество обобщённых флагов, которые $E$-соизмеримы с  $\mathcal{F}$.

\begin{proposition}[\cite{Dimitrov-Penkov}]
\label{P2.3-1}
Множество $\mathbf{X}=\mathbf{X}(\mathcal{F},E)$ обладает естественной структурой инд-многообразия.
Далее, $\mathbf{X}$  является $\mathbf{G}(E)$-однородным инд-многообразием и отображение $g\mapsto g\mathcal{F}$ индуцирует изоморфизм инд-многооб\-разий $\mathbf{G}(E)/\mathbf{P}_\mathcal{F}\stackrel{\sim}{\to}\mathbf{X}$.
\end{proposition}

%Let $\mathfrak{S}_\infty$ be the group of permutations
%$w:\NN\to\NN$ such that $w(\ell)=\ell$ for all $\ell\gg 1$.



\begin{proposition}[\cite{Fresse-Penkov}]
\label{P2-3.2}
Пусть $\sigma:\NN\to (A,\prec)$ и $\tau:\NN\to (B,\prec)$ --- два сюръективных отображения  на полностью упорядоченные множества.
\begin{itemize}
\item[\rm (a)] Каждый $E$-согласованный обобщённый флаг в $\mathbf{X}(\mathcal{F}_\sigma,E)$ имеет вид $\mathcal{F}_{\sigma w}$ для $w\in \mathfrak{S}_\infty$.
Более того, $\mathcal{F}_{\sigma w}=\mathcal{F}_{\sigma w'}\Leftrightarrow
w'w^{-1}\in\mathrm{Stab}_\sigma:=\{v\in\mathfrak{S}_\infty:\sigma v=\sigma\}$.
\item[\rm (b)] Предположим, что $\mathcal{F}_\tau$ максимален (то есть, $\tau$ --- биекция), и, таким образом, $\mathbf{B}_{\mathcal{F}_\tau}$ --- расщепляющая Борелевская инд-подгруппа.
Каждая $\mathbf{B}_{\mathcal{F}_\tau}$-орбита  в $\mathbf{X}(\mathcal{F}_\sigma,E)$
содержит единственный элемент вида $\mathcal{F}_{\sigma w}$ для $w\in\mathfrak{S}_\infty/\mathrm{Stab}_\sigma$.
\item[\rm (c)] В частности, если $\mathcal{F}_\sigma$ и $\mathcal{F}_\tau$ оба максимальны (то есть, $\sigma,\tau$ --- биекции), то
\begin{eqnarray*}
\displaystyle \mathbf{X}(\mathcal{F}_\tau,E)\times\mathbf{X}(\mathcal{F}_\sigma,E)=\bigsqcup_{w\in\mathfrak{S}_\infty}(\pmb{\mathbb{O}}_{\tau,\sigma})_w ,
\end{eqnarray*}
где
\begin{eqnarray*}
(\pmb{\mathbb{O}}_{\tau,\sigma})_w:=\{(g\mathcal{F}_\tau,g\mathcal{F}_{\sigma w}):g\in\mathbf{G}(E)\}
\end{eqnarray*}
--- разложение пространства $\mathbf{X}(\mathcal{F}_\tau,E)\times\mathbf{X}(\mathcal{F}_\sigma,E)$ на $\mathbf{G}(E)$-орбиты.
\end{itemize}
\end{proposition}

\begin{remark}
Орбита $(\pmb{\mathbb{O}}_{\tau,\sigma})_w$
из предложения \ref{P2-3.2}\,{\rm (c)}
на самом деле состоит из всех пар обобщённых флагов $(\mathcal{F}_\tau(\underline{v}),\mathcal{F}_{\sigma w}(\underline{v}))$
слабо согласованных с базисом $\underline{v}=(v_1,v_2,\ldots)$.
\end{remark}




Предположим, $\mathbf{V}$ обладает невырожденной  симметрической или симплектической  формой $\omega$, значения которой на  элементах базиса $E$ могут быть получены из матрицы $\Omega$ в (\ref{omega}).

%\begin{notation}
%Given a generalized flag $\mathcal{F}$ let $\mathcal{F}^\perp:=\{F^\perp:F\in\mathcal{F}\}$,
%where we denote $F^\perp:=\{x\in\mathbf{V}:\omega(x,y)=0\ \forall y\in F\}$.
%
%If $\mathcal{F}=\mathcal{F}_\sigma$ is an $E$-compatible generalized flag, then so is $\mathcal{F}^\perp$.
%Specifically assume that $\sigma:\NN\to(A,\prec)$.
%Let $(A^*,\prec^*)$ be the totally ordered set given by $A^*=A$
%and $a\prec^*a'$ whenever $a\succ a'$.
%In view of $\Omega$
%for each $j\in\NN$ there is a unique $j^*\in\NN$ such that $\omega(e_j,e_{j^*})\not=0$.
%Then
%$\mathcal{F}^\perp=\mathcal{F}_{\sigma^*}$
%where $\sigma^*:\NN\to(A^*,\prec^*)$, $j\mapsto \sigma(j^*)$.
%\end{notation}

\begin{definition}
Обобщённый флаг $\mathcal{F}$ назывется {\em $\omega$-изотропным}, если отображение $F\mapsto F^\perp:=\{x\in\mathbf{V}:\omega(x,y)=0\ \forall y\in F\}$
является инволюцией обобщённого флага  $\mathcal{F}$.
\end{definition}

\begin{proposition}[\cite{Dimitrov-Penkov}]
\label{P2.3-3}
Пусть $\mathcal{F}$ есть  $\omega$-изотропный обобщённый флаг слабо согласованный с $E$.
Множество $\mathbf{X}_\omega(\mathcal{F},E)$ всех
$\omega$-изотропных обобщённых флагов, которые $E$-соизмеримы с $\mathcal{F}$,
является $\mathbf{G}(E,\omega)$-однородным, замкнутым инд-подмногообразием в $\mathbf{X}(\mathcal{F},E)$.
\end{proposition}

Наконец, мы хотели бы отметить, что одна из основных особенностей  инд-групп это факт, что их Борелевские подгруппы не $\mathrm{Aut}(\mathbf{G})$-сопряжены.  Вот три примера максимальных обобщённых флагов в $\mathbf{V}$, согласованных с базисом $E$, и таких, что их стабилизаторы в $\mathbf{G}(E)$ попарно не $\mathrm{Aut}(\mathbf{G})$-сопряжены. Более подробное обсуждение этих примеров см. в \cite{Dimitrov-Penkov}.

\begin{example}
{\rm (a)} Пусть  $\sigma_1:\mathbb{N}^*\to(\mathbb{N}^*,<)$, $\ell\mapsto\ell$.
Обобщённый флаг $\mathcal{F}_{\sigma_1}$ это восходящая цепочка подпространств
$\mathcal{F}_{\sigma_1}=\{0=F_0\subset F_1\subset F_2\subset\ldots\}$, изоморфная упорядоченному множеству $(\mathbb{N},<)$ .
\\
{\rm (b)} Пусть $\sigma_2:\mathbb{N}^*\to\Big(\{\frac{1}{n}:n\in\mathbb{Z}^*\},<\Big)$, $\ell\mapsto \frac{(-1)^\ell}{\ell}$.
Обобщённый флаг $\mathcal{F}_{\sigma_2}$ это цепь вида
$\mathcal{F}_{\sigma_2}=\{0=F_0\subset F_1\subset\ldots\subset F_{-2}\subset F_{-1}=\mathbf{V}\}$
не изоморфная как упорядоченное множество подмножеству в $(\mathbb{Z},<)$.
\\
{\rm (c)} Пусть $\sigma_3:\mathbb{N}^*\to(\mathbb{Q},<)$ - биекция.
В этом случе нет подпространств $F\in\mathcal{F}_{\sigma_3}$, у которых есть одновременно  непосредственно последующее и непосредственно предыдущее подпространство.

\end{example}

%our main purpose is...


\section{Параметризация орбит в конечномерном случае}

\label{S3}

В разделах ~\ref{S3.1}-\ref{S3.3}, мы приводим явную параметризацию $K$- и $G^0$-орбит в конечномерном случае. Все доказательства представленны в разделе  ~\ref{S3.5}.

\subsection{Типы A1 и A2}
\label{S3.1}
Обозначения такие же, как и в разделе~\ref{S:A1A2}.
Пространство $V=V_n:=\langle e_1,\ldots,e_n\rangle_\mathbb{C}$ обладает симметрической или симплектической билинейной формой $\omega(x,y)={}^tx\cdot \Omega\cdot y$ и сопряжением $\gamma(x)=\Omega\overline{x}$, которые  являются ограничениями на $V$ отображений $\omega,\gamma$, введённых в разделе \ref{S2.1}.
Это позволяет нам определить две инволюции многообразия флагов $X$:
\[
\mathcal{F}=(F_0,\ldots,F_n)\mapsto\mathcal{F}^\perp:=(F_n^\perp,\ldots,F_0^\perp)
\quad\mbox{и}\quad
\mathcal{F}\mapsto\gamma(\mathcal{F}):=(\gamma(F_0),\ldots,\gamma(F_n)),
\]
где $F^\perp\subset V$ обозначает подпространство, ортогональное  $F$ по отношению к форме $\omega$.

Положим $K=\{g\in\mathrm{GL}(V):\mbox{$g$ сохраняет $\omega$}\}$
и $G^0=\{g\in\mathrm{GL}(V):\gamma g=g\gamma\}$.

Через $\mathfrak{I}_n\subset\mathfrak{S}_n$ мы обозначим подмножество инволюций.
Если $n=2m$ чётное, тогда $\mathfrak{I}'_n\subset\mathfrak{I}_n$
будет подмножеством инволюций  без неподвижных точек.

\begin{definition}
\label{DA12}
Пусть $w\in\mathfrak{I}_n$. Зададим $\epsilon:=1$ для типа A1 и
$\epsilon:=-1$ для типа A2. Базис $(v_1,\ldots,v_n)$ в $V$ такой, что
\[\omega(v_k,v_\ell)=\left\{\begin{array}{ll} 1, & \mbox{если $w_k=\ell\geq k$} \\ \epsilon, & \mbox{если $w_k=\ell< k$} \\ 0, & \mbox{если $w_k\not=\ell$} \end{array}\right.
\ \mbox{ для всех $k,\ell\in\{1,\ldots,n\}$}\] называется {\em
$w$-двойственным}. Базис $(v_1,\ldots,v_n)$ в $V$ такой, что 
\[
\gamma(v_k)=\left\{\begin{array}{ll} \epsilon v_{w_k}, & \mbox{если $w_k\geq k$}
\\ v_{w_k}, & \mbox{если $w_k< k$}
\end{array}\right.
\ \mbox{ для всех $k\in\{1,\ldots,n\}$}
\]
называется {\em $w$-сопряжённым}. Положим
\begin{eqnarray*}
 & \mathcal{O}_w:=\{\mathcal{F}(v_1,\ldots,v_n):\mbox{$(v_1,\ldots,v_n)$ ---  $w$-двойственный базис}\}, \\[1mm]
 & \mathfrak{O}_w:=\{\mathcal{F}(v_1,\ldots,v_n):\mbox{$(v_1,\ldots,v_n)$ --- $w$-сопряжённый базис}\}.
\end{eqnarray*}
\end{definition}

%In type A2 there exist $w$-dual and/or $w$-conjugate bases only if there is no $k$ such that $w_k=k$, i.e., $w\in\mathfrak{I}'_n$. Part {\rm (a)} of the following proposition shows in particular that there always exist $w$-dual bases and $w$-conjugate bases for all $w\in\mathfrak{I}_n$ in type A1 (resp. $w\in\mathfrak{I}_n'$ in type A2); it is actually much more precise.


\begin{proposition}
\label{P1} Пусть $\mathfrak{I}_n^\epsilon=\mathfrak{I}_n$ для типа A1
и $\mathfrak{I}_n^\epsilon=\mathfrak{I}'_n$ для типа A2. Напомним обозначения
$\mathbb{O}_w$ и $w_0$, введённые  в разделе \ref{S2.2}.
\begin{itemize}
\item[\rm (a)]
Для каждого $w\in\mathfrak{I}_n^\epsilon$ мы имеем
$\mathcal{O}_w\not =\emptyset$, $\mathfrak{O}_w\not=\emptyset$
и
$\mathcal{O}_w\cap\mathfrak{O}_w $ $=\{\mathcal{F}(v_1,\ldots,v_n):\mbox{$(v_1,\ldots,v_n)$  $w$-двойствен и $w$-сопряжён}\}\not=\emptyset.$
\item[\rm (b)]
Для каждого $w\in \mathfrak{I}_n^\epsilon$,
\[
\mathcal{O}_w=\{\mathcal{F}\in
X:(\mathcal{F}^\perp,\mathcal{F})\in\mathbb{O}_{w_0w}\}
\quad\mbox{и}\quad \mathfrak{O}_w=\{\mathcal{F}\in
X:(\gamma(\mathcal{F}),\mathcal{F})\in\mathbb{O}_{w}\}.
\]
\item[\rm (c)]
Подмножества $\mathcal{O}_w$ ($w\in\mathfrak{I}_n^\epsilon$) суть в точности $K$-орбиты из $X$. Подмножества $\mathfrak{O}_w$
($w\in\mathfrak{I}_n^\epsilon$) суть в точности $G^0$-орбиты из $X$.
\item[\rm (d)]
Отображение $\mathcal{O}_w\mapsto\mathfrak{O}_w$ есть двойственность Мацуки.
\end{itemize}
\end{proposition}

\subsection{Тип A3}
\label{S3.2}
Обозначения такие же как и в разделе~\ref{S:A3}: пространство $V=V_n=\langle e_1,\ldots,e_n\rangle_\mathbb{C}$ обладает Эрмитовой формой $\phi(x,y)={}^t\overline{x}\Phi y$ и сопряжением $\delta(x)=\Phi x$, где $\Phi$ это диагональная матрица, состоящая из $\epsilon_1,\ldots,\epsilon_n\in\{+1,-1\}$ (эта матрица ни что иное, как левый верхний $n\times n$-угол  матрицы $\Phi$ из раздела \ref{S2.1}).

Зададим $V_+=\langle e_k:\epsilon_k=1\rangle_\mathbb{C}$ и $V_-=\langle e_k:\epsilon_k=-1\rangle_\mathbb{C}$.  Тогда $V=V_+\oplus V_-$.
Пусть $K=\{g\in\mathrm{GL}(V):\delta g=g\delta\}=\mathrm{GL}(V_+)\times\mathrm{GL}(V_-)$ и $G^0=\{g\in \mathrm{GL}(V):\mbox{$g$ сохраняет $\phi$}\}$.

Как и в разделе \ref{S3.1} мы получаем две инволюции многообразия флагов $X$:
\[\mathcal{F}=(F_0,\ldots,F_n)\mapsto\delta(\mathcal{F}):=(\delta(F_0),\ldots,\delta(F_n))\quad\mbox{и}\quad\mathcal{F}\mapsto\mathcal{F}^\dag:=(F_n^\dag,\ldots,F_0^\dag),\]
где $F^\dag\subset V$ обозначение для подпространства, ортогонального  $F\subset V$ относительно формы $\phi$.
Эрмитова форма на факторе $F/(F\cap F^\dag)$ индуцированная формой $\phi$ невырождена; мы обозначим её сигнатуру через $\varsigma(\phi:F)$. Пусть задан $\mathcal{F}=(F_0,\ldots,F_n)\in X$, и пусть
\[\varsigma(\phi:\mathcal{F}):=\big(\varsigma(\phi:F_\ell)\big)_{\ell=1}^n\in(\{0,\ldots,n\}^2)^n.\]
Тогда
\[\varsigma(\delta:\mathcal{F}):=\big((\dim F_\ell\cap V_+,\dim F_\ell\cap V_-)\big)_{\ell=1}^n\in(\{0,\ldots,n\}^2)^n\]
отмечает относительное положение $\mathcal{F}$ по отношению к подпространствам
$V_+$ и
$V_-$.

\medskip

\paragraph{\bf Комбинаторные обозначения}
Мы назовём {\em инволюцией со знаками} пару $(w,\varepsilon)$, состоящую из инволюции $w\in\mathfrak{I}_n$ и значений
$\varepsilon_k\in\{+1,-1\}$, прикреплённых к заданным позициям
$k\in\{\ell:w_\ell=\ell\}$.
(Эквивалентным образом, $\varepsilon$ является отображением $\{\ell:w_\ell=\ell\}\to\{+1,-1\}$.)

Удобно представлять $w$ в виде графа $l(w)$ (называемого {\em шаблоном связей}) с $n$ вершинами
$1,2,\ldots,n$ и дугами $(k,w_k)$, соединяющими $k$ и $w_k$ всякий раз, когда $k<w_k$.
 {\em Шаблон связей со знаками $l(w,\varepsilon)$} получается из графа $l(w)$, если пометить каждую вершину
$k\in\{\ell:w_\ell=\ell\}$ знаком $\varepsilon_k$.

Например, следующий шаблон связей со знаками (где нумерация вершин  неявна)
\[
\begin{picture}(140,30)(-20,-5)
\courbe{-16}{32}{0}{8}{10}{12}{16} \courbe{0}{80}{0}{36}{40}{44}{24}
\courbe{96}{112}{0}{103}{104}{105}{6} \put(-18,-3){$\bullet$}
\put(-2,-3){$\bullet$} \put(14,-3){$\bullet$} \put(13,-9){$+$}
\put(30,-3){$\bullet$} \put(46,-3){$\bullet$} \put(45,-9){$-$}
\put(62,-3){$\bullet$} \put(61,-9){$+$} \put(78,-3){$\bullet$}
\put(94,-3){$\bullet$} \put(110,-3){$\bullet$}
\end{picture}
\]
представляет элемент $(w,\varepsilon)$, где
$w=(1;4)(2;7)(8;9)\in\mathfrak{I}_9$ и
$(\varepsilon_3,\varepsilon_5,\varepsilon_6)=(+1,-1,+1)$.

Определим $\varsigma(w,\varepsilon):=\{(p_\ell,q_\ell)\}_{\ell=1}^n$
как пару последовательностей, получаемых следующим образом:
\[
p_\ell\mbox{ (соответственно, $q_\ell$)}=\mbox{(количество знаков $+$
\mbox{ (соответственно, знаков $-$ )}}
\]
\[
\mbox{и дуг проходящих через первые  $\ell$ вершин
$l(w,\varepsilon)$).}
\]

Предполагая $n=p+q$, обозначим через $\mathfrak{I}_n(p,q)$  множество инволюций со знаками сигнатуры $(p,q)$, то есть таких, что $(p_n,q_n)=(p,q)$.
Отметим, что элементы из $\mathfrak{I}_n(p,q)$ совпадают с множеством сигнатуры $(p,q)$
в смысле \cite{MO,Y}.

Например, для вышеуказанной пары
$(w,\varepsilon)$ мы получаем
$(w,\varepsilon)\in\mathfrak{I}_9(5,4)$ и
\[
\varsigma(w,\varepsilon)=\big((0,0),(0,0),(1,0),(2,1),(2,2),(3,2),(4,3),(4,3),(5,4)\big).
\]

\begin{definition}
\label{D2}
\label{DA3}
Для данной инволюции со знаками
$(w,\varepsilon)$ мы будем говорить, что базис $(v_1,\ldots,v_n)$ в $V$  \emph{$(w,\epsilon)$-сопряжён,} если
\[
\delta(v_k)=\left\{\begin{array}{ll} \varepsilon_kv_{w_k}, & \mbox{если
$w_k=k$}
\\ v_{w_k}, & \mbox{если $w_k\not=k$}
\end{array}\right.
\ \mbox{ для всех $k\in\{1,\ldots,n\}$\,.}
\]
Базис
$(v_1,\ldots,v_n)$ такой, что
\[\phi(v_k,v_\ell)=\left\{\begin{array}{ll} \varepsilon_k, & \mbox{если $w_k=\ell=k$} \\ 1, & \mbox{если $w_k=\ell\not=k$} \\ 0, & \mbox{если $w_k\not=\ell$} \end{array}\right.
\ \mbox{ для всех $k,\ell\in\{1,\ldots,n\}$}\] называется {\em
$(w,\varepsilon)$-двойственным}.
Положим
\begin{eqnarray*}
 & \mathcal{O}_{(w,\varepsilon)}:=\{\mathcal{F}(v_1,\ldots,v_n):\mbox{$(v_1,\ldots,v_n)$ --- $(w,\varepsilon)$-сопряжённый базис}\}, \\[1mm]
 & \mathfrak{O}_{(w,\varepsilon)}:=\{\mathcal{F}(v_1,\ldots,v_n):\mbox{$(v_1,\ldots,v_n)$ --- $(w,\varepsilon)$-двойственный базис}\}.
\end{eqnarray*}
\end{definition}



\begin{proposition}
\label{P2}
В дополнение к обозначениям выше, положим $(p,q)=\\ (\dim V_+,\dim V_-)$. Тогда:
\begin{itemize}
\item[\rm (a)] Для каждой $(w,\varepsilon)\in\mathfrak{I}_n(p,q)$ подмножества $\mathcal{O}_{(w,\varepsilon)}$ и $\mathfrak{O}_{(w,\varepsilon)}$ непусты, и
$\mathcal{O}_{(w,\varepsilon)}\cap\mathfrak{O}_{(w,\varepsilon)}= $
\[\{\mathcal{F}(\underline{v}):\mbox{$\underline{v}=(v_k)_{k=1}^n$  $(w,\varepsilon)$-двойственен и $(w,\varepsilon)$-сопряжён}\} \not=\emptyset.\]
\item[\rm (b)]
Для каждого $(w,\varepsilon)\in \mathfrak{I}_n(p,q)$,
\begin{eqnarray*}
 & \mathcal{O}_{(w,\varepsilon)}=\big\{\mathcal{F}\in
X: (\delta(\mathcal{F}),\mathcal{F})\in\mathbb{O}_{w}\mbox{ и
}\varsigma(\delta:\mathcal{F})=\varsigma(w,\varepsilon)\big\}, \\
 &\mathfrak{O}_{(w,\varepsilon)}=\big\{\mathcal{F}\in X:
(\mathcal{F}^\dag,\mathcal{F})\in\mathbb{O}_{w_0w}\mbox{ и
}\varsigma(\phi:\mathcal{F})=\varsigma(w,\varepsilon)\big\}.
\end{eqnarray*}
\item[\rm (c)]
Подмножества $\mathcal{O}_{(w,\varepsilon)}$
($(w,\varepsilon)\in\mathfrak{I}_n(p,q)$) в точности  $K$-орбиты
в $X$. Подмножества $\mathfrak{O}_{(w,\varepsilon)}$
($(w,\varepsilon)\in\mathfrak{I}_n(p,q)$) в точности
$G^0$-орбиты в $X$.
\item[\rm (d)]
Отображение
$\mathcal{O}_{(w,\varepsilon)}\mapsto\mathfrak{O}_{(w,\varepsilon)}$
есть двойственность Мацуки.
\end{itemize}
\end{proposition}

\subsection{Типы B, C, D}

\label{S3.3}

В этом разделе мы будем предполагать, что пространство $V=V_n=\langle e_1,\ldots,e_n\rangle_\mathbb{C}$ обладает симметрической или симплектической формой $\omega$, действие которой на базисе $(e_1,\ldots,e_n)$ описывается матрицей $\Omega$
из (\ref{omega}).
%for instance be described in terms of the matrix $\Omega=\Omega_n$ given by
%\[\begin{array}[t]{c}
%\Omega=I_n \\[1mm]
%\mbox{(type BD1)}
%\end{array},
%\quad
%\begin{array}[t]{c}
%\Omega=\left(\begin{array}{cccc}
%\begin{smallmatrix} 0 & 1 \\ -1 & 0 \end{smallmatrix} & & (0) \\
% & \ddots \\
%(0) & & \begin{smallmatrix} 0 & 1 \\ -1 & 0 \end{smallmatrix}
%\end{array}\right)
%\\
%\mbox{(type C)}
%\end{array},
%\quad
%\begin{array}[t]{c}
%\Omega=\left(\begin{array}{cccc}
%\begin{smallmatrix} 0 & 1 \\ 1 & 0 \end{smallmatrix} & & (0) \\
% & \ddots \\
%(0) & & \begin{smallmatrix} 0 & 1 \\ 1 & 0 \end{smallmatrix}
%\end{array}\right)
%\\
%\mbox{(type D3)}
%\end{array}.\]
Мы рассматриваем группу $G=G(V,\omega)=\{g\in\mathrm{GL}(V):\mbox{$g$ сохраняет $\omega$}\}$ и многообразие изотропных флагов $X_\omega=\{\mathcal{F}\in X:\mathcal{F}^\perp=\mathcal{F}\}$
(см. раздел \ref{S2.2}).

В дополнение мы предположим, что $V$ обладает Эрмитовой формой $\phi$, сопряжением $\delta$ и разложением $V=V_+\oplus V_-$ (как и в разделе \ref{S3.2}) таким, что
\begin{itemize}
\item в случаях BD1 и C2 ограничения формы $\omega$ на $V_+$ и $V_-$ невырождены, то есть, $V_+^\perp=V_-$,
\item в случаях C1 и D3 $V_+$ и $V_-$ --- Лагранжевы по отношению $\omega$, то есть, $V_+^\perp=V_+$ и $V_-^\perp=V_-$.
\end{itemize}
Положим $K:=\{g\in G:g\delta=\delta g\}$ и $G^0:=\{g\in G:\mbox{$g$ сохраняет $\phi$}\}$.

\medskip

\paragraph{\bf Комбинаторные обозначения}
Напомним, что $w_0(k)=n-k+1$.
Пусть $(\eta,\epsilon)\in\{1,-1\}^2$.
Инволюция со знаками $(w,\varepsilon)$ называется {\em $(\eta,\epsilon)$-симметричной}, если выполнены следующие условия:
\begin{itemize}
\item[(i)] $ww_0=w_0w$ (так, что  $w_0$ сохраняет множество $\{\ell:w_\ell=\ell\}$);
\item[(ii)] $\varepsilon_{w_0(k)}=\eta \varepsilon_k$ для всех $k\in\{\ell:w_\ell=\ell\}$;
\end{itemize}
и в случае, когда $\eta\not=\epsilon$:
\begin{itemize}
\item[(iii)]$w_k\not=w_0(k)$ для всех $k$.
\end{itemize}
Предположим, что $n=p+q$, и через $\mathfrak{I}_n^{\eta,\epsilon}(p,q)\subset\mathfrak{I}_n(p,q)$ обозначим подмножество инволюций со знаками сигнатуры $(p,q)$, которые $(\eta,\epsilon)$-симметричны.

В частности, инволюция со знаками $(w,\varepsilon)$  $(1,1)$-симметрична, когда шаблон связей со знаками $l(w,\varepsilon)$ симметричен относительно обращения нумерации вершин;
инволюция со знаками $(w,\varepsilon)$  $(1,-1)$-симметрична, когда  $l(w,\varepsilon)$ симметричен и не имеет симметричных дуг (то есть, соединённых $k$ и $n-k+1$);
$(w,\varepsilon)$  $(-1,-1)$-симметрична,
когда $l(w,\varepsilon)$ антисимметричен
в смысле того, что зеркальный образ $l(w,\varepsilon)$ это шаблон связей со знаками с теми же дугами но противоположными знаками; $(w,\varepsilon)$  $(-1,1)$-симметрична, когда $l(w,\varepsilon)$ антисимметричен и не имеет симметричных дуг. Например:
\begin{eqnarray*}
 & \begin{array}[t]{c}
\begin{picture}(140,30)(-20,-5)
\courbe{-16}{32}{0}{8}{10}{12}{16} \courbe{0}{96}{0}{40}{48}{56}{26}
\courbe{64}{112}{0}{86}{88}{90}{16} \put(-18,-3){$\bullet$}
\put(-2,-3){$\bullet$} \put(14,-3){$\bullet$} \put(13,-9){$+$}
\put(30,-3){$\bullet$} \put(46,-3){$\bullet$} \put(45,-9){$-$}
\put(62,-3){$\bullet$} \put(77,-9){$+$} \put(78,-3){$\bullet$}
\put(94,-3){$\bullet$} \put(110,-3){$\bullet$}
\end{picture}
\\[1mm]
\mbox{$(w,\varepsilon)\in\mathfrak{I}^{1,1}_9(5,4),$}
\end{array}
\qquad
\begin{array}[t]{c}
\begin{picture}(140,30)(-20,-5)
\courbe{-16}{32}{0}{8}{10}{12}{16} \courbe{0}{112}{0}{48}{56}{64}{28}
\courbe{80}{128}{0}{102}{104}{106}{16} \put(-18,-3){$\bullet$}
\put(-2,-3){$\bullet$} \put(14,-3){$\bullet$} \put(13,-9){$+$}
\put(30,-3){$\bullet$} \put(46,-3){$\bullet$} \put(45,-9){$-$}
\put(62,-3){$\bullet$} \put(61,-9){$+$} \put(78,-3){$\bullet$}
\put(94,-3){$\bullet$} \put(93,-9){$-$} \put(110,-3){$\bullet$} \put(126,-3){$\bullet$}
\end{picture}
\\[1mm]
\mbox{$(w,\varepsilon)\in\mathfrak{I}^{-1,-1}_{10}(5,5)$,}
\end{array} \\
 & \begin{array}[t]{c}
\begin{picture}(140,30)(-20,-5)
\courbe{-16}{32}{0}{8}{10}{12}{16}
\courbe{80}{128}{0}{102}{104}{106}{16} \put(-18,-3){$\bullet$}
\put(-2,-3){$\bullet$} \put(-3,-9){$-$} \put(14,-3){$\bullet$} \put(13,-9){$+$}
\put(30,-3){$\bullet$} \put(46,-3){$\bullet$} \put(45,-9){$+$}
\put(62,-3){$\bullet$} \put(61,-9){$+$} \put(78,-3){$\bullet$}
\put(94,-3){$\bullet$} \put(93,-9){$+$} \put(110,-3){$\bullet$} \put(109,-9){$-$} \put(126,-3){$\bullet$}
\end{picture}\\[1mm]
\mbox{$(w,\varepsilon)\in\mathfrak{I}^{1,-1}_{10}(6,4),$}
\end{array}
\qquad
\begin{array}[t]{c}
\begin{picture}(140,30)(-20,-5)
\courbe{-16}{32}{0}{8}{10}{12}{16}
\courbe{80}{128}{0}{102}{104}{106}{16} \put(-18,-3){$\bullet$}
\put(-2,-3){$\bullet$} \put(-3,-9){$-$} \put(14,-3){$\bullet$} \put(13,-9){$+$}
\put(30,-3){$\bullet$} \put(46,-3){$\bullet$} \put(45,-9){$-$}
\put(62,-3){$\bullet$} \put(61,-9){$+$} \put(78,-3){$\bullet$}
\put(94,-3){$\bullet$} \put(93,-9){$-$} \put(110,-3){$\bullet$} \put(109,-9){$+$} \put(126,-3){$\bullet$}
\end{picture}
\\[1mm]
\mbox{$(w,\varepsilon)\in\mathfrak{I}^{-1,1}_{10}(5,5)$.}
\end{array}
\end{eqnarray*}



\begin{proposition}
\label{P3}
Положим $(p,q)=(\dim V_+,\dim V_-)$ (так, что $p=q=\frac{n}{2}$ для типов C1 и D3).
Зададим
$(\eta,\epsilon)=(1,1)$ для типа BD1,
$(\eta,\epsilon)=(1,-1)$ для типа C2,
$(\eta,\epsilon)=(-1,-1)$ для типа C1 и
$(\eta,\epsilon)=(-1,1)$ для типа D3.
\begin{itemize}
\item[(a)] Для каждого $(w,\varepsilon)\in\mathfrak{I}_n^{\eta,\epsilon}(p,q)$,
рассматривая базис $\underline{v}=(v_1,\ldots,v_n)$ в $V$ такой, что
\begin{equation}
\label{3.3.18}
\omega(v_k,v_\ell)=\left\{\begin{array}{ll} 0, & \mbox{если $\ell\not=n-k+1$} %\\ \pm 1 & \mbox{if $\ell=n-k+1$.}
\\
1, & \mbox{если $\ell=n-k+1$ и $w_k,w_\ell\in[k,\,\ell]$ ($k\leq\ell$)} \\
\epsilon, & \mbox{если $\ell=n-k+1$ и $w_k,w_\ell\in[\ell,\,k]$ ($\ell\leq k$)} \\
\eta, & \mbox{если $\ell=n-k+1$ и $k,\ell\in]w_k,\,w_\ell[$} \\
\eta\epsilon, & \mbox{если $\ell=n-k+1$ и $k,\ell\in]w_\ell,\,w_k[$,} \\
\end{array}\right.
\end{equation}
мы получаем
\begin{eqnarray*}
 & & \mathcal{O}^{\eta,\epsilon}_{(w,\varepsilon)}:=\\
 & &\mathcal{O}_{(w,\varepsilon)}\cap X_\omega=\{\mathcal{F}(\underline{v}):\underline{v}\mbox{  $(w,\varepsilon)$-сопряжён и удовлетворяет (\ref{3.3.18})}\}\not=\emptyset, \\
 & & \mathfrak{O}^{\eta,\epsilon}_{(w,\varepsilon)}:=\\
 & &\mathfrak{O}_{(w,\varepsilon)}\cap X_\omega=\{\mathcal{F}(\underline{v}):\underline{v}\mbox{ $(w,\varepsilon)$-двойственен и удовлетворяет (\ref{3.3.18})}\}\not=\emptyset, \\
 & & \mathcal{O}_{(w,\varepsilon)}^{\eta,\epsilon}\cap\mathfrak{O}_{(w,\varepsilon)}^{\eta,\epsilon}=\\
 & &=\{\mathcal{F}(\underline{v}):\underline{v}\mbox{ $(w,\varepsilon)$-сопряжён, $(w,\varepsilon)$-двойственен и удовлетворяет (\ref{3.3.18})}\}\not=\emptyset.
\end{eqnarray*}
\item[(b)] Подмножества $\mathcal{O}_{(w,\varepsilon)}^{\eta,\epsilon}$ ($(w,\varepsilon)\in\mathfrak{I}_n^{\eta,\epsilon}(p,q)$) суть в точности $K$-орбиты в $X_\omega$.
Подмножества $\mathfrak{O}_{(w,\varepsilon)}^{\eta,\epsilon}$ ($(w,\varepsilon)\in\mathfrak{I}_n^{\eta,\epsilon}(p,q)$) суть в точности $G^0$-орбиты в $X_\omega$.
\item[(c)] Отображение $\mathcal{O}_{(w,\varepsilon)}^{\eta,\epsilon}\mapsto \mathfrak{O}_{(w,\varepsilon)}^{\eta,\epsilon}$ есть двойственность Мацуки.
\end{itemize}
\end{proposition}

\subsection{Замечания}

Положим $X_0:=X$ для типа A и $X_0:=X_\omega$ для типов B, C, D.

\begin{remark}
\label{R3.4.1}
Характеризация $K$-орбит в разделах \ref{P1} --- \ref{P3} может быть получена следующим единым образом. Если $\mathcal{F}\in X$, мы будем писать
$\sigma(\mathcal{F})=\mathcal{F}^\perp$ для типов A1--A2 и
$\sigma(\mathcal{F})=\delta(\mathcal{F})$ для типов A3, BD1, C1--C2, D3. Пусть $P\subset
G$ --- параболическая подгруппа, содержащая $K$ и минимальная с этим свойством. Два флага $\mathcal{F}_1,\mathcal{F}_2\in X_0$ принадлежат одной  $K$-орбите тогда и только тогда
$(\sigma(\mathcal{F}_1),\mathcal{F}_1)$ и
$(\sigma(\mathcal{F}_2),\mathcal{F}_2)$ принадлежат одной  $P$-орбите   для диагонального действия $P$ на $X_0\times X_0$.
\end{remark}


\begin{remark}
	
	(Открытые $K$-орбиты.)
\label{R-open}
	В обозначениях из замечания \ref{R3.4.1} отображение $\sigma_0:X_0\to X\times X$, $\mathcal{F}\mapsto(\sigma(\mathcal{F}),\mathcal{F})$ является замкнутым вложением.
%For every $K$-orbit $\mathcal{O}\subset X_0$
%there is a unique $G$-orbit $\mathbb{O}_w\subset X\times X$
%such that $\mathcal{O}\subset\sigma_0^{-1}(\mathbb{O}_w)$.

Для типов A и C многообразие флагов $X_0$ неприводимо.
В частности, существует единственная $G$-орбита $\mathbb{O}_w\subset X\times X$
такая, что пересечение $\mathbb{O}_w\cap \sigma_0(X_0)$ открыто в $\sigma_0(X_0)$;
оно соответствует элементу $w\in\mathfrak{S}_n$ максимальному для порядка Брюа такому, что $\mathbb{O}_w$ пересекается с $\sigma_0(X_0)$.
Во всех случаях найдётся единственная  $K$-орбита $\mathcal{O}\subset X_0$ такая, что $\sigma_0(\mathcal{O})\subset\mathbb{O}_w$, поэтому это (единственная) открытая $K$-орбита в $X_0$.
Таким образом, мы получаем следующий список открытых $K$-орбит для типов A1--A3, C1--C2:
\begin{itemize}
\item[\rm A1:] $\mathcal{O}_\mathrm{id}$;
\item[\rm A2:] $\mathcal{O}_{v_0}$, где $v_0=(1;2)(3;4)\cdots(n-1;n)$;
\item[\rm A3:] $\mathcal{O}_{(w_0^{(t)},\varepsilon)}$, где
$t=\min\{p,q\}$, $\varepsilon\equiv\mathrm{sign}(p-q)$
и $w^{(t)}_0=\prod\limits_{k=1}^t(k;n-k+1)$;
\item[\rm C1:] $\mathcal{O}^{-1,-1}_{(w_0,\emptyset)}$;
\item[\rm C2:] $\mathcal{O}^{1,-1}_{(\hat{w}_0^{(t)},\varepsilon)}$,
где $t=\min\{p,q\}$, $\varepsilon\equiv\mathrm{sign}(p-q)$ и $\hat{w}_0^{(t)}=v_0^{(t)}w_0^{(t)}v_0^{(t)}$,
где $v_0^{(t)}=(1;2)(3;4)\cdots(t-1;t)$.
\end{itemize}

Если $n=\dim V$ чётно и форма $\omega$ ортогональна, то многообразие $X_\omega$ иимеет две связные компоненты.
В самом деле, для каждого изотропного флага $\mathcal{F}=(F_k)_{k=0}^n\in X_\omega$ существет единственный изотропный флаг $\tilde{\mathcal{F}}=(\tilde{F}_k)_{k=0}^n\in X_\omega$
такой, что $F_k=\tilde{F}_k$ для всех $k\not=m:=\frac{n}{2}$, $\tilde{F}_m\not=F_m$.
Тогда отображение $\tilde{I}:\mathcal{F}\mapsto\tilde{\mathcal{F}}$ --- автоморфизм
 $X_\omega$, который отображает одну связную компоненту $X_\omega$ на другую.
Если $\mathcal{F}=\mathcal{F}(v_1,\ldots,v_n)$ для базиса $\underline{v}=(v_1,\ldots,v_n)$ такого, что
\[\omega(v_k,v_\ell)\not=0\Leftrightarrow \ell=n-k+1,\]
то
$\tilde{I}(\mathcal{F}(\underline{v}))=\mathcal{F}(\underline{\tilde{v}})$,
где $\underline{\tilde{v}}$ --- базис полученный $\underline{v}$ перестановкой двух средних векторов $v_m,v_{m+1}$.
Если $\underline{v}$ --- $(w,\varepsilon)$-сопряжённый, то $\underline{\tilde{v}}$ --- $\tilde{i}(w,\varepsilon)$-сопряжённый, где
$\tilde{i}(w,\varepsilon):=\big((m;m+1)w(m;m+1),\varepsilon\circ(m;m+1)\big)$.
Следовательно $\tilde{I}$ отображает  $K$-орбиту $\mathcal{O}^{\eta,\epsilon}_{(w,\varepsilon)}$
на $\mathcal{O}^{\eta,\epsilon}_{\tilde{i}(w,\varepsilon)}$.

Для типа D3
$X_\omega$ имеет в точности две открытые $K$-орбиты.
Действительно, перестановка $w=\hat{w}_0:=w_0v_0$
максимальна для порядка такого, что $\mathbb{O}_w\cap\sigma_0(X_0)$ непусто,
следовательно, $\sigma_0^{-1}(\mathbb{O}_{\hat{w}_0})$ открыта.
Перестановка $\hat{w}_0$ не имеет неподвижных точек, если $m:=\frac{n}{2}$ чётное; если $m:=\frac{n}{2}$ нечётное, то $\hat{w}_0$ фиксирует $m$ и $m+1$.
В первом случае
$\sigma_0^{-1}(\mathbb{O}_{\hat{w}_0})=\mathcal{O}^{-1,1}_{(\hat{w}_0,\emptyset)}$ ---
единая $K$-орбита, и
$\tilde{I}(\mathcal{O}^{-1,1}_{(\hat{w}_0,\emptyset)})=\mathcal{O}^{-1,1}_{\tilde{i}(\hat{w}_0,\emptyset)}$ ---
вторая открытая $K$-орбита.
Во втором случае
$\sigma_0^{-1}(\mathbb{O}_{\hat{w}_0^{(m-1)}})=\mathcal{O}^{-1,1}_{(\hat{w}_0,\varepsilon)}\cup\mathcal{O}^{-1,1}_{(\hat{w}_0,\tilde\varepsilon)}$, где $(\varepsilon_m,\varepsilon_{m+1})=(\tilde\varepsilon_{m+1},\tilde\varepsilon_m)=(+1,-1)$, это объединение двух различных $K$-орбит, являющихся образами друг друга под действием  $\tilde{I}$.

В случае BD1 многообразие $X_\omega$ может быть разложимым,
но $w=w^{(t)}_0$, для $t:=\min\{p,q\}$, это единственный максимальный элемент в $\mathfrak{S}_n$
такой, что $\mathbb{O}_w\cap\sigma_0(X_0)$ непусто.
Тогда $\sigma_0^{-1}(\mathbb{O}_w)$
состоит из единой $\tilde{I}$-инвариантной открытой $K$-орбиты, обозначаемой
$\mathcal{O}^{1,1}_{(w^{(t)}_0,\varepsilon)}$ для $\varepsilon\equiv\mathrm{sign}(p-q)$.
Следовательно многообразие флагов $X_\omega$  имеет единственную открытую орбиту  $K$-орбиту (которая несвязана в случае, когда $n$ чётное).
\end{remark}



\begin{remark}
	
	(Замкнутые $K$-орбиты.)
\label{R-closed}
Мы будем использовать обозначения из замечаний \ref{R3.4.1}--\ref{R-open}.
Как видно из предложений \ref{P1}--\ref{P3}, в каждом случае можно найти единственную перестановку $w_{\mathrm{min}}\in\mathfrak{S}_n$ такую, что $\mathbb{O}_{w_{\mathrm{min}}}\cap\sigma_0(X_0)$ замкнуто;
в самом деле, $w_{\mathrm{min}}=\mathrm{id}$ кроме типа BD1 для $p,q$ нечётных:
в этом случае$w_{\mathrm{min}}=(\frac{n}{2};\frac{n}{2}+1)$.
Для каждой $K$-орбиты $\mathcal{O}\subset X_0$ имеет место следующая эквивалентность:
\[
\mbox{$\mathcal{O}$ замкнута}\quad \Leftrightarrow\quad \sigma_0(\mathcal{O})\subset\mathbb{O}_{w_{\mathrm{min}}}
\]
(см. \cite{Brion,Richardson-Springer}). Ввиду этой эквивалентности, мы получаем следующий полный список замкнутых $K$-орбит  $X_0$ для разных типов.
Для типов A1 и A2, $\mathcal{O}_{w_0}$ --- единственная замкнутая $K$-орбита.
Для типа A3 замкнутые $K$-орбиты суть в точности орбиты
$\mathcal{O}_{(\mathrm{id},\varepsilon)}$ для всех пар вида $(\mathrm{id},\varepsilon)\in\mathfrak{I}_n(p,q)$;
существует $\binom{n}{p}$ таких орбит.
Для типов B, C, D, замкнутые $K$-орбиты это орбиты
$\mathcal{O}^{\eta,\epsilon}_{(\mathrm{id},\varepsilon)}$ для всех пар вида $(\mathrm{id},\varepsilon)\in\mathfrak{I}^{\eta,\epsilon}_n(p,q)$,
кроме случая BD1, при условии $n=:2m$ чётное и $p,q$ нечётные;
в этом случае замкнутые $K$-орбиты
это орбиты $\mathcal{O}^{1,1}_{((m;m+1),\varepsilon)}$
для всех пар вида $((m;m+1),\varepsilon)\in\mathfrak{I}^{1,1}_n(p,q)$.
Существует $\binom{\lfloor\frac{p}{2}\rfloor+\lfloor\frac{q}{2}\rfloor}{\lfloor\frac{p}{2}\rfloor}$ замкнутых орбит в случаях BD1 и C2, и $2^{\frac{n}{2}}$ замкнутых орбит для типов C1 и D3.
\end{remark}


\begin{remark} Предложения
\ref{P1}--\ref{P3} показывают, в частности, что {\em специальные элементы $X_0$},  в понимании Мацуки \cite{Matsuki, Matsuki3}, суть в точности флаги $\mathcal{F}\in X_0$ вида
$\mathcal{F}=\mathcal{F}(v_1,\ldots,v_n)$, где $(v_1,\ldots,v_n)$ ---
одновременно двойственный и сопряжённый базис пространства $V$, относительно некоторой инволюции $w\in\mathfrak{I}_n^\epsilon$ для типов A1
и A2, и  некоторой инволюции со знаками
$(w,\varepsilon)\in\mathfrak{I}_n(p,q)$ для типов A3, B--D. В самом деле,  согласно \cite{Matsuki, Matsuki3}, множество
$\mathcal{S}\subset X_0$ специальных элементов совпадает с
\[\bigcup_{\mathcal{O}\in X_0/K}\mathcal{O}\cap\Xi(\mathcal{O}),\]
где отображение $X_0/K\to X_0/G^0$, $\mathcal{O}\mapsto\Xi(\mathcal{O})$
это двойственность Мацуки.
\end{remark}

\subsection{Доказательства}

\label{S3.5}

\begin{proof}[Доказательство предложения \ref{P1}\,{\rm (a)}]
Мы будем писать $w=(a_1;b_1)\cdots(a_m;b_m)$ с $a_1<\ldots<a_m$ и
$a_k<b_k$ для всех $k$; положим $c_1<\ldots<c_{n-2m}$ будут элементами множества $\{k:w_k=k\}$.
Для типа A2 мы имеем $n=2m$, и $(e_1,\ldots,e_n)$ есть одновременно
$(1;2)(3;4)\cdots(n-1;n)$-двойственный базис и
$(1;2)(3;4)\cdots(n-1;n)$-сопряжённый базис; тогда базис
$\left\{e'_1,\ldots,e'_n\right\}$, где
\[e'_{a_\ell}=e_{2\ell-1}\quad\mbox{и}\quad e'_{b_\ell}=e_{2\ell}\quad\mbox{для всех $\ell\in\{1,\ldots,m\}$},\]
является  одновременно $w$-двойственным и $w$-сопряжённым.
Для типа A1,
заменяя $e_\ell$ и $e_{\ell^*}$ на $\frac{e_\ell+e_{\ell^*}}{\sqrt{2}}$ и $\frac{e_\ell-e_{\ell^*}}{i\sqrt{2}}$ в тех случаях, когда $\ell<\ell^*$,
мы можем считать, что
базис $(e_1,\ldots,e_n)$  одновременно $\mathrm{id}$-двойственен
и $\mathrm{id}$-сопряжён. Для всех $\ell\in\{1,\ldots,m\}$ и
$k\in\{1,\ldots,n-2m\}$,
\[e'_{a_\ell}=\frac{e_{2\ell-1}+ie_{2\ell}}{\sqrt{2}},\quad e'_{b_\ell}=\frac{e_{2\ell-1}-ie_{2\ell}}{\sqrt{2}}\quad\mbox{и}\quad e'_{c_k}=e_{2m+k}.\]
Тогда $(e'_1,\ldots,e'_n)$ --- одновременно  $w$-двойственный и
$w$-сопряжённый базис.
В обоих случаях мы заключаем, что
\begin{equation}
\label{3.1.4}
\emptyset\not=\{\mathcal{F}(v_1,\ldots,v_n):\mbox{$(v_1,\ldots,v_n)$ --- $w$-двойственный и $w$-сопряжённый}\}\subset\mathcal{O}_w\cap\mathfrak{O}_w.
\end{equation}


Теперь проверим обратное включение. Предположим, что
$\mathcal{F}\:=\:(F_0,\ldots,F_n)\in\mathcal{O}_w \cap \mathfrak{O}_w$.
Пусть $(v_1,\:\ldots,\:v_n)$ --- $w$-двойственный базис такой, что $ \mathcal{F} = \mathcal{F}(v_1,\ldots,v_n)$.
Поскольку $\mathcal{F}\in\mathfrak{O}_w$, мы имеем
\begin{equation}
\label{1-new}
w_k=\min\{\ell=1,\ldots,n:\gamma(F_k)\cap F_\ell\not=\gamma(F_{k-1})\cap F_\ell\}.
\end{equation}
Для всех $\ell\in\{0,\ldots,n\}$ мы построим  $w$-двойственный базис
$(v_1^{(\ell)},\ldots,v_n^{(\ell)})$ в $V$ такой, что
\begin{eqnarray}
\label{2-new}
& F_k=\langle v_1^{(\ell)},\ldots,v_k^{(\ell)}\rangle_\mathbb{C}
\quad\mbox{
для всех $k\in\{1,\ldots,n\}$}
\end{eqnarray}
и
\begin{eqnarray}
\label{3-new}
&
\gamma(v_k^{(\ell)})=\left\{\begin{array}{ll}
\epsilon v_{w_k}^{(\ell)}, & \mbox{если $w_k\geq k$,} \\
v_{w_k}^{(\ell)}, & \mbox{если $w_k<k$}
\end{array}\right.
\quad\mbox{
для всех $k\in\{1,\ldots,\ell\}$.}
\end{eqnarray}
Это будет означать, что
$\mathcal{F}=\mathcal{F}(v_1^{(n)},\ldots,v_n^{(n)})$ для базиса $(v_1^{(n)},\ldots,v_n^{(n)})$ одновременно $w$-двойственного и $w$-сопряжённого,
то есть, полностью докажет (a).

Наша конструкция будет строиться по индукции, начиная с $(v_1^{(0)},\ldots,v_n^{(0)})=(v_1,\ldots,v_n)$.
Пусть $\ell\in\{1,\ldots,n\}$, и предположим, что базис $(v_1^{(\ell-1)},\ldots,v_n^{(\ell-1)})$ построен.
Мы выделяем три случая.

\smallskip
\noindent
{\it Случай 1: $w_\ell<\ell$.}

Неравенство $w_\ell<\ell=w(w_\ell)$ влечёт $\gamma(v_{w_\ell}^{(\ell-1)})=\epsilon v_{\ell}^{(\ell-1)}$, откуда
$\gamma(v_{\ell}^{(\ell-1)})=v_{w_\ell}^{(\ell-1)}$, так как $\gamma^2=\epsilon\mathrm{id}$.
Поэтому базис $(v_1^{(\ell)},\ldots,v_n^{(\ell)}):=(v_1^{(\ell-1)},\ldots,v_n^{(\ell-1)})$ удовлетворяет условиям (\ref{2-new}) и (\ref{3-new}).

\smallskip
\noindent
{\it Случай 2: $w_\ell=\ell$.}

Этот случай имеет место только для типа A1.
С одной стороны, (\ref{1-new}) даёт
\[\gamma(v_\ell^{(\ell-1)})\in\langle v_1^{(\ell-1)},\ldots,v_\ell^{(\ell-1)},v_{w_1}^{(\ell-1)},\ldots,v_{w_{\ell-1}}^{(\ell-1)}\rangle_\mathbb{C}.\]
С другой стороны, так как базис $(v_1^{(\ell-1)},\ldots,v_n^{(\ell-1)})$ является $w$-двойственным,
мы получаем
\[v_\ell^{(\ell-1)}\in\langle v_1^{(\ell-1)},\ldots,v_{\ell-1}^{(\ell-1)},v_{w_1}^{(\ell-1)},\ldots,v_{w_{\ell-1}}^{(\ell-1)}\rangle_\mathbb{C}^\perp\,.\]
Следовательно, поскольку $\gamma$ сохраняет ортогональность по отношению к форме $\omega$, то
\begin{eqnarray*}
\gamma(v_\ell^{(\ell-1)}) & \in & \langle \gamma(v_1^{(\ell-1)}),\ldots,\gamma(v_{\ell-1}^{(\ell-1)}),\gamma(v_{w_1}^{(\ell-1)}),\ldots,\gamma(v_{w_{\ell-1}}^{(\ell-1)})\rangle_\mathbb{C}^\perp \\
 & & =\langle v_1^{(\ell-1)},\ldots,v_{\ell-1}^{(\ell-1)},v_{w_1}^{(\ell-1)},\ldots,v_{w_{\ell-1}}^{(\ell-1)}\rangle_\mathbb{C}^\perp.
\end{eqnarray*}
Всё это вместе  указывает на существование ненулевого комплексного  числа  $\lambda$ такого, что $\gamma(v_\ell^{(\ell-1)})=\lambda v_\ell^{(\ell-1)}$.
Так как $\gamma$ --- инволюция, мы заключаем, что $\lambda\in\{+1,-1\}$.
В дополнение мы знаем, что
\[\lambda=\omega(\gamma(v_\ell^{(\ell-1)}),v_\ell^{(\ell-1)})={}^t\overline{v_\ell^{(\ell-1)}} v_\ell^{(\ell-1)}\in\mathbb{R}^+.\]
Следовательно $\gamma(v_\ell^{(\ell-1)})=v_\ell^{(\ell-1)}$, и мы можем положить
$$(v_1^{(\ell)},\ldots,v_n^{(\ell)}):=(v_1^{(\ell-1)},\ldots,v_n^{(\ell-1)}).$$

\smallskip
\noindent
{\it Случай 3: $w_\ell>\ell$.}

Из (\ref{1-new}) мы получаем
\[\gamma(v_\ell^{(\ell-1)})\in
\langle v_k^{(\ell-1)}:1\leq k\leq w_\ell\rangle_\mathbb{C}+\langle v_{w_k}^{(\ell-1)}:1\leq k\leq \ell-1\rangle_\mathbb{C}.
\]
С другой стороны, рассуждая, как и в случае 2, мы видим, что
\[\gamma(v_\ell^{(\ell-1)})\in\langle v_1^{(\ell-1)},\ldots,v_{\ell-1}^{(\ell-1)},v_{w_1}^{(\ell-1)},\ldots,v_{w_{\ell-1}}^{(\ell-1)}\rangle_\mathbb{C}^\perp.\]
Поэтому выполняется равенство
\begin{equation}
\label{4-new}
\gamma(v_\ell^{(\ell-1)})=\sum_{k\in I}\lambda_kv_k^{(\ell-1)}
\quad\mbox{для некоторых $\lambda_k\in\mathbb{C}$ ,}
\end{equation}
где $I:=\{k:\ell\leq k\leq w_\ell\mbox{ and }\ell\leq w_k\}\subset\hat{I}:=\{k:\ell\leq k\mbox{ and }\ell\leq w_k\}$.
Используя (\ref{4-new}), факт, что базис $(v_1^{(\ell-1)},\ldots,v_n^{(\ell-1)})$  $w$-двойственен, и определения $\omega$ и $\gamma$, мы видим, что
\begin{equation}
\label{5-new}
\lambda_{w_\ell}=\omega(v_\ell^{(\ell-1)},\gamma(v_\ell^{(\ell-1)}))=\epsilon \mbox{\:}^tv_\ell^{(\ell-1)}\overline{v_\ell^{(\ell-1)}}=\epsilon\alpha,
\end{equation}
где $\alpha\in\mathbb{R}$, $\alpha>0$.
Положим
\begin{eqnarray*}
 && v_\ell^{(\ell)}:=\frac{1}{\sqrt{\alpha}}v_\ell^{(\ell-1)},\quad
v_{w_\ell}^{(\ell)}:=\frac{\epsilon}{\sqrt{\alpha}}\gamma(v_\ell^{(\ell-1)}),\\
 && v_k^{(\ell)}:=v_k^{(\ell-1)}-\frac{\omega(v_k^{(\ell-1)},\gamma(v_\ell^{(\ell-1)}))}{\lambda_{w_\ell}}v_\ell^{(\ell-1)}\ \mbox{ для всех $k\in \hat{I}\setminus\{\ell,w_\ell\}$}, \\
&& v_k^{(\ell)}:=v_k^{(\ell-1)}\ \mbox{ для всех $k\in\{1,\ldots,n\}\setminus \hat{I}$.}
\end{eqnarray*}
Основываясь на равенствах (\ref{4-new}) и (\ref{5-new}), легко проверить, что $(v_1^{(\ell)},\ldots,v_n^{(\ell)})$ является  $w$-двойственным базисом,  удовлетворяющим (\ref{2-new}) и (\ref{3-new}).
Это завершает рассмотрение случая 3.
\end{proof}

\begin{proof}[Доказательство предложений \ref{P1}\,{\rm (b)}--{\rm (d)}]
Пусть $\mathcal{F}\in\mathcal{O}_w$, то есть,
$\mathcal{F}=\mathcal{F}(v_1,\ldots,v_n)$ для некоторого $w$-двойственного базиса
$(v_1,\ldots,v_n)$ в $V$. Из определения
$w$-двойственного базиса мы видим, что
\begin{eqnarray*}
\langle v_1,\ldots,v_{n-k}\rangle_\mathbb{C}^\perp & = & \langle v_j:w_j\notin\{1,\ldots,n-k\}\rangle_\mathbb{C} \\
 & = & \langle v_j:w_j\in\{n-k+1,\ldots,n\}\rangle_\mathbb{C} \\
 & = & \langle v_j:(w_0w)_j\in\{1,\ldots,k\}\rangle_\mathbb{C}\,.
\end{eqnarray*}
Следовательно
\[\dim\langle v_1,\ldots,v_{n-k}\rangle_\mathbb{C}^\perp\cap\langle v_1,\ldots,v_\ell\rangle_\mathbb{C}=\big|\big\{j\in\{1,\ldots,\ell\}:(w_0w)_j\in\{1,\ldots,k\}\big\}\big|\]
для всех $k,\ell\in\{1,\ldots,n\}$, что даёт равенство
$w(\mathcal{F}^\perp,\mathcal{F})=w_0w$, и поэтому включение
\begin{equation}
\label{1}
\mathcal{O}_w\subset\{\mathcal{F}\in
X:(\mathcal{F}^\perp,\mathcal{F})\in\mathbb{O}_{w_0w}\}.
\end{equation}

Пусть $\mathcal{F}=\mathcal{F}(v_1,\ldots,v_n)\in\mathfrak{O}_w$ для
$w$-сопряжённого базиса $(v_1,\ldots,v_n)$ в $V$. Из определения $w$-сопряжённого базиса мы получаем
\[\gamma(\langle v_1,\ldots,v_k\rangle_\mathbb{C})=\langle v_{w_j}:j\in\{1,\ldots,k\}\rangle_\mathbb{C}\,.\]
Поэтому
\[\dim\gamma(\langle v_1,\ldots,v_k\rangle_\mathbb{C})\cap\langle v_1,\ldots,v_\ell\rangle_\mathbb{C}=\]

\[\big|\big\{j\in\{1,\ldots,\ell\}:w^{-1}_j\in\{1,\ldots,k\}\big\}\big|\]
для всех $k,\ell\in\{1,\ldots,n\}$,
 откуда
$w(\gamma(\mathcal{F}),\mathcal{F})=w^{-1}=w$ (где $w$ --- инволюция).\\ Это влечёт включение
\begin{equation}
\label{2} \mathfrak{O}_w\subset\{\mathcal{F}\in
X:(\gamma(\mathcal{F}),\mathcal{F})\in\mathbb{O}_{w}\}.
\end{equation}

Ясно, что группа $K$ действует транзитивно на множестве
$w$-двойственных базисов, следовательно, $\mathcal{O}_w$ является  $K$-орбитой. Более того,
(\ref{1}) влечёт то, что орбиты $\mathcal{O}_w$ (для
$w\in\mathfrak{I}_w^\epsilon$) попарно различны. Аналогично, подмножества $\mathfrak{O}_w$ (для $w\in\mathfrak{I}_w^\epsilon$) суть попарно различные $G^0$-орбиты.

Обозначим через $L_k$ матрицу размера $k\times k$ с $1$ на антидиагонали и $0$ в остальных местах. Пусть
$\underline{v}=(v_1,\ldots,v_n)$ является $w_0$-двойственным базисом, другими словами,
\[
\left\{
\begin{array}{ll}
\omega(v_k,v_{n+1-k})=\left\{\begin{array}{ll}
1, & \mbox{если $k\leq\frac{n+1}{2}$} \\[1mm] \epsilon, & \mbox{если $k>\frac{n+1}{2}$}
\end{array}\right. \\[4mm]
\omega(v_k,v_\ell)=0,\quad\mbox{если $\ell\not=n+1-k$;}
\end{array}
\right.
\]
поэтому  $L:=\left(\omega(v_k,v_\ell)\right)_{1\leq k,\ell\leq n}$ это следующая матрица:
\[
L=L_n\quad \mbox{(тип A1)}\quad\mbox{или}\quad L=\left(\begin{matrix} 0 & L_m \\ -L_m & 0 \end{matrix}\right)\quad\mbox{(тип A2, $n=2m$)}.
\]

Флаг $\mathcal{F}_0:=\mathcal{F}(v_1,\ldots,v_n)$ удовлетворяет условию $\mathcal{F}_0^\perp=\mathcal{F}_0$. По Ричардсону-Спрингеру \cite{Richardson-Springer} каждая $K$-орбита
$\mathcal{O}\subset X$ содержит элемент вида
$g\mathcal{F}_0$ с  $g\in G$ такой, что
$h:=L{}^t[g]_{\underline{v}}L^{-1}[g]_{\underline{v}}\in N$, где
$[g]_{\underline{v}}$ обозначает матрицу $g$ в базисе
$\underline{v}$ и $N$ обозначает группу обратимых $n\times
n$-матриц с ровно одним ненулевым коэффициентом в каждой строке и кажом столбце. Заметим, что матрица
$Lh={}^t[g]_{\underline{v}}L[g]_{\underline{v}}$ тоже принадлежит   $N$
(так же, как и $L$ ) и является симметрической в случае A1 и антисимметрической в
случае A2. Вследствие этого, существует $w\in\mathfrak{I}_n$ и константы
$t_1,\ldots,t_n\in\mathbb{C}^*$ такие, что матрица
$Lh=:\left(a_{k,\ell}\right)_{1\leq k,\ell\leq n}$ состоит из коэффициентов:
\[
a_{k,\ell}=0,\ \mbox{если $\ell\not=w_k$,}\qquad
a_{k,w_k}=\left\{\begin{array}{ll} t_k, & \mbox{если $w_k\geq k$}
\\ \epsilon t_k, & \mbox{если $w_k\leq k$.}
\end{array}\right.
\]
Так как $\epsilon=-1$ для типа A2, должно выполняться $w_k\not=k$ для всех $k$,
поэтому $w\in\mathfrak{I}'_n$. Следовательно, в обоих случаях
$w\in\mathfrak{I}^\epsilon_n$.
Для всех $k\in\{1,\ldots,n\}$, мы
выберем $s_k=s_{w_k}\in\mathbb{C}^*$ такие, что $s_k^{-2}=t_k$
(отметим, что $t_{w_k}=t_k$).
Таким образом,
\[g\mathcal{F}_0=\mathcal{F}(s_1gv_1,\ldots,s_ngv_n)\,,\]
и для всех $k,\ell\in\{1,\ldots,n\}$ мы имеем
\[
\omega(s_kgv_k,s_\ell g
v_\ell)=s_ks_\ell\omega(gv_k,gv_\ell)=s_ks_\ell
a_{k,\ell}=\left\{\begin{array}{ll} 1, & \mbox{если $\ell=w_k\geq k$}
\\ \epsilon, & \mbox{если $\ell=w_k< k$} \\ 0, & \mbox{если
$\ell\not=w_k$}.
\end{array}\right.
\]
Поэтому $g\mathcal{F}_0\in\mathcal{O}_w$. Это даёт
$\mathcal{O}=\mathcal{O}_w$.

Мы доказали, что подмножества $\mathcal{O}_w$ (для
$w\in\mathfrak{I}^\epsilon_w$) суть в точности $K$-орбиты в $X$. В частности, $X=\bigcup_{w\in\mathfrak{I}^\epsilon_w}\mathcal{O}_w$ так, что включение (\ref{1}) на самом деле является равенством. По двойственности Мацуки число $G^0$-орбит в $X$ равно числу
$K$-орбит, поэтому подмножества $\mathfrak{O}_w$ (для
$w\in\mathfrak{I}^\epsilon_w$) суть в точности $G^0$-орбиты в $X$.
Тем самым в (\ref{2})  имеет место равенство. Таким образом, мы доказали части
{\rm (b)} и {\rm (c)} утверждения.

Из части (a) следует что,  для каждого $w\in\mathfrak{I}_n^\epsilon$, пересечение $\mathcal{O}_w\cap\mathfrak{O}_w$ непусто и состоит из единой $K\cap G^0$-орбиты.
Это показывает, что орбита $\mathfrak{O}_w$ двойственна по Мацуки орбите $\mathcal{O}_w$ (см. \cite{Matsuki3}), и часть {\rm (d)} утверждения также доказана.
\end{proof}

\begin{proof}[Доказательство предложения \ref{P2}\,{\rm (a)}]
Мы представим $w$ в виде произвдения попарно коммутирующих транспозиций
$w=(a_1;b_1)\cdots(a_m;b_m)$,
и пусть $c_{m+1}<\ldots<c_p$ --- элементы в $\{k:w_k=k,\ \varepsilon_k=+1\}$
и $d_{m+1}<\ldots<d_q$ - элементы в $\{k:w_k=k,\ \varepsilon_k=-1\}$.
Представим $\{e_1,\ldots,e_n\}=\{e_1^+,\ldots,e_p^+\}\cup\{e_1^-,\ldots,e_q^-\}$
так, чтобы $V_+=\langle e_\ell^+:\ell=1,\ldots,p\rangle_\mathbb{C}$
и $V_-=\langle e_\ell^-:\ell=1,\ldots,q\rangle_\mathbb{C}$.
Полагая
\begin{eqnarray*}
& v_{a_k}:=\frac{e^+_k+e^-_k}{\sqrt{2}}\,,\,\,\ v_{b_k}:=\frac{e^+_k-e^-_k}{\sqrt{2}}\ \mbox{ для всех $k\in\{1,\ldots,m\}$,} \\
& v_{c_k}:=e^+_k\ \mbox{ для всех $k\in\{m+1,\ldots,p\}$, \ и }\ v_{d_k}:=e^-_k\ \mbox{ для всех $k\in\{m+1,\ldots,q\}$,}
\end{eqnarray*}
легко видеть, что $(v_1,\ldots,v_n)$ ---  
$(w,\varepsilon)$-двойственный и $(w,\varepsilon)$-сопряжённый базис в $V$.  Следовательно
\begin{equation}
\label{3.5.14}
\emptyset\not=\{\mathcal{F}(\underline{v}):\mbox{$\underline{v}$  $(w,\varepsilon)$-двойственен и $(w,\varepsilon)$-сопряжён}\}\subset \mathcal{O}_{(w,\varepsilon)}\cap\mathfrak{O}_{(w,\varepsilon)}.
\end{equation}

	Чтобы показать обратное  включение, рассмотрим
$\mathcal{F}=(F_0,\ldots,F_n)\in \mathcal{O}_{(w,\varepsilon)}\cap\mathfrak{O}_{(w,\varepsilon)}$.
С одной стороны, так как $\mathcal{F}\in\mathfrak{O}_{(w,\varepsilon)}$, существует  $(w,\varepsilon)$-двойственный базис $(v_1,\ldots,v_n)$ такой, что
$\mathcal{F}=\mathcal{F}(v_1,\ldots,v_n)$. С другой стороны, то, что $\mathcal{F}\in\mathcal{O}_{(w,\varepsilon)}$ влечёт
\begin{equation}
\label{8-new}
w_k=\min\{\ell=1,\ldots,n:\delta(F_k)\cap F_\ell\not=\delta(F_{k-1})\cap F_\ell\}\ \mbox{ для всех $k\in\{1,\ldots,n\}$.}
\end{equation}

Для всех $\ell\in\{0,\ldots,n\}$ мы построим $(w,\varepsilon)$-двойственный базис
$(v_1^{(\ell)},\ldots,v_n^{(\ell)})$ такой, что
\begin{eqnarray}
\label{9-new}
 & F_k=\langle v_1^{(\ell)},\ldots,v_k^{(\ell)}\rangle_\mathbb{C}\quad\mbox{для всех $k\in\{1,\ldots,n\}$}
\\
\label{10-new}
 \mbox{и} & \delta(v_k^{(\ell)})=\left\{\begin{array}{ll}
v_{w_k}^{(\ell)}, & \mbox{если $w_k\not=k$,} \\
\varepsilon_k v_{k}^{(\ell)}, & \mbox{если $w_k=k$} \\
\end{array}\right.\quad\mbox{для всех $k\in\{1,\ldots,\ell\}$.}
\end{eqnarray}
Это даст нам
базис $(v_1^{(n)},\ldots,v_n^{(n)})$, являющийся одновременно $(w,\varepsilon)$-двойственным и $(w,\varepsilon)$-сопряжённым и таким, что
$\mathcal{F}=\mathcal{F}(v_1^{(n)},\ldots,v_n^{(n)})$, и, таким образом, завершит доказательство части {\rm (a)}.

Построение будет посредством индукции по $\ell\in\{0,\ldots,n\}$, с базой индукции
$(v_1^{(0)},\ldots,v_n^{(0)}):=(v_1,\ldots,v_n)$. Предположим, что для
$\ell\in\{1,\ldots,n\}$  базис
$(v_1^{(\ell-1)},\ldots,v_n^{(\ell-1)})$ уже построен. Мы различаем три случая.

\medskip
\noindent {\it Случай 1: $w_\ell<\ell$.}

Поскольку в этом случае $w_\ell\leq\ell-1$ и $w(w_\ell)=\ell$, мы получаем
$\delta(v_{w_\ell}^{(\ell-1)})=v_{\ell}^{(\ell-1)}$ и поэтому
$\delta(v_{\ell}^{(\ell-1)})=v_{w_\ell}^{(\ell-1)}$ (так как
$\delta$ инволюция). Поэтому базис
$(v_1^{(\ell)},\ldots,v_n^{(\ell)}):=(v_1^{(\ell-1)},\ldots,v_n^{(\ell-1)})$
удовлетворяет условиям (\ref{9-new}) и (\ref{10-new}).

\medskip
\noindent {\it Случай 2: $w_\ell=\ell$.}

Используя (\ref{8-new}), мы получаем
\[\delta(v_\ell^{(\ell-1)})\in\langle v_1^{(\ell-1)},v_2^{(\ell-1)},\ldots,v_\ell^{(\ell-1)}\rangle_\mathbb{C}+\langle v_{w_1}^{(\ell-1)},\ldots,v_{w_{\ell-1}}^{(\ell-1)}\rangle_\mathbb{C}\,.\]
С другой стороны, из того, что базис
$(v_1^{(\ell-1)},\ldots,v_n^{(\ell-1)})$
$(w,\varepsilon)$-сопряжён следует
\begin{align}v_\ell^{(\ell-1)}\in\langle v_1^{(\ell-1)},\ldots,v_{\ell-1}^{(\ell-1)},v_{w_1}^{(\ell-1)},\ldots,v_{w_{\ell-1}}^{(\ell-1)}\rangle_\mathbb{C}^\dag.
\label{E:induction}
\end{align}
Поскольку $\delta$ сохраняет ортогональность по отношению к форме
$\phi$, и  $\delta(v_k^{(\ell-1)})=v_{w_k}^{(\ell-1)}$
для всех $k\in\{1,\ldots,\ell-1\}$ (по предположению индукции), (\ref{E:induction}) влечёт
\[\delta(v_\ell^{(\ell-1)})\in\langle v_1^{(\ell-1)},\ldots,v_{\ell-1}^{(\ell-1)},v_{w_1}^{(\ell-1)},\ldots,v_{w_{\ell-1}}^{(\ell-1)}\rangle_\mathbb{C}^\dag.\]
В итоге мы получаем
\[\delta(v_\ell^{(\ell-1)})=\lambda
v_\ell^{(\ell-1)}\quad\mbox{для некоторого $\lambda\in\mathbb{C}^*$.}\]
Так как $\delta$ ---  инволюция, мы заключаем, что $\lambda\in\{+1,-1\}$.
Далее, из того, что
$\phi(v_\ell^{(\ell-1)},v_\ell^{(\ell-1)})=\varepsilon_\ell$ вытекает
\[\lambda\varepsilon_\ell=\phi(v_\ell^{(\ell-1)},\delta(v_\ell^{(\ell-1)}))={}^t\overline{v_\ell^{(\ell-1)}}\Phi\Phi v_\ell^{(\ell-1)}={}^t\overline{v_\ell^{(\ell-1)}}v_\ell^{(\ell-1)}\geq 0.\]
Наконец, мы получаем $\lambda=\varepsilon_\ell$. Отсюда следует, что базис
$(v_1^{(\ell)},\ldots,v_n^{(\ell)}):=(v_1^{(\ell-1)},\ldots,v_n^{(\ell-1)})$
удовлетворяет (\ref{9-new}) и (\ref{10-new}).

\medskip
\noindent {\it Случай 3: $w_\ell>\ell$.}

Вспоминая (\ref{8-new}), тот факт, что
$(v_1^{(\ell-1)},\ldots,v_n^{(\ell-1)})$
$(w,\varepsilon)$-двойственен, предположение индукции и  факт, что $\delta$ сохраняет ортогональность по отношению к форме $\phi$, мы получаем, как и в случае 2 ,
\begin{eqnarray*}
\delta(v_\ell^{(\ell-1)}) & \in & \big(\langle v_k^{(\ell-1)}:1\leq
k\leq w_\ell\rangle_\mathbb{C}+\langle v_{w_k}^{(\ell-1)}:1\leq k\leq
\ell-1\rangle_\mathbb{C}\big) \\
 & & \cap\,\langle
v_1^{(\ell-1)},\ldots,v_{\ell-1}^{(\ell-1)},v_{w_1}^{(\ell-1)},\ldots,v_{w_{\ell-1}}^{(\ell-1)}\rangle_\mathbb{C}^\dag.
\end{eqnarray*} Следовательно
\begin{equation} \label{11-new}
\delta(v_\ell^{(\ell-1)})=\sum_{k\in
I}\lambda_kv_k^{(\ell-1)}\quad\mbox{с $\lambda_k\in\mathbb{C}$,}
\end{equation}
где $I:=\{k:\ell\leq k\leq w_\ell,\ \ell\leq w_k\}\subset\hat{I}:=\{k:\ell\leq k,\ \ell\leq w_k\}$.
Это влечёт
\[\lambda_{w_\ell}=\phi(v_\ell^{(\ell-1)},\delta(v_\ell^{(\ell-1)}))={}^t\overline{v_\ell^{(\ell-1)}}\Phi\Phi v_\ell^{(\ell-1)}={}^t\overline{v_\ell^{(\ell-1)}}v_\ell^{(\ell-1)}\in\mathbb{R}_+^*.\]
Нетрудно проверить, что базис
$(v_1^{(\ell)},\ldots,v_n^{(\ell)})$, определённый как
\begin{eqnarray*}
 & & v_\ell^{(\ell)}:=\frac{1}{\sqrt{\lambda_{w_\ell}}}v_\ell^{(\ell-1)},\quad
v_{w_\ell}^{(\ell)}:=\frac{1}{\sqrt{\lambda_{w_\ell}}}\delta(v_\ell^{(\ell-1)}),
\\
 & & v_k^{(\ell)}:=v_k^{(\ell-1)}-\frac{\phi(v_k^{(\ell-1)},\delta(v_\ell^{(\ell-1)}))}{\lambda_{w_\ell}}v_\ell^{(\ell-1)}
\quad\mbox{для всех $k\in \hat{I}\setminus\{\ell,w_\ell\}$,} \\
 &  & v_k^{(\ell)}:=v_k^{(\ell-1)}\quad\mbox{для всех $k\in\{1,\ldots,n\}\setminus \hat{I}$,}
\end{eqnarray*}
 является $(w,\varepsilon)$-двойственным и удовлетворяет условиям
(\ref{9-new}) и (\ref{10-new}).
\end{proof}

\begin{proof}[Доказательство предложения \ref{P2}\,{\rm (b)}--{\rm (d)}]
Пусть $\mathcal{F}=\mathcal{F}(v_1,\ldots,v_n)$, где
$(v_1,\ldots,v_n)$ --- $(w,\varepsilon)$-сопряжённый базис. Тогда, по определению, мы имеем
\[\delta(\langle v_1,\ldots,v_k\rangle_\mathbb{C})=\langle v_{w_j}:j\in\{1,\ldots,k\}\rangle_\mathbb{C}\,,\]
и поэтому \begin{eqnarray*} \dim \delta(\langle
v_1,\ldots,v_k\rangle_\mathbb{C})\cap \langle v_1,\ldots,v_\ell\rangle_\mathbb{C} & = &
|\{j\in\{1,\ldots,\ell\}:w^{-1}_j\in\{1,\ldots,k\}\}| \\
  & = &
|\{j\in\{1,\ldots,\ell\}:w_j\in\{1,\ldots,k\}\}|
\end{eqnarray*}
для всех $k,\ell\in\{1,\ldots,n\}$. Далее, для
$\varepsilon\in\{+1,-1\}$ мы получаем
\begin{eqnarray*}
\langle
v_1,\ldots,v_\ell\rangle_\mathbb{C}\cap\ker(\delta-\varepsilon\mathrm{id}) & =
&
\langle v_j:1\leq w_j=j\leq \ell\mbox{ и }\varepsilon_j=\varepsilon\rangle_\mathbb{C} \\
 & & +\langle
v_j+\varepsilon v_{w_j}:1\leq w_j<j\leq \ell\rangle_\mathbb{C}\,.
\end{eqnarray*}
Следовательно
\[
\big(\dim \langle
v_1,\ldots,v_\ell\rangle_\mathbb{C}\cap V_+,\dim \langle
v_1,\ldots,v_\ell\rangle_\mathbb{C}\cap V_-\big)_{\ell=1}^n=\varsigma(w,\varepsilon).
\]
В итоге это даёт включение
\begin{equation}
\label{3} \mathcal{O}_{(w,\varepsilon)}\subset\big\{\mathcal{F}\in
X: (\delta(\mathcal{F}),\mathcal{F})\in\mathbb{O}_{w}\mbox{ и
}\varsigma(\delta:\mathcal{F})=\varsigma(w,\varepsilon)\big\}.
\end{equation}

Теперь пусть $(v_1,\ldots,v_n)$ --- $(w,\varepsilon)$-двойственный базис. Тогда
\begin{eqnarray*}
\langle v_1,\ldots,v_{n-k}\rangle_\mathbb{C}^\dag\cap\langle
v_1,\ldots,v_\ell\rangle_\mathbb{C} & = & \langle
v_j:j\in\{1,\ldots,\ell\}\mbox{ и }w_j>n-k\rangle_\mathbb{C} \\
 & = & \langle
v_j:j\in\{1,\ldots,\ell\}\mbox{ и }(w_0w)_j\leq k\rangle_\mathbb{C}\,,
\end{eqnarray*}
откуда
\[\dim\langle v_1,\ldots,v_{n-k}\rangle_\mathbb{C}^\dag\cap\langle v_1,\ldots,v_\ell\rangle_\mathbb{C}=|\{j\in\{1,\ldots,\ell\}:(w_0w)_j\in\{1,\ldots,k\}|\]
для всех $k,\ell\in\{1,\ldots,n\}$. В частности, мы видим, что
\[\langle v_1,\ldots,v_\ell\rangle_\mathbb{C}=\langle v_1,\ldots,v_\ell\rangle_\mathbb{C}\cap\langle v_1,\ldots,v_\ell\rangle_\mathbb{C}^\dag\oplus\langle v_j:
j\in\{1,\ldots,\ell\}\mbox{ и }w_j\leq \ell\rangle_\mathbb{C}.\] Это влечёт, что векторы $v_j$ (для $1\leq w_j=j\leq \ell$) и
$\frac{1}{\sqrt{2}}(v_j\pm v_{w_j})$ (для $1\leq w_j<j\leq \ell$)
составляют базис факторпространства $\langle
v_1,\ldots,v_\ell\rangle_\mathbb{C}/\langle v_1,\ldots,v_\ell\rangle_\mathbb{C}\cap\langle
v_1,\ldots,v_\ell\rangle_\mathbb{C}^\dag$. Этот базис $\phi$-ортогонален и,
поскольку $(v_1,\ldots,v_n)$  $(w,\varepsilon)$-двойственен, мы имеем
\[
\textstyle \phi(v_j,v_j)=\varepsilon_j,\mbox{ если $w_j=j$};\quad
\left\{\begin{array}{l}
\phi\big(\frac{v_j+v_{w_j}}{\sqrt{2}},\frac{v_j+v_{w_j}}{\sqrt{2}}\big)=1,\\[2mm]
\phi\big(\frac{v_j-v_{w_j}}{\sqrt{2}},\frac{v_j-v_{w_j}}{\sqrt{2}}\big)=-1
\end{array}\right.\mbox{ если $w_j<j$.}
\]
Поэтому сигнатура $\phi$ на $\langle
v_1,\ldots,v_\ell\rangle_\mathbb{C}/\langle v_1,\ldots,v_\ell\rangle_\mathbb{C}\cap\langle
v_1,\ldots,v_\ell\rangle_\mathbb{C}^\dag$ это пара
\[
\begin{array}{rll}
\big( & |\{j:w_j=j\leq\ell,\
\varepsilon_j=+1\}|+|\{j:w_j<j\leq\ell\}|,
\\[2mm]
 & |\{j:w_j=j\leq\ell,\
\varepsilon_j=-1\}|+|\{j:w_j<j\leq\ell\}| & \big),\end{array}
\]
которая совпадает с  $\ell$-ым элементом последовательности
$\varsigma(w,\varepsilon)$. Наконец, мы получаем включение
\begin{equation}
\label{4} \mathfrak{O}_{(w,\varepsilon)}\subset\big\{\mathcal{F}\in
X: (\mathcal{F}^\dag,\mathcal{F})\in\mathbb{O}_{w_0w}\mbox{ и
}\varsigma(\phi:\mathcal{F})=\varsigma(w,\varepsilon)\big\}.
\end{equation}

Ясно, что $K$ (соответственно, $G^0$) действует транзитивно на множестве
$(w,\varepsilon)$-сопряжённых базисов (соответственно, $(w,\varepsilon)$-двойственных
базисов). Поэтому подмножества $\mathcal{O}_{(w,\varepsilon)}$ (соответственно,
$\mathfrak{O}_{(w,\varepsilon)}$) суть  $K$-орбиты (соответственно,
$G^0$-орбиты). Далее, принимая во внимание (\ref{3}) и (\ref{4}) мы видим, что эти орбиты попарно различны.

Пусть $\mathcal{O}$ это  $K$-орбита в $X$. Отметим, что базис $(e_1,\ldots,e_n)$ в
$V$ удовлетворяет $\delta(e_j)=\pm
e_j$ для всех $j$, следовательно, флаг
$\mathcal{F}_0:=\mathcal{F}(e_1,\ldots,e_n)$ удовлетворяет
$\delta(\mathcal{F}_0)=\mathcal{F}_0$. Из
\cite{Richardson-Springer} следует, что  $K$-орбита $\mathcal{O}$ содержит элементы вида $g\mathcal{F}_0$ для некоторого $g\in G$ такого, что
$h:=\Phi g^{-1}\Phi g\in N$ где, как и в доказательстве предложения
\ref{P1}, $N\subset G$ обозначает подгруппу матриц с ровно
одним ненулевым коэффициентом в каждой строке и каждом столбце. Так как
$\Phi\in N$, то тоже $\Phi h\in N$. Следовательно существует перестановка
$w\in\mathfrak{S}_n$ и константы
$t_1,\ldots,t_n\in\mathbb{C}^*$ такие, что матрица $\Phi
h=:\big(a_{k,\ell}\big)_{1\leq k,\ell\leq n}$ составлена из
\[a_{k,\ell}=0,\ \mbox{ если $\ell\not=w_k$,}\quad a_{k,w_k}=t_k\quad\mbox{ для всех $k,\ell\in\{1,\ldots,n\}$}.\]
Соотношение $\Phi h=g^{-1}\Phi g$ показывает, что $(\Phi h)^2=1_n$. Это
даёт $w^2=\mathrm{id}$ и $t_kt_{w_k}=1$ для всех $k$; поэтому
\[t_{w_k}=t_k^{-1}\ \mbox{ при $w_k\not=k$}\quad\mbox{и}\quad \varepsilon_k:=t_k\in\{+1,-1\}\ \mbox{ при
$w_k=k$}.\] В дополнение, поскольку матрица $\Phi h$ сопряжена матрице $\Phi$, её собственные значения $+1$ и $-1$ имеют соответственно кратности $p$ и
$q$, откуда
\[(w,\varepsilon)\in\mathfrak{I}_n(p,q).\]
Для каждого $k\in\{1,\ldots,n\}$ с $w_k<k$, мы
выберем  $s_k\in\mathbb{C}^*$ так, чтобы $t_k=s_k^2$, и положим
$s_{w_k}=s_k^{-1}$ (так, чтобы $s_{w_k}^2=t_k^{-1}=t_{w_k}$). Далее,
для каждого $k\in\{1,\ldots,n\}$ с $w_k=k$, положим $s_k=1$. Равенство $\Phi g=g\Phi h$ даёт
\[
\delta(g(s_ke_k))=s_k\Phi g e_k=s_k g(\Phi h)e_k=s_k
g(t_{w_k}e_{w_k})=s_{w_k}^{-1}
g(s_{w_k}^2e_{w_k})=g(s_{w_k}e_{w_k})
\]
для всех $k\in\{1,\ldots,n\}$ таких, что $w_k\not=k$, и
\[
\delta(g(s_k e_k))=\delta(g(e_k))=\Phi g e_k=g(\Phi
h)e_k=g(\varepsilon_k e_k)=\varepsilon_k g(e_k)=\varepsilon_k g(s_k
e_k)
\]
для всех $k\in\{1,\ldots,n\}$ таких, что $w_k=k$. Следовательно, множество
$(g(s_1 e_1),\ldots,\\g(s_n e_n))$ является $(w,\varepsilon)$-сопряжённым базисом в
$V$. Таким образм
\[g\mathcal{F}_0=g\mathcal{F}(e_1,\ldots,e_n)=g\mathcal{F}(s_1 e_1,\ldots,s_n e_n)=\mathcal{F}(g(s_1 e_1),\ldots,g(s_n e_n))\in\mathcal{O}_{(w,\varepsilon)}.\]
Следовательно $\mathcal{O}=\mathcal{O}_{(w,\varepsilon)}$.

Мы заключаем, что подмножества $\mathcal{O}_{(w,\varepsilon)}$ (для
$(w,\varepsilon)\in\mathfrak{I}_n(p,q)$) суть в точности $K$-орбиты
в $X$. Тогда из двойственности Мацуки вытекает, что подмножества
$\mathfrak{O}_{(w,\varepsilon)}$ (для
$(w,\varepsilon)\in\mathfrak{I}_n(p,q)$) являются в точности
$G^0$-орбитами в $X$. Этот факт влечёт, в частности, что (\ref{3}) и (\ref{4}) на самом деле являются равенствами. В итоге мы доказали части {\rm (b)} и {\rm (c)} утверждения.

Наконец, часть {\rm (a)} показывет, что для каждой
$(w,\varepsilon)\in\mathfrak{I}_n(p,q)$ пересечение
$\mathcal{O}_{(w,\varepsilon)}\cap\mathfrak{O}_{(w,\varepsilon)}$
состоит из единой $K\cap G^0$-орбиты, и откуда следует, что орбиты  $\mathcal{O}_{(w,\varepsilon)}$ и
$\mathfrak{O}_{(w,\varepsilon)}$ двойственны по Мацуки (см.
\cite{Matsuki,Matsuki3}). Это доказывает часть {\rm (d)} утверждения. Доказательство предложения \ref{P2} закончено.
\end{proof}

\begin{proof}[Доказательство предложения \ref{P3}]
Доказательство опирается на следующие два технических утверждения.

\medskip
\item
{\it Утверждение 1:}
Для каждой инволюции со знаками $(w,\varepsilon)\in\mathfrak{I}_n(p,q)$ мы имеем
$\mathcal{O}_{(w,\varepsilon)}\cap X_\omega=\emptyset$, если $(w,\varepsilon)\notin\mathfrak{I}_n^{\eta,\epsilon}(p,q)$.

\medskip
\noindent
{\it Утверждение 2:}
Для каждой инволюции со знаками $(w,\varepsilon)\in\mathfrak{I}_n^{\eta,\epsilon}(p,q)$ есть базис $\underline{v}=(v_1,\ldots,v_n)$, являющийся одновременно $(w,\varepsilon)$-двойственным и $(w,\varepsilon)$-сопряжённым, и удовлетворяющим (\ref{3.3.18}).

\medskip

Опираясь на утверждения 1 и 2, продолжим доказательство. Для каждой $(w,\varepsilon)\in\mathfrak{I}_n(p,q)$ включения
\begin{eqnarray}
\label{proof.P3.1} & \{\mathcal{F}(\underline{v}):\mbox{$\underline{v}$  $(w,\varepsilon)$-сопряжён
 и удовлетворяет (\ref{3.3.18})}\}\subset\mathcal{O}_{(w,\varepsilon)}\cap X_\omega,
 \\
\label{proof.P3.2} & \{\mathcal{F}(\underline{v}):\mbox{$\underline{v}$  $(w,\varepsilon)$-двойственен
 и удовлетворяет (\ref{3.3.18})}\}\subset\mathfrak{O}_{(w,\varepsilon)}\cap X_\omega, \\
\label{proof.P3.3}  & \{\mathcal{F}(\underline{v}):\mbox{$\underline{v}$  $(w,\varepsilon)$-двойственен и $(w,\varepsilon)$-сопряжён и удовлетворяет} \\ \nonumber
& \mbox{ условию  (\ref{3.3.18}) }\}\\ \nonumber
& \subset
\mathcal{O}_{(w,\varepsilon)}\cap\mathfrak{O}_{(w,\varepsilon)}\cap X_\omega
\end{eqnarray}
очевидно выполняются. Поэтому утверждение 2 показывает, что
$\mathcal{O}_{(w,\varepsilon)}^{\eta,\epsilon}$, $\mathfrak{O}_{(w,\varepsilon)}^{\eta,\epsilon}$, и
$\mathcal{O}_{(w,\varepsilon)}^{\eta,\epsilon}\cap\mathfrak{O}_{(w,\varepsilon)}^{\eta,\epsilon}$ все непусты, когда $(w,\varepsilon)\in\mathfrak{I}_n^{\eta,\epsilon}(p,q)$.
По утверждению 1, лемме \ref{lemma-2.2.1}, и предложению \ref{P2}\,{\rm (c)},  $K$-орбиты в $X_\omega$ суть в точности подмножества $\mathcal{O}_{(w,\varepsilon)}^{\eta,\epsilon}$.
С другой стороны, подмножества $\mathfrak{O}_{(w,\varepsilon)}\cap X_\omega$ (для $(w,\varepsilon)\in\mathfrak{I}_n(p,q)$)  $G^0$-инвариантны и попарно различны. По двойственности Мацуки существует биекция между $K$-орбитами и $G^0$-орбитами. Это заставляет $\mathfrak{O}_{(w,\varepsilon)}^{\eta,\epsilon}=\mathfrak{O}_{(w,\varepsilon)}\cap X_\omega$ быть единой $G^0$-орбитой при $(w,\varepsilon)\in\mathfrak{I}_n^{\eta,\epsilon}(p,q)$, и  $\mathfrak{O}_{(w,\varepsilon)}\cap X_\omega$ быть пустым при $(w,\varepsilon)\notin\mathfrak{I}_n^{\eta,\epsilon}(p,q)$.
Это доказывает предложение \ref{P3}\,{\rm (b)}.

Поскольку орбиты $\mathcal{O}_{(w,\varepsilon)},\mathfrak{O}_{(w,\varepsilon)}\subset X$ двойственны по Мацуки (см. предложение \ref{P2}\,{\rm (d)}), их пересечение $\mathcal{O}_{(w,\varepsilon)}\cap\mathfrak{O}_{(w,\varepsilon)}$ компактно, поэтому то же верно для всех пересечений $\mathcal{O}_{(w,\varepsilon)}^{\eta,\epsilon}\cap
\mathfrak{O}_{(w,\varepsilon)}^{\eta,\epsilon}$ при $(w,\varepsilon)\in\mathfrak{I}_n^{\eta,\epsilon}(p,q)$.  Это влечёт то, что $\mathcal{O}_{(w,\varepsilon)}^{\eta,\epsilon}$ и $\mathfrak{O}_{(w,\varepsilon)}^{\eta,\epsilon}$ двойственны по Мацуки (см. \cite{Gindikin-Matsuki}), и, соответственно, выполнение части {\rm (c)} утверждения .

Пусть $(w,\varepsilon)\in\mathfrak{I}_n^{\eta,\epsilon}(p,q)$. Поскольку орбиты $\mathcal{O}_{(w,\varepsilon)}^{\eta,\epsilon}$ и $\mathfrak{O}_{(w,\varepsilon)}^{\eta,\epsilon}$ двойственны по Мацуки, их пересечение --- единая $K\cap G^0$-орбита. Множество в левой части  (\ref{proof.P3.3}) непусто (по утверждению 2) и $K\cap G^0$-инвариантно, поэтому равенство из (\ref{proof.P3.3}) выполнено. По аналогии, множества в левых частях  (\ref{proof.P3.1}) и (\ref{proof.P3.2}) непусты (по утверждению 2) и, соответственно, $K$- и $G^0$-инвариантны. Поскольку $\mathcal{O}_{(w,\varepsilon)}^{\eta,\epsilon}=\mathcal{O}_{(w,\varepsilon)}\cap X_\omega$ и $\mathfrak{O}_{(w,\varepsilon)}^{\eta,\epsilon}=\mathfrak{O}_{(w,\varepsilon)}\cap X_\omega$  являются, соответственно, $K$-орбитой и $G^0$-орбитой, мы получаем, что (\ref{proof.P3.1}) и (\ref{proof.P3.2}) суть равенства. Это показывает часть {\rm (a)} утверждения.

Таким образом, доказательство предложения \ref{P3} будет завершено, когда мы установим утверждения 1 и 2.

\medskip
\noindent
{\it Доказательства утверждения 1.}
Заметим, что для двух подпространств $A,B\subset V$ мы имеем
$A^\perp+B^\perp=(A\cap B)^\perp$, поэтому
\begin{equation}
\label{proof-P3.4}
\dim A^\perp\cap B^\perp+\dim A+\dim B=\dim A\cap B+\dim V.
\end{equation}
Заметим также, что отображение $\delta$
самосопряжено (для типов BD1 и C2) или антисамосопряжено (для типов C1 и D3) по отношению к форме $\omega$, поэтому равенство $\delta(A)^\perp=\delta(A^\perp)$ выполняется для любого подпространства $A\subset V$ и для всех типов.

Пусть $(w,\varepsilon)\in\mathfrak{I}_n(p,q)$ такая, что $\mathcal{O}_{(w,\varepsilon)}\cap X_\omega\not=\emptyset$.
Положим $\mathcal{F}=(F_0,\ldots,F_n)\in \mathcal{O}_{(w,\varepsilon)}\cap X_\omega$.

Применяя (\ref{proof-P3.4}) к $A=\delta\left( F_k\right)$ и $B=F_\ell$ для $1\leq k,\ell\leq n$, получаем
\begin{equation}
\label{proof-P3.5}
\dim \delta\left( F_{n-k}\right)\cap F_{n-\ell}+k+\ell=\dim \delta \left(F_k\right)\cap F_\ell+n.
\end{equation}
С другой стороны, поскольку $\mathcal{F}\in\mathcal{O}_{(w,\varepsilon)}$, предложение \ref{P2}\,{\rm (b)} даёт
\begin{equation}
\label{proof-P3.6}
\dim \delta\left( F_{n-k}\right)\cap F_{n-\ell}=|\{j=1,\ldots,n-\ell:1\leq w_j\leq n-k\}|
\end{equation}
и
\begin{eqnarray}
\label{proof-P3.7}
\dim\delta \left(F_k\right)\cap F_\ell & = & |\{j=1,\ldots,\ell:1\leq w_j\leq k\}| \\
\nonumber & = & \ell-|\{j=1,\ldots,\ell:w_j\geq k+1\}| \\
\nonumber & = & \ell-(n-k-|\{j\geq \ell+1:w_j\geq k+1\}|) \\
\nonumber & = & \ell+k-n+|\{j=1,\ldots,n-\ell:w_0ww_0(j)\leq n-k\}|
\end{eqnarray}
для всех $k,\ell\in\{1,\ldots,n\}$. Сравнивая (\ref{proof-P3.5})--(\ref{proof-P3.7}), мы заключаем, что $w=w_0ww_0$.

Пусть $k\in\{1,\ldots,n\}$ такое, что $w_k=k$. Поскольку $ww_0=w_0w$, мы имеем $w_{n-k+1}=n-k+1$.
Применяя (\ref{proof-P3.4}) для $A=F_k$ (соответственно, $A=F_{k-1}$) и $B=V_+$, мы получаем
\[1+\dim F_{k-1}\cap V_+-\dim F_k\cap V_+=\dim F_{n-k+1}\cap V_--\dim F_{n-k}\cap V_-\]
для типов BD1 и C2 (где $V_+^\perp=V_-$), откуда
\begin{eqnarray*}
\varepsilon_k=1 & \Leftrightarrow &  \dim F_k\cap V_+=\dim F_{k-1}\cap V_++1 \\
 & \Leftrightarrow &  \dim F_{n-k+1}\cap V_-=\dim F_{n-k}\cap V_-\Leftrightarrow\varepsilon_{n-k+1}=1
\end{eqnarray*}
в этом случае.  Для типов C1 и D3  (где $V_+^\perp=V_+$), мы получаем
\[1+\dim F_{k-1}\cap V_+-\dim F_k\cap V_+=\dim F_{n-k+1}\cap V_+-\dim F_{n-k}\cap V_+\,,\]
откуда также
\[\varepsilon_k=1\Leftrightarrow \varepsilon_{n-k+1}=-1\,.\]

На этом этапе мы получаем, что инволюция со знаками $(w,\varepsilon)$ удовлетворяет условиям {\rm (i)}--{\rm (ii)} из раздела \ref{S3.3}.
Чтобы заключить, что $(w,\varepsilon)\in\mathfrak{I}_n^{\eta,\epsilon}(p,q)$, остаётся проверить, что для типов C2 и D3 мы имеем $w_k\not=n-k+1$ при всех $k\leq\frac{n}{2}$. Рассуждая от противного, предположим, что $w_k=n-k+1$. Поскольку $\mathcal{F}\in\mathcal{O}_{(w,\varepsilon)}$, то существует $(w,\varepsilon)$-сопряжённый базис $\underline{v}=(v_1,\ldots,v_n)$ такой, что $\mathcal{F}=\mathcal{F}(\underline{v})$.
Таким образом, $\delta(v_k)=v_{n-k+1}$ и мы можем написать $v_k=v_k^++v_k^-$ и $v_{n-k+1}=v_k^+-v_k^-$.
Для типа C2 мы имеем $V_+^\perp=V_-$ и $\omega$ антисимметрична, поэтому
\[\omega(v_k^++v_k^-,v_k^+-v_k^-)=\omega(v_k^+,v_k^+)-\omega(v_k^-,v_k^-)=0-0=0.\]
Для типа D3 мы имеем $V_+^\perp=V_+$, $V_-^\perp=V_-$, и $\omega$ симметрична, и поэтому
\[\omega(v_k^++v_k^-,v_k^+-v_k^-)=-\omega(v_k^+,v_k^-)+\omega(v_k^-,v_k^+)=0.\]
В обоих случаях мы выводим
\[F_{n-k+1}=F_{n-k}+\langle v_{n-k+1}\rangle_{\mathbb{C}}\subset F_k^\perp+F_{k-1}^\perp\cap\langle v_k\rangle_\mathbb{C}^\perp=F_k^\perp=F_{n-k},\]
что есть противоречие.
Этим завершается доказательство утверждения 1.

\medskip
\noindent
{\it Доказательство утверждения 2.}
Для $k\in\{1,\ldots,n\}$ положим $k^*=n-k+1$. Запишем $\omega$ как
\[w=(c_1;c'_1)\cdots(c_s;c'_s)(c'^*_1;c^*_1)\cdots(c'^*_s;c^*_s)(d_1;d^*_1)\cdots(d_t;d^*_t),\]
где $c_1<\ldots<c_s<c^*_s<\ldots<c^*_1$, $c_j<c'_j\not=c^*_j$ для всех $j$, $d_1<\ldots<d_t<d^*_t<\ldots<d^*_1$. Отметим, что $t=0$ для типов C2 и D3.
Далее, мы обозначим
\begin{eqnarray*}
 & \{a_1<\ldots<a_{p-t-2s}\}:=\{k:w_k=k,\ \varepsilon_k=1\},\\
 & \{b_1<\ldots<b_{q-t-2s}\}:=\{k:w_k=k,\ \varepsilon_k=-1\}.
\end{eqnarray*}

Зафиксируем $\phi$-ортонормальный базис $V_+$
\[x_1^+,\ldots,x_t^+,y^+_1,\ldots,y^+_s,y^{+*}_s,\ldots,y^{+*}_1,z^+_1,\ldots,z^+_{p-t-2s}\]
и $(-\phi)$-ортонормальный базис $V_-$
\[x_1^-,\ldots,x_t^-,y^-_1,\ldots,y^-_s,y^{-*}_s,\ldots,y^{-*}_1,z^-_1,\ldots,z^-_{q-t-2s},\]
такие, что в случаях BD1 и C2 (где ограничение $\omega$ на $V_+$ и $V_-$ невырождено)  имеем
\begin{eqnarray*}
 & \omega(x^+_j,x^+_j)=\omega(x^-_j,x^-_j)=1,\\
 & \omega(y^+_j,y^{+*}_j)=\omega(y^-_j,y^{-*}_j)=1,\quad\omega(y^{+*}_j,y^+_j)=\omega(y^{-*}_j,y^-_j)=\epsilon, \\
 & \omega(z^+_j,z^+_{\ell})=\left\{\begin{array}{ll} 1, & \mbox{если $j\leq \ell=p-t-2s+1-j$} \\ \epsilon, & \mbox{если $j>\ell=p-t-2s+1-j, $} \end{array}\right.
\\
 & \omega(z^-_j,z^-_{\ell})=\left\{\begin{array}{ll} 1, & \mbox{если $j\leq \ell=q-t-2s+1-j$} \\ \epsilon, & \mbox{если $j>\ell=q-t-2s+1-j, $} \end{array}\right.
\end{eqnarray*}
и остальные значения $\omega$ на базисе  равны $0$.  Для типов C1 и D3 (где $V_+^\perp=V_+$, $V_-^\perp=V_-$, и, в частности, $p=q=\frac{n}{2}$ в этом случае) мы требуем, чтобы
\begin{eqnarray*}
 & \omega(x^+_j,x^-_j)=i,\quad \omega(x^-_j,x^+_j)=\epsilon i,\\
 & \omega(y^+_j,y^{-*}_j)=\omega(y^-_j,y^{+*}_j)=1,\quad\omega(y^{+*}_j,y^-_j)=\omega(y^{-*}_j,y^+_j)=\epsilon, \\
 & \omega(z^+_j,z^-_{\ell})=\epsilon\omega(z^-_\ell,z^+_j)=\left\{\begin{array}{ll} 1, & \mbox{если $\ell=\tilde{j}:=\frac{n}{2}-t-2s+1-j$ и $a_j<b_{\tilde{j}}$} \\ \epsilon, & \mbox{если $\ell=\tilde{j}:=\frac{n}{2}-t-2s+1-j$ и $a_j>b_{\tilde{j}}$,} \end{array}\right.
\end{eqnarray*}
и чтобы  остальные значения $\omega$ на базисе были равны $0$.
В отличие от значений $\omega(z_j^\pm,z_\ell^\pm)$ в случаях BD1,C2, значения $\omega(z_j^+,z_\ell^-)$ в случаях C1,D3 не подлежат ограничанию но выбраны так, чтобы базис $(v_1,\ldots,v_n)$ ниже удовлетворял (\ref{3.3.18}).

Во всех случаях мы строим базис $(v_1,\ldots,v_n)$ следующим образом
\[v_{d_j}=\frac{x^+_j+ix^-_j}{\sqrt{2}},\quad v_{d_j^*}=\frac{x^+_j-ix^-_j}{\sqrt{2}},\]
\[v_{c_j}=\frac{y^+_j+y^-_j}{\sqrt{2}},\quad v_{c'_j}=\frac{y^+_j-y^-_j}{\sqrt{2}},\quad v_{c^*_j}=\frac{y^{+*}_j+y^{-*}_j}{\sqrt{2}},\quad v_{c'^*_j}=\frac{y^{+*}_j-y^{-*}_j}{\sqrt{2}},\]
\[v_{a_j}=z^+_j,\quad\mbox{и}\quad v_{b_j}=z_j^-.\]
Просто проверить, что базис $(v_1,\ldots,v_n)$ является одновременно $(w,\varepsilon)$-двойственным и $(w,\varepsilon)$-сопряжённым и, что он удовлетворяет (\ref{3.3.18}). Это завершает доказательство утверждения 2.
\end{proof}



\section{Двойственность орбит в инд-многообразиях обобщённых флагов}

\label{S4}

Следуя схеме раздела ~\ref{S3}, мы теперь изложим наши результаты о двойственности орбит в бесконечномерном случае.  Все доказательства представлены в разделе~\ref{S4.proofs}.
\subsection{Типы A1 и A2}

\label{S4.1}

Обозначения такие же, как  в разделе \ref{S:A1A2}.
Для каждого $\ell\in\NN$ существует единственное $\ell^*\in\NN$
такое, что $\omega(e_\ell,e_{\ell^*})\not=0$, и
это даёт биекцию $\iota:\NN\to\NN$, $\ell\mapsto\ell^*$.

Обозначим через $\mathfrak{I}_\infty(\iota)$ - множество инволлюций $w:\NN\to\NN$
таких, что $w(\ell)=\ell^*$ для всех кроме конечного числа $\ell\in\NN$.
В частности, мы имеем $w\iota\in\mathfrak{S}_\infty$ для всех $w\in\mathfrak{I}_\infty(\iota)$.
Пусть $\mathfrak{I}'_\infty(\iota)\subset \mathfrak{I}_\infty(\iota)$ ---
подмножество инволюций без неподвижных точек (то есть, таких, что $w(\ell)\not=\ell$ для всех $\ell\in\NN$).

Пусть $\sigma:\NN\to(A,\prec)$ --- биекция на полностью упорядоченное множество,
и  рассмотрим инд-многообразие обобщённых флагов $\mathbf{X}(\mathcal{F}_\sigma,E)$.
В предложении \ref{P4-1.1} ниже мы покажем, что $\mathbf{K}$-орбиты
и $\mathbf{G}^0$-орбиты в $\mathbf{X}(\mathcal{F}_\sigma,E)$ параметризуются элементами $\mathfrak{I}_\infty(\iota)$ для типа A1, и элементами $\mathfrak{I}'_\infty(\iota)$ для типа A2.

\begin{definition}
Пусть $w\in\mathfrak{I}_\infty(\iota)$.
И $\underline{v}=(v_1,v_2,\ldots)$ --- базис $\mathbf{V}$ такой, что
\begin{equation}
\label{4.1.1}
v_\ell=e_\ell\quad\mbox{для всех кроме конечного числа $\ell\in\NN$}.
\end{equation}
Назовём $\underline{v}$ {\em $w$-двойственным}, если в дополнение к (\ref{4.1.1}) $\underline{v}$ удовлетворяет
\[\omega(v_\ell,v_k)=\left\{
\begin{array}{ll}
0, & \mbox{если $\ell\not=w_k$,} \\
\pm1, & \mbox{если $\ell=w_k$}
\end{array}
\quad\mbox{для всех $k,\ell\in\NN$,}
\right.\]
и назовём $\underline{v}$ {\em $w$-сопряжённым}, если в дополнение к (\ref{4.1.1})
\[\gamma(v_k)=\pm v_{w_k}\quad\mbox{для всех $k\in\NN$.}\]
Положим
$\bs{\mathcal{O}}_w:=\{\mathcal{F}_{\sigma}(\underline{v}):\mbox{$\underline{v}$ --- $w$-двойственный}\}$ и $\bs{\mathfrak{O}}_w:=\{\mathcal{F}_{\sigma}(\underline{v}):\mbox{$\underline{v}$ ---\\ $w$-сопряжён-}\\\mbox{ный}\}$, где $\bs{\mathcal{O}}_w$ и $\bs{\mathfrak{O}}_w$ --- подмножества инд-многообразия $\mathbf{X}(\mathcal{F}_\sigma,E)$.
\end{definition}

\paragraph{\bf Обозначения.}
{\rm (a)}
Мы будем использовать сокращенное обозначение $\mathbf{X}:=\mathbf{X}(\mathcal{F}_\sigma,E)$. \\
{\rm (b)}
Если $\mathcal{F}$ --- обобщённый флаг слабо согласованный с $E$,
то $\mathcal{F}^\perp:=\{F^\perp:F\in\mathcal{F}\}$ также является обобщённым флагом слабо согласованным с  $E$.

Обозначим через $(A^*,\prec^*)$  полностью упорядоченное множество такое, что $A^*=A$ как множество и  $a\prec^*a'$, тогда и только когда $a\succ a'$.
Определим $\sigma^\perp:\NN\to(A^*,\prec^*)$, полагая $\sigma^\perp(\ell)=\sigma(\ell^*)$.
Тогда
$\mathcal{F}_\sigma^\perp=\mathcal{F}_{\sigma^\perp}$.
Отметим, что $\mathcal{F}^\perp$ --- $E$-соизмерим с $\mathcal{F}_{\sigma^\perp}$,
когда $\mathcal{F}$ --- $E$-соизмерим с $\mathcal{F}_\sigma$.
Следовательно, отображение
\[\mathbf{X}\to\mathbf{X}^\perp:=\mathbf{X}(\mathcal{F}_{\sigma^\perp},E),\ \mathcal{F}\mapsto\mathcal{F}^\perp\]
корректно определено.
Мы  будем использовать сокращённое обозначение $\pmb{\mathbb{O}}^\perp_w:=(\pmb{\mathbb{O}}_{\sigma^\perp,\sigma})_w$ для всех $w\in\mathfrak{S}_\infty$. \\
{\rm (c)}
Заметим далее, что $\gamma(\mathcal{F}_\sigma)=\mathcal{F}_{\sigma\circ\iota}$, и что $\gamma(\mathcal{F})\in\mathbf{X}^\gamma:=\mathbf{X}(\mathcal{F}_{\sigma\circ\iota},E)$ , когда $\mathcal{F}\in\mathbf{X}$.
Мы  будем использовать сокращённое обозначение $\pmb{\mathbb{O}}^\gamma_w:=(\pmb{\mathbb{O}}_{\sigma\circ\iota,\sigma})_w$ для всех $w\in\mathfrak{S}_\infty$.

\medskip
Таким образом,
$\displaystyle \mathbf{X}^\perp\times\mathbf{X}=\bigsqcup_{w\in\mathfrak{S}_\infty}\pmb{\mathbb{O}}^\perp_w$
и
$\displaystyle \mathbf{X}^\gamma\times\mathbf{X}=\bigsqcup_{w\in\mathfrak{S}_\infty}\pmb{\mathbb{O}}^\gamma_w$
(см. предложение \ref{P2-3.2}).


\begin{proposition}
\label{P4-1.1}
Пусть $\mathfrak{I}^\epsilon_\infty(\iota)=\mathfrak{I}_\infty(\iota)$ для типа A1
и $\mathfrak{I}^\epsilon_\infty(\iota)=\mathfrak{I}'_\infty(\iota)$ для типа A2.
\begin{itemize}
\item[\rm (a)] Для всех $w\in\mathfrak{I}^\epsilon_\infty(\iota)$, имеем \[\bs{\mathcal{O}}_w\cap\bs{\mathfrak{O}}_w=\{\mathcal{F}_\sigma(\underline{v}):\mbox{$\underline{v}$  $w$-двойственен и $w$-сопряжён}\}\not=\emptyset.\]
\item[\rm (b)] Для каждого $w\in\mathfrak{I}^\epsilon_\infty(\iota)$, имеем
\[
\bs{\mathcal{O}}_w=\{\mathcal{F}\in\mathbf{X}:(\mathcal{F}^\perp,\mathcal{F})\in\pmb{\mathbb{O}}^\perp_{w\iota}\} \quad\mbox{и}\quad
\bs{\mathfrak{O}}_w=\{\mathcal{F}\in\mathbf{X}:(\gamma(\mathcal{F}),\mathcal{F})\in\pmb{\mathbb{O}}^\gamma_{w\iota}\}.
\]
\item[\rm (c)] Подмножества $\bs{\mathcal{O}}_w$ (для $w\in\mathfrak{I}^\epsilon_\infty(\iota)$) суть в точности $\mathbf{K}$-орбиты в $\mathbf{X}$.
Подмножества $\bs{\mathfrak{O}}_w$ (для $w\in\mathfrak{I}^\epsilon_\infty(\iota)$) суть в точности $\mathbf{G}^0$-орбиты в $\mathbf{X}$.
Далее, $\bs{\mathcal{O}}_w\cap\bs{\mathfrak{O}}_w$ есть единая $\mathbf{K}\cap\mathbf{G}^0$-орбита.
\end{itemize}
\end{proposition}

\subsection{Тип A3}

\label{S4.2}

Обозначения такие же, как и в разделе \ref{S:A3}.
В частности, мы фиксируем разложение $\NN=N_+\sqcup N_-$, задающее $\Phi$ как и в  (\ref{phi}), и мы рассматриваем соответствующую
Эрмитову форму $\phi$ и инволюцию $\delta$ на $\mathbf{V}$.

Обозначим через $\mathfrak{I}_\infty(N_+,N_-)$  множество пар $(w,\varepsilon)$, состоящих из инволюции $w:\NN\to\NN$ и отображения $\varepsilon:\{\ell:w_\ell=\ell\}\to\{1,-1\}$ таких, что подмножества
\[N'_\pm=N'_\pm(w,\varepsilon):=\{\ell\in N_\pm:(w_\ell,\varepsilon_\ell)=(\ell,\pm1)\}
\]
удовлетворяют
\[|N_\pm\setminus N'_\pm|=|\{\ell\in N_\mp:(w_\ell,\varepsilon_\ell)=(\ell,\pm 1)\}|+\frac{1}{2}|\{\ell\in\NN:w_\ell\not=\ell\}|<\infty.\]
В частности, $w\in\mathfrak{S}_\infty$.

Зафиксируем биекцию $\sigma:\NN\to (A,\prec)$  на полностью упорядоченное множество.
Мы покажем в предложении \ref{P4.2-2}, что  $\mathbf{K}$-орбиты и $\mathbf{G}^0$-орбиты инд-многообразия $\mathbf{X}:=\mathbf{X}(\mathcal{F}_\sigma,E)$ параметризуются элементами из $\mathfrak{I}_\infty(N_+,N_-)$.

\begin{definition}
Пусть $(w,\varepsilon)\in\mathfrak{I}_\infty(N_+,N_-)$.
Базис $\underline{v}=(v_1,v_2,\ldots)$ в $\mathbf{V}$ такой, что $v_\ell=e_\ell$ для всех кроме конечного множества $\ell\in\NN$, назовём {\em $(w,\varepsilon)$-сопряжённым}, если
\[\delta(v_k)=\left\{
\begin{array}{ll}
v_{w_k}, & \mbox{если $w_k\not=k$,} \\
\varepsilon_kv_k, & \mbox{если $w_k=k$}
\end{array}
\right.
\quad\mbox{для всех $k\in\NN$,}\]
и  {\em $(w,\varepsilon)$-двойственным}, если
\[
\phi(v_k,v_\ell)=\left\{
\begin{array}{ll}
0, & \mbox{если $\ell\not=w_k$,} \\
1, & \mbox{если $\ell=w_k\not=k$,} \\
\varepsilon_k, & \mbox{если $\ell=w_k=k$}
\end{array}
\right.
\quad\mbox{для всех $k,\ell\in\NN$}.
\]
Положим
$\bs{\mathcal{O}}_{(w,\varepsilon)}:=\{\mathcal{F}_\sigma(\underline{v}):\mbox{$\underline{v}$ --- $(w,\varepsilon)$-сопряжённый}\}$,
 $\bs{\mathfrak{O}}_{(w,\varepsilon)}:=\\\{\mathcal{F}_\sigma(\underline{v}):\mbox{$\underline{v}$ --- $(w,\varepsilon)$-двойственный}\}$.

\end{definition}

\paragraph{\bf Примечание}
{\rm (a)} Заметим, что каждое подпространство в обобщённом флаге $\mathcal{F}_\sigma$ --- $\delta$-инвариантно, то есть, $\delta(\mathcal{F}_\sigma)=\mathcal{F}_\sigma$.
Отображение $\mathbf{X}\to\mathbf{X}$, $\mathcal{F}\mapsto\delta(\mathcal{F})$ корректно определено.
%We abbreviate $\pmb{\mathbb{O}}_w:=(\pmb{\mathbb{O}}_{\sigma,\sigma})_w$
%for all $w\in\mathfrak{S}_\infty$.
\\
{\rm (b)} Обозначим $F^\dag=\{x\in\mathbf{V}:\phi(x,y)=0\ \forall y\in F\}$ и
$\mathcal{F}^\dag:=\{F^\dag:F\in \mathcal{F}\}$; при этом, $\mathcal{F}^\dag$ является обобщённым флагом слабо согласованным с  $E$, когда это же верно для $\mathcal{F}$.

Как и в разделе \ref{S4.1}, мы обозначим через $(A^*,\prec^*)$ полностью упорядочнное множество
такое, что $A^*=A$ и $a\prec^*a'$, когда $a\succ a'$.
Легко видеть, что $\mathcal{F}_\sigma^\dag=\mathcal{F}_{\sigma^\dag}$,
где $\sigma^\dag:\NN\to(A^*,\prec^*)$ такое, что $\sigma^\dag(\ell)=\sigma(\ell)$ для всех $\ell\in\NN$, и поэтому мы получаем корректно определённое отображение
\[\mathbf{X}\to\mathbf{X}^\dag:=\mathbf{X}(\mathcal{F}_{\sigma^\dag},E),\ \mathcal{F}\mapsto\mathcal{F}^\dag.\]
{\rm (c)}
Обозначим $\pmb{\mathbb{O}}_w:=(\pmb{\mathbb{O}}_{\sigma,\sigma})_w$ и
$\pmb{\mathbb{O}}_w^\dag:=(\pmb{\mathbb{O}}_{\sigma^\dag,\sigma})_w$; тогда
\[\mathbf{X}\times\mathbf{X}=\bigsqcup_{w\in\mathfrak{S}_\infty}\pmb{\mathbb{O}}_w
\quad\mbox{и}\quad
\mathbf{X}^\dag\times\mathbf{X}=\bigsqcup_{w\in\mathfrak{S}_\infty}\pmb{\mathbb{O}}^\dag_w
\]
(см. предложение \ref{P2-3.2}).

\begin{proposition}
\label{P4.2-2}
\begin{itemize}
\item[\rm (a)] Для каждой $(w,\varepsilon)\in\mathfrak{I}_\infty(N_+,N_-)$ мы имеем
\[\bs{\mathcal{O}}_{(w,\varepsilon)}\cap\bs{\mathfrak{O}}_{(w,\varepsilon)}=\{\mathcal{F}_\sigma(\underline{v}):\mbox{$\underline{v}$  $(w,\varepsilon)$-сопряжён и $(w,\varepsilon)$-двойственен}\}\not=\emptyset.\]
\item[\rm (b)] Пусть $(w,\varepsilon)\in\mathfrak{I}_\infty(N_+,N_-)$
и $\mathcal{F}=\{F'_a,F''_a:a\in A\}\in\mathbf{X}$. Тогда
$\mathcal{F}\in\bs{\mathcal{O}}_{(w,\varepsilon)}$ (соответственно, $\mathcal{F}\in\bs{\mathfrak{O}}_{(w,\varepsilon)}$), если и только если
\[(\delta(\mathcal{F}),\mathcal{F})\in\pmb{\mathbb{O}}_w\quad\mbox{(соответственно, $(\mathcal{F}^\dag,\mathcal{F})\in\pmb{\mathbb{O}}_w^\dag$)}\]
и для всех $\ell\in\NN$ выполнены условия
\[
\dim F''_{\sigma(\ell)}\cap \mathbf{V}_\pm/F'_{\sigma(\ell)}\cap \mathbf{V}_\pm
=
\left\{\begin{array}{ll}
1, & \mbox{если $\sigma(w_\ell)\prec\sigma(\ell)$ или $(w_\ell,\varepsilon_\ell)=(\ell,\pm1)$,} \\
0 & \mbox{в противном случае,}
\end{array}\right.
\]
где $\mathbf{V}_\pm=\langle e_\ell:\ell\in N_\pm\rangle_\mathbb{C}$
(соответственно, для достаточно большого $n\in\NN$
\[
\varsigma(\phi:F''_{\sigma(\ell)}\cap V_n)=\varsigma(\phi:F'_{\sigma(\ell)}\cap V_n)+\left\{
\begin{array}{ll}
(1,1), & \mbox{если $\sigma(w_\ell)\prec\sigma(\ell)$,} \\
(1,0), & \mbox{если $(w_\ell,\varepsilon_\ell)=(\ell,1)$,} \\
(0,1), & \mbox{если $(w_\ell,\varepsilon_\ell)=(\ell,-1)$,} \\
(0,0), & \mbox{если $\sigma(w_\ell)\succ\sigma(\ell)$,}
\end{array}
\right.
\]
где
$V_n=\langle e_k:k\leq n\rangle_\mathbb{C}$
и
$\varsigma(\phi:F)$ обозначает сигнатуру $\phi$ на $F/F\cap F^\dag$).
\item[\rm (c)] Подмножества $\bs{\mathcal{O}}_{(w,\varepsilon)}$ ($(w,\varepsilon)\in\mathfrak{I}_\infty(N_+,N_-)$) суть в точности $\mathbf{K}$-орбиты в $\mathbf{X}$. Подмножества $\bs{\mathfrak{O}}_{(w,\varepsilon)}$ ($(w,\varepsilon)\in\mathfrak{I}_\infty(N_+,N_-)$) суть в точности $\mathbf{G}^0$-орбиты в $\mathbf{X}$.
Далее, $\bs{\mathcal{O}}_{(w,\varepsilon)}\cap\bs{\mathfrak{O}}_{(w,\varepsilon)}$ есть единая $\mathbf{K}\cap\mathbf{G}^0$-орбита.
\end{itemize}
\end{proposition}

\subsection{Типы B, C, D}

\label{S4.3}

Предположим, что $\mathbf{V}$ обладает невырожденной симметрической или симплектической формой $\omega$, определяемой матрицей $\Omega$ как и в (\ref{omega}).
Пусть $\iota:\NN\to\NN$, $\ell\mapsto\ell^*$
удовлетворяет $\omega(e_\ell,e_{\ell^*})\not=0$ для всех $\ell$.

Зафиксируем  разложение $\NN=N_+\sqcup N_-$  такое, что $N_+,N_-$ или оба $\iota$-инвариантны или такие, что  $\iota(N_+)=N_-$.
Как и ранее, пусть $\phi$ и $\delta$ будут соответственно Эрмитовой формой и инволюцией в $\mathbf{V}$,
согласованными с этим разложением.
В следующей таблице представлены различные случаи.
%\begin{center}
%\begin{tabular}{|c|c|c|}
%\cline{2-3}
%\multicolumn{1}{c|}{ } & $\omega$ symmetric ($\epsilon:=1$) & $\omega$ symplectic %($\epsilon:=-1$) \\
%\hline
%$\iota(N_\pm)\subset N_\pm$ ($\eta:=1$) & type BD1 & type C2 \\
%\hline
%$\iota(N_\pm)=N_\mp$ ($\eta:=-1$) & type D3 & type C1 \\
%\hline
%\end{tabular}
%\end{center}
\begin{center}
\begin{tabular}{|c|c|c|}
\cline{2-3}
\multicolumn{1}{c|}{ } & \begin{tabular}{c} $\omega$ симметрическая \\ $\epsilon=1$ \end{tabular} & \begin{tabular}{c} $\omega$ симплектическая \\ $\epsilon=-1$ \end{tabular} \\
\hline
\begin{tabular}{c} $\iota(N_\pm)\subset N_\pm$ \\ $\eta=1$ \end{tabular} & тип BD1 & тип C2 \\
\hline
\begin{tabular}{c} $\iota(N_\pm)=N_\mp$ \\ $\eta=-1$ \end{tabular} & тип D3 & тип C1 \\
\hline
\end{tabular}
\end{center}

Обозначим через $\mathfrak{I}_\infty^{\eta,\epsilon}(N_+,N_-)\subset
\mathfrak{I}_\infty(N_+,N_-)$ подмножество пар
$(w,\varepsilon)$ такое, что
\begin{itemize}
\item[\rm (i)] $\iota w=w\iota$ (поэтому множество $\{\ell:w_\ell=\ell\}$  $\iota$-инвариантно);
\item[\rm (ii)] $\varepsilon_{\iota(k)}=\eta\varepsilon_k$ для всех $k\in\{\ell:w_\ell=\ell\}$;
\end{itemize}
и, если $\eta\epsilon=-1$:
\begin{itemize}
\item[\rm (iii)] 

 $w_k\not=\iota(k)$ для всех $k\in\mathbb{N}^*$.
\end{itemize}

Пусть $\mathcal{F}_\sigma$ является $\omega$-изотропным максимальным  обобщённым флагом, согласованным с $E$.
Таким образом $\sigma:\NN\to(A,\prec)$ --- биекция на
полностью упорядоченное множество $(A,\prec)$, обладающее (инволютивным) антиавтоморфизмом упорядоченных множеств
$\iota_A:(A,\prec)\to (A,\prec)$ таким, что $\sigma\iota=\iota_A\sigma$.
Следующее утверждение показывает, что $\mathbf{K}$-орбиты и $\mathbf{G}^0$-орбиты
инд-многообразия $\mathbf{X}_\omega:=\mathbf{X}_\omega(\mathcal{F}_\sigma,E)$
нумеруются элементами множества $\mathfrak{I}_\infty^{\eta,\epsilon}(N_+,N_-)$.

\begin{proposition}
\label{P4.3}
Рассмотрим базис $\underline{v}=(v_1,v_2,\ldots)$ в $\mathbf{V}$ такой, что
\begin{equation}
\label{4-3.1}
\omega(v_k,v_\ell)\not=0,\quad\mbox{если и только если}\quad \ell=\iota(k).
\end{equation}
\begin{itemize}
\item[\rm (a)] Для каждой пары $(w,\varepsilon)\in\mathfrak{I}^{\eta,\epsilon}_\infty(N_+,N_-)$ мы имеем
\begin{eqnarray*}
 & \bs{\mathcal{O}}_{(w,\varepsilon)}^{\eta,\epsilon}:=\bs{\mathcal{O}}_{(w,\varepsilon)}\cap\mathbf{X}_\omega=\{\mathcal{F}_\sigma(\underline{v}):
 \mbox{$\underline{v}$  $(w,\varepsilon)$-сопряжён} 
 \mbox{ и удовлетворяет (\ref{4-3.1})}\}\not=\emptyset, \\
 & \bs{\mathfrak{O}}_{(w,\varepsilon)}^{\eta,\epsilon}:=\bs{\mathfrak{O}}_{(w,\varepsilon)}\cap\mathbf{X}_\omega=\{\mathcal{F}_\sigma(\underline{v}):\mbox{$\underline{v}$ $(w,\varepsilon)$-двойственен} \mbox{ и удовлетворяет (\ref{4-3.1})}\}\not=\emptyset, \\
 & \bs{\mathcal{O}}_{(w,\varepsilon)}^{\eta,\epsilon}\cap\bs{\mathfrak{O}}_{(w,\varepsilon)}^{\eta,\epsilon}=\{\mathcal{F}_\sigma(\underline{v}):\mbox{$\underline{v}$ $(w,\varepsilon)$-сопряжён, $(w,\varepsilon)$-двойственен}\\
 & \mbox{и удовлетворяет (\ref{4-3.1})}\}\not=\emptyset.
\end{eqnarray*}
\item[\rm (b)] Подмножества $\bs{\mathcal{O}}_{(w,\varepsilon)}^{\eta,\epsilon}$
($(w,\varepsilon)\in\mathfrak{I}_\infty^{\eta,\epsilon}(N_+,N_-)$) суть
в точности $\mathbf{K}$-орбиты в $\mathbf{X}_\omega$.
Подмножества $\bs{\mathfrak{O}}_{(w,\varepsilon)}^{\eta,\epsilon}$
($(w,\varepsilon)\in\mathfrak{I}_\infty^{\eta,\epsilon}(N_+,N_-)$) суть
в точности $\mathbf{G}^0$-орбиты в $\mathbf{X}_\omega$.
Далее, $\bs{\mathcal{O}}_{(w,\varepsilon)}^{\eta,\epsilon}\cap\bs{\mathfrak{O}}_{(w,\varepsilon)}^{\eta,\epsilon}$ есть единая $\mathbf{K}\cap\mathbf{G}^0$-орбита.
\end{itemize}
\end{proposition}

\subsection{Структура инд-многообразий}

\label{S4.4}

В этом разделе мы напоминаем структуру инд-многообразий  на $\mathbf{X}$ и $\mathbf{X}_\omega$, см. \cite{Dimitrov-Penkov}.

Напомним, что $E=(e_1,e_2,\ldots)$ --- (счётный) базис в $\mathbf{V}$.
Зафиксируем $E$-соизмеримый максимальный обобщённый флаг $\mathcal{F}_\sigma$, соответствующий некоторой биекции $\sigma:\NN\to(A,\prec)$ на полностью упорядоченное множество, и положим
$\mathbf{X}=\mathbf{X}(\mathcal{F}_\sigma,E)$.

Пусть $V_n:=\langle e_1,\ldots,e_n\rangle_\mathbb{C}$ и
пусть $X_n$ обозначает многообразие полных флагов в $V_n$ определённых как  в (\ref{2.2.X}).
Имеются естественные включения
$V_n\subset V_{n+1}$
и
\begin{equation}
\label{4.4-1}
\mathrm{GL}(V_n)\cong\{g\in \mathrm{GL}(V_{n+1}):g(V_n)=V_n,\ g(e_{n+1})=e_{n+1}\}\subset \mathrm{GL}(V_{n+1}).
\end{equation}
Зададим $\mathrm{GL}(V_n)$-эквивариантное вложение
\[\iota_n=\iota_n(\sigma):X_n\to X_{n+1},\ (F_k)_{k=0}^n\mapsto (F'_k)_{k=0}^{n+1},\]
полагая
\[F'_k:=\left\{
\begin{array}{ll}
F_k & \mbox{, если $a_k\prec\sigma(n+1)$} \\
F_{k-1}\oplus\langle e_{n+1}\rangle_\mathbb{C} & \mbox{, если $a_k\succeq \sigma(n+1),$}
\end{array}
\right.  \]
где $a_1\prec a_2\prec\ldots\prec a_{n+1}$ --- элементы множества $\{\sigma(\ell):1\leq \ell\leq n+1\}$, записанные в порядке возрастания.
Таким образом, мы имеем цепь вложений (являющихся морфизмами алгебраических многообразий)
%\[
%\xymatrix{
%\cdots \ar@{^{(}->}[r] & X_{n-1} \ar@{^{(}->}[r]^{\iota_{n-1}}
% & X_{n} \ar@{^{(}->}[r]^{\iota_{n}}
% & X_{n+1} \ar@{^{(}->}[r]^{\iota_{n+1}}
% & \cdots
%}
%\]
\[\cdots \hookrightarrow X_{n-1} \stackrel{\iota_{n-1}}{\hookrightarrow} X_n \stackrel{\iota_{n}}{\hookrightarrow} X_{n+1} \stackrel{\iota_{n+1}}{\hookrightarrow} \cdots,\]
и $\mathbf{X}$ получается как прямой предел
\[\mathbf{X}=\mathbf{X}(\mathcal{F}_\sigma,E)=\lim_\to X_n.\]
В частности, для каждого $n$ мы получаем вложение $\hat{\iota}_n:X_n\hookrightarrow \mathbf{X}$, и, отождествляя $X_n$ со своим образом в $\mathbf{X}$, мы можем представить $\mathbf{X}$ в виде объединения
$\mathbf{X}=\bigcup_{n\geq 1}X_n$.
При этом, каждый обобщённый флаг $\mathcal{F}\in\mathbf{X}$ принадлежит всем $X_n$, начиная с некоторого номера $n_\mathcal{F}$. Например, $\mathcal{F}_\sigma\in X_n$ для всех $n\geq 1$.

Базис $\underline{v}=(v_1,\ldots,v_n)$ в $V_n$ может быть дополнен до базиса в $\mathbf{V}$, обозначаемого $\underline{\hat{v}}:=(v_1,\ldots,v_n,e_{n+1},e_{n+2},\ldots)$ так, что
\begin{equation}
\label{completed-basis}
\hat{\iota}_n(\mathcal{F}(v_{\tau_1},\ldots,v_{\tau_n}))=\mathcal{F}_\sigma(\underline{\hat{v}})
\end{equation}
(в обозначениях из разделов \ref{S2.2}--\ref{S2.3}),
где $\tau=\tau^{(n)}\in\mathfrak{S}_n$ --- перестановка такая, что
$\sigma(\tau^{(n)}_1)\prec\ldots\prec\sigma(\tau^{(n)}_n)$.

Напомним, что инд-топология на $\mathbf{X}$ определяется объявлением подмножества $\mathbf{Z}\subset\mathbf{X}$ открытым (соответственно, замкнутым), если каждое пересечение $\mathbf{Z}\cap X_n$ открыто (соответственно, замкнуто).

Ясно, что структура инд-многообразия на $\mathbf{X}$ не меняется, если последовательность $(X_n,\iota_n)_{n\geq 1}$ заменить на подпоследовательность $(X_{n_k},\iota'_k)_{k\geq 1}$,
где $\iota'_k:=\iota_{n_{k+1}-1}\circ\cdots\circ\iota_{n_k}$.

Для типа A3 (используя обозначения из раздела \ref{S2.1}) подпространство $V_n\subset\mathbf{V}$ обладает ограничениями $\phi$ и $\delta$
, поэтому мы можем определить $K_n,G_n^0\subset \mathrm{GL}(V_n)$ как в разделе \ref{S3.2}
и с условием, чтобы включения в (\ref{4.4-1}) задавали  естественные включения
$K_n\subset K_{n+1}$ и $G_n^0\subset G_{n+1}^0$.

Далее предположим, что пространство $\mathbf{V}$ обладает невырожденной симплектической или симметрической формой $\omega$ определяемой матрицей $\Omega$ из (\ref{omega}).
Блоки $J_1,J_2,\ldots$ в матрице $\Omega$ имеют размер $1$ или $2$.
Зададим $n_k:=|J_1|+\ldots+|J_k|$
так, чтобы ограничение $\omega$ на каждое подпространство $V_{n_k}$ было невырождено.
Поэтому для типов A1, A2, BD1, C1, C2, и D3 мы можем определить
подгруппы $K_{n_k},G_{n_k}^0\subset \mathrm{GL}(V_{n_k})$
как в разделе \ref{S3}
и так, чтобы включения в (\ref{4.4-1}) задавали естественные включения
\[K_{n_k}\subset K_{n_{k+1}}\quad\mbox{и}\quad G^0_{n_k}\subset G^0_{n_{k+1}}.\]

Далее, подмногообразие $(X_{n_k})_\omega\subset X_{n_k}$ изотропных флагов (относительно формы $\omega$) может быть определено как и в (\ref{2.2.Xomega}).
Для $\omega$-изотропного обобщённого флага $\mathcal{F}_\sigma$  включение
$\iota'_k:X_{n_k}\hookrightarrow X_{n_{k+1}}$ отображает $(X_{n_k})_\omega$ в $(X_{n_{k+1}})_\omega$, и мы имеем
\[\mathbf{X}_\omega=\mathbf{X}_\omega(\mathcal{F}_\sigma,E)=\bigcup_{k\geq 1}(X_{n_k})_\omega\quad\mbox{и}\quad (X_{n_k})_\omega=\mathbf{X}_\omega\cap X_{n_k}\ \mbox{для всех $k\geq 1$}.\]
В частности, $\mathbf{X}_\omega$ --- замкнутое инд-подмногообразие $\mathbf{X}$
(как утверждалось в предложении \ref{P2.3-3}).

\subsection{Доказательства}

\label{S4.proofs}

\begin{proof}[Доказательство предложения \ref{P4-1.1}]
Пусть $\mathcal{F}=\{F'_a,F''_a:a\in A\}=\mathcal{F}_\sigma(\underline{v})$ для базиса $\underline{v}=(v_1,v_2,\ldots)$ в $\mathbf{V}$.
Пусть $w\in\mathfrak{I}_\infty^\epsilon(\iota)$.
Если базис $\underline{v}$ является $w$-двойственным,  то
\[
(F'_a)^\perp=\langle v_\ell:\sigma(w_\ell)\succeq a\rangle_\mathbb{C}
\quad\mbox{и}\quad
(F''_a)^\perp=\langle v_\ell:\sigma(w_\ell)\succ a\rangle_\mathbb{C}\,,
\]
поэтому $\mathcal{F}^\perp=\mathcal{F}_{\sigma^\perp\iota w}(\underline{v})$;
это даёт
$(\mathcal{F}^\perp,\mathcal{F})\in\pmb{\mathbb{O}}^\perp_{w\iota}$.
Если $\underline{v}$ --- $w$-сопряжённый, то
\[
\gamma(F'_a)=\langle v_\ell:\sigma(w_\ell)\prec a\rangle_\mathbb{C}
\quad\mbox{и}\quad
\gamma(F''_a)=\langle v_\ell:\sigma(w_\ell)\preceq a\rangle_\mathbb{C}\,,
\]
следовательно $\gamma(\mathcal{F})=\mathcal{F}_{\sigma w}(\underline{v})$
и
$(\gamma(\mathcal{F}),\mathcal{F})\in\pmb{\mathbb{O}}^\gamma_{w\iota}$.
Это доказывает включение в предложении \ref{P4-1.1}\,{\rm (b)}.
Отметим что эти включения указывают, в частности,  на то, что подмножества $\bs{\mathcal{O}}_w$, так же как и $\bs{\mathfrak{O}}_w$, попарно различны.

Для $w\in\mathfrak{I}_{n_k}^\epsilon$ мы определим $\hat{w}:\mathbb{N}^*\to\mathbb{N}^*$ следующим образом:
\[\hat{w}(\ell)=\left\{
\begin{array}{ll}
\tau w\tau^{-1}(\ell), & \mbox{если $\ell\leq n_k$,} \\
\iota(\ell), & \mbox{если $\ell\geq n_k+1$,}
\end{array}
\right.\]
где $\tau=\tau^{(n_k)}:\{1,\ldots,n_k\}\to\{1,\ldots,n_k\}$ --- перестановка такая, что $\sigma(\tau_1)\prec\ldots\prec\sigma(\tau_{n_k})$.
Легко видеть, что мы получаем корректно определённое (инъективное) отображение
$j_k:\mathfrak{I}_{n_k}^\epsilon\to\mathfrak{I}_\infty^\epsilon(\iota)$, $j_k(w):=\hat{w}$,
и
\begin{equation}
\label{4.5.0}
\mathfrak{I}_\infty^\epsilon(\iota)=\bigcup_{k\geq 1}j_k(\mathfrak{I}_{n_k}^\epsilon).\end{equation}
Далее, если
дан базис $\underline{v}=(v_1,\ldots,v_{n_k})$ в $V_{n_k}$ и
базис $\underline{\hat{v}}$ в $\mathbf{V}$, получаемый добавлением векторов $e_\ell$ для $\ell\geq n_k+1$,
то импликация
\begin{eqnarray}
\label{4.5.1}
 & \mbox{$(v_{\tau_1},\ldots,v_{\tau_{n_k}})$ --- $w$-двойственный (соответственно, $w$-сопряжённый)} \\
\nonumber
 & \Rightarrow \
\mbox{$\underline{\hat{v}}$ --- $\hat{w}$-двойственный (соответственно, $\hat{w}$-сопряжённый)}
\end{eqnarray}
очевидно следует из нашей конструкции.
Отметим, что
\begin{equation}
\label{T-proof.1}
\bs{\mathcal{O}}_{\hat{w}}\cap X_{n_k}=\mathcal{O}_{w}
\quad
\mbox{и}
\quad
\bs{\mathfrak{O}}_{\hat{w}}\cap X_{n_k}=\mathfrak{O}_{w},
\end{equation}
где $\mathcal{O}_w,\mathfrak{O}_w\subset X_{n_k}$ --- орбиты из  определения \ref{DA12};
в самом деле, включения $\supset$ в (\ref{T-proof.1}) следуют из
(\ref{completed-basis}) и (\ref{4.5.1}),
в то время как включения $\subset$ следуют из
предложения \ref{P1}\,{\rm (c)}
и факта, что подмножества $\bs{\mathcal{O}}_{\hat{w}}$, так же как и $\bs{\mathfrak{O}}_{\hat{w}}$,
попарно различны.
Части {\rm (a)} и {\rm (c)} предложения \ref{P4-1.1} теперь следуют
из (\ref{4.5.0})--(\ref{T-proof.1}) и предложения \ref{P1}\,{\rm (a)}, {\rm (c)}. Из предложения \ref{P4-1.1}\,{\rm (a)} мы выводим, что равенства в предложениях \ref{P4-1.1}\,{\rm (b)} справедливы, и завершаем доказательство.
\end{proof}

\begin{proof}[Доказательство предложения \ref{P4.2-2}]
Для каждого $n\geq 1$ положим $p_n=|N_+\cap\{1,\ldots,n\}|$  и $q_n=|N_-\cap\{1,\ldots,n\}|$.

Пусть $\mathcal{F}=\{F'_a,F''_a:a\in A\}=\mathcal{F}_\sigma(\underline{v})$ для некоторого базиса $\underline{v}=(v_1,v_2,\ldots)$  $\mathbf{V}$.
Пусть $(w,\varepsilon)\in\mathfrak{I}_\infty(N_+,N_-)$.
Если $\underline{v}$ $(w,\varepsilon)$-сопряжён, то
\[
\delta(F'_a)=\langle v_\ell:\sigma(w_\ell)\prec a\rangle_\mathbb{C}
\quad\mbox{и}\quad
\delta(F''_a)=\langle v_\ell:\sigma(w_\ell)\preceq a\rangle_\mathbb{C}
\]
так, что $(\delta(\mathcal{F}),\mathcal{F})=(\mathcal{F}_{\sigma w}(\underline{v}),\mathcal{F}_\sigma(\underline{v}))\in\pmb{\mathbb{O}}_w$.
Далее,
\[\left\{
\begin{array}{ll}
\mbox{$F''_{\sigma(\ell)}\cap\mathbf{V}_+/F'_{\sigma(\ell)}\cap\mathbf{V}_+=\langle v_\ell\rangle_\mathbb{C}$, $F''_{\sigma(\ell)}\cap\mathbf{V}_-=F'_{\sigma(\ell)}\cap\mathbf{V}_-$} & \mbox{\!\!\!, если $(w_\ell,\varepsilon_\ell)=(\ell,+1)$,} \\[1.5mm]
\mbox{$F''_{\sigma(\ell)}\cap\mathbf{V}_-/F'_{\sigma(\ell)}\cap\mathbf{V}_-=\langle v_\ell\rangle_\mathbb{C}$, $F''_{\sigma(\ell)}\cap\mathbf{V}_+=F'_{\sigma(\ell)}\cap\mathbf{V}_+$} & \mbox{\!\!\!, если $(w_\ell,\varepsilon_\ell)=(\ell,-1)$,} \\[1.5mm]
\mbox{$F''_{\sigma(\ell)}\cap\mathbf{V}_+/F'_{\sigma(\ell)}\cap\mathbf{V}_+=\langle v_\ell+v_{w_\ell}\rangle_\mathbb{C}$\,,} \\[1mm]
\mbox{\qquad $F''_{\sigma(\ell)}\cap\mathbf{V}_-/F'_{\sigma(\ell)}\cap\mathbf{V}_-=\langle v_\ell- v_{w_\ell}\rangle_\mathbb{C}$} & \mbox{\!\!\!, если $\sigma(w_\ell)\prec\sigma(\ell)$,} \\[1.5mm]
F''_{\sigma(\ell)}\cap \mathbf{V}_+=F'_{\sigma(\ell)}\cap \mathbf{V}_+,\ F''_{\sigma(\ell)}\cap \mathbf{V}_+=F'_{\sigma(\ell)}\cap \mathbf{V}_+ & \mbox{\!\!\!, если $\sigma(w_\ell)\succ\sigma(\ell)$,}
\end{array}
\right.\]
что доказывает формулу для $\dim F''_{\sigma(\ell)}\cap \mathbf{V}_\pm/F'_{\sigma(\ell)}\cap \mathbf{V}_\pm$ из предложения \ref{P4.2-2}\,{\rm (b)}.
Если $\underline{v}$  $(w,\varepsilon)$-двойственен, то мы получаем одновременно
\[
(F'_a)^\dag=\langle v_\ell:\sigma(w_\ell)\succeq a\rangle_\mathbb{C}
\quad\mbox{и}\quad
(F''_a)^\dag=\langle v_\ell:\sigma(w_\ell)\succ a\rangle_\mathbb{C}\,.
\]
Поэтому $(\mathcal{F}^\dag,\mathcal{F})=(\mathcal{F}_{\sigma^\dag w}(\underline{v}),\mathcal{F}_\sigma(\underline{v}))\in\pmb{\mathbb{O}}_w^\dag$.
Для достаточно большого $n\geq 1$  мы имеем
$(w_\ell,\varepsilon_\ell)=(\ell,\pm1)$ для всех $\ell\in N_\pm\cap\{n+1,n+2,\ldots\}$
и $v_\ell=e_\ell$ для всех $\ell\geq n+1$.
Таким образом, пара $(\check{w},\check{\varepsilon}):=(w|_{\{1,\ldots,n\}},\varepsilon|_{\{1,\ldots,n\}})$ принадлежит $\mathfrak{I}_n(p_n,q_n)$, в то время как согласно (\ref{completed-basis})  имеем
\[\mathcal{F}=\mathcal{F}(v_{\tau_1},\ldots,v_{\tau_n}).\]
Базис $(v_{\tau_1},\ldots,v_{\tau_n})$ в $V_n$  $(\tau^{-1}\check{w}\tau,\check{\varepsilon}\tau)$-двойственен, если $\underline{v}$  $(w,\varepsilon)$-двойственен; последняя формула в предложении \ref{P4.2-2}\,{\rm (b)} теперь следует из предложения \ref{P2}\,{\rm (b)} и данного наблюдения. В итоге, это доказывает условие необходимости в предложении \ref{P4.2-2}\,{\rm (b)}, которое гарантирует, в частности, что подмножества $\bs{\mathcal{O}}_{(w,\varepsilon)}$, так же как и подмножества $\bs{\mathfrak{O}}_{(w,\varepsilon)}$,
попарно различны.
Достаточность в предложении \ref{P4.2-2}\,{\rm (b)} получится как только мы докажем
предложение \ref{P4.2-2}\,{\rm (a)}.

Для $(w,\varepsilon)\in\mathfrak{I}_n(p_n,q_n)$
положим
\begin{equation}
\label{4.5.4}
\hat{w}(\ell)=\left\{
\begin{array}{ll}
\tau w\tau^{-1}(\ell) & \mbox{, если $\ell\leq n$,} \\
\ell & \mbox{, если $\ell\geq n+1$}
\end{array}
\right.\quad\mbox{для всех $\ell\in\mathbb{N}^*$,}
\end{equation}
где $\tau=\tau^{(n)}\in\mathfrak{S}_n$ как и в (\ref{completed-basis}), и
\begin{equation}
\label{4.5.5}\hat{\varepsilon}(\ell)=\left\{
\begin{array}{ll}
\varepsilon\tau^{-1}(\ell) & \mbox{, если $\ell\leq n$,} \\
1 & \mbox{, если $\ell\geq n+1$, $n\in N_+$,} \\
-1 & \mbox{, если $\ell\geq n+1$, $n\in N_-$}
\end{array}
\right.
\end{equation}
для всех $\ell\in\mathbb{N}^*$ таких, что $\hat{w}_\ell=\ell$.
Легко проверить, что $(\hat{w},\hat{\varepsilon})\in\mathfrak{I}_\infty(N_+,N_-)$,
и, что полученное таким образом отображение $j_n:\mathfrak{I}_n(p_n,q_n)\to\mathfrak{I}_\infty(N_+,N_-)$ инъективно
и
\[
\mathfrak{I}_\infty(N_+,N_-)=\bigcup_{n\geq 1}j_n(\mathfrak{I}_n(p_n,q_n)).
\]
Более того, из нашего построения следует, что, если дан базис в $\underline{v}=(v_1,\ldots,v_n)$ в $V_n$ и базис $\underline{\hat{v}}$ в $\mathbf{V}$ получается добавлением векторов $e_\ell$ для $\ell\geq n+1$, то верна импликация:
\begin{eqnarray*}
 & \mbox{$(v_{\tau_1},\ldots,v_{\tau_n})$ --- $(w,\varepsilon)$-сопряжённый (соответственно, двойственный)} \\
 & \Rightarrow \
\mbox{$\underline{\hat{v}}$ --- $(\hat{w},\hat{\varepsilon})$-сопряжённый (соответственно, двойственный)}.
\end{eqnarray*}
Как и в доказательстве предложения \ref{P4-1.1} мы получаем равенства
\begin{equation}
\label{T-proof.2}
\bs{\mathcal{O}}_{(\hat{w},\hat{\varepsilon})}\cap X_{n}=\mathcal{O}_{(w,\varepsilon)}
\quad
\mbox{и}
\quad
\bs{\mathfrak{O}}_{(\hat{w},\hat{\varepsilon})}\cap X_{n}=\mathfrak{O}_{(w,\varepsilon),}
\end{equation}
где $\mathcal{O}_{(w,\varepsilon)},\mathfrak{O}_{(w,\varepsilon)}\subset X_{n}$ как и в определении \ref{DA3}.
Части {\rm (a)} и {\rm (c)} предложения \ref{P4.2-2} тогда следуют
из предложений \ref{P2}\,{\rm (a)} и {\rm (c)}.
\end{proof}

\begin{proof}[Доказательство предложения \ref{P4.3}]
Пусть $n\in\{n_1,n_2,\ldots\}$ (где $n_k=|J_1|+\ldots+|J_k|$ как и ранее) и $(p_n,q_n)=(|N_+\cap\{1,\ldots,n\}|,|N_-\cap\{1,\ldots,n\}|)$
и пусть $\tau=\tau^{(n)}:\{1,\ldots,n\}\to\{1,\ldots,n\}$
--- перестановка такая, что $\sigma(\tau_1)\prec\ldots\prec\sigma(\tau_n)$.
Поскольку обобщённый флаг $\mathcal{F}_\sigma$  $\omega$-изотропен,
должно быть выполнено
\[\iota(\tau_\ell)=\tau_{n-\ell+1}\quad\mbox{для всех $\ell\in\{1,\ldots,n\}$}.\]
Это наблюдение легко влечёт то, что отображение $j_n$ определённое в доказательстве
предложения \ref{P4.2-2}
ограничивается до корректно определённого инъективного отображения
\[j_n:\mathfrak{I}^{\eta,\epsilon}_n(p_n,q_n)\to \mathfrak{I}^{\eta,\epsilon}_\infty(N_+,N_-)\]
такого, что
\[\mathfrak{I}^{\eta,\epsilon}_\infty(N_+,N_-)=\bigcup_{k\geq 1}
j_{n_k}(\mathfrak{I}^{\eta,\epsilon}_{n_k}(p_{n_k},q_{n_k})).\]
Из (\ref{T-proof.2}) для $(\hat{w},\hat{\varepsilon})=j_n(w,\varepsilon)$ мы получаем
\begin{equation}
\label{T-proof.3}
\bs{\mathcal{O}}^{\eta,\epsilon}_{(\hat{w},\hat{\varepsilon})}\cap (X_{n})_\omega=\mathcal{O}^{\eta,\epsilon}_{(w,\varepsilon)}
\quad
\mbox{и}
\quad
\bs{\mathfrak{O}}^{\eta,\epsilon}_{(\hat{w},\hat{\varepsilon})}\cap (X_{n})_\omega=\mathfrak{O}^{\eta,\epsilon}_{(w,\varepsilon)}.
\end{equation}
Предложение \ref{P4.3} легко следует из этого факта и предложения \ref{P3}.
\end{proof}



\section{Следствия}

\label{S5}

%Note that our parametrization of orbits does not depend on the choice of the $E$-compatible maximal generalized flag $\mathcal{F}_\sigma$, which yields:
\begin{corollary}
	Отображение двойственности $\Xi$ из теоремы \ref{theorem-1} $(b)$ зависит от выбора $\mathbf{G}$, $\mathbf{B}$, $\mathbf{K}$ и $\mathbf{G}^0$, но не от выбора базиса $E$, используемого выше в построении $\mathbf{G}$, $\mathbf{B}$, $\mathbf{K}$ и $\mathbf{G}^0$. В частности, $\Xi$ не зависит от исчерпания $\mathbf{X}= \bigcup_{n \geq 1} X_n$, определяемого $E$ и упоминемого в теореме \ref{theorem-1} $(b)$.
\end{corollary}

\begin{proof}
	Утверждение следует сразу из коммутативности диаграммы (\ref{diag})  и из наблюдения, что, если  $\mathbf{X}= \bigcup_{n \geq 1} X_n$ и  $\mathbf{X}= \bigcup_{n' \geq 1} X'_{n'}$, суть два исчерпания инд-многообразия  $\mathbf{X}$, то, для любых $n_0$ и $n'_0$ существуют $n_1$ и $n'_1$ такие, что $X_{n_0}\cup X'_{n'_0} \subset X_{n_1}$ и $X_{n_0}\cup X'_{n'_0} \subset X'_{n'_1}$.
\end{proof}
Второе следствие утверждает, что параметризация $\mathbf{K}$- и $\mathbf{G}^0$-орбит на $\mathbf{G}/\mathbf{B}$ на самом деле зависит только от тройки $\left(\mathbf{G},\mathbf{K},\mathbf{G}^0\right)$, но не зависит от выбора инд-многообразия  $\mathbf{G}/\mathbf{B}$.

\begin{corollary}
\label{C1}
Пусть $E,\mathbf{G},\mathbf{K},\mathbf{G}^0$ --- такие же, как  в  разделе \ref{S2.1}.
Пусть $\mathcal{F}_{\sigma_j}$ ($j=1,2$) --- два
$E$-соизмеримых
максимальных обобщённых флага, которые $\omega$-изотропны для типов B,C,D,
и пусть $\mathbf{X}_j=\mathbf{G}/\mathbf{B}_{\mathcal{F}_{\sigma_j}}$.
Тогда существуют естественные биекции
\[
\mathbf{X}_1/\mathbf{K}\cong\mathbf{X}_2/\mathbf{K}
\quad\mbox{и}\quad
\mathbf{X}_1/\mathbf{G}^0\cong\mathbf{X}_2/\mathbf{G}^0,
\]
 которые коммутируют с двойственностью из теоремы \ref{theorem-1}.
\end{corollary}

Далее, прямой подсчёт параметров даёт:

\begin{corollary}
\label{C2}
В следствии \ref{C1}
множества орбит $\mathbf{X}_j/\mathbf{K}$ и $\mathbf{X}_j/\mathbf{G}^0$
всегда бесконечны.
\end{corollary}

%Corollary \ref{C2} makes the existence of open (resp., closed) orbits not automatic.
%The next statement answers this question.

Важно отметить, что, несмотря на следствие \ref{C1}, топологические свойства орбит на $\mathbf{G}/\mathbf{B}$ не одинаковы для различных Борелевских инд-подгрупп $\mathbf{B}\subset\mathbf{G}$.  Следующее следствие устанавливает критерии для существования открытых и замкнутых орбит на $\mathbf{G}/\mathbf{B}=\mathbf{X}\left(\mathcal{F}_\sigma,E\right)$.

\begin{corollary}
\label{C3}
Пусть $E,\mathbf{G},\mathbf{K},\mathbf{G}^0$ --- такие, как в разделе \ref{S2.1},
и пусть $\mathcal{F}_\sigma$ --- $E$-соизмеримый максимальный обобщённый флаг, $\omega$-изотропный для типов B,C,D, где $\sigma:\mathbb{N}^*\to (A,\prec)$ --- биекция на полностью упорядоченное множество.
Пусть $\mathbf{X}=\mathbf{G}/\mathbf{B}_{\mathcal{F}_\sigma}$; то есть, $\mathbf{X}=\mathbf{X}(\mathcal{F}_\sigma,E)$ для типа A, и $\mathbf{X}=\mathbf{X}_\omega(\mathcal{F}_\sigma,E)$ для типов B,C,D.

\begin{itemize}
\item[\rm ($\mbox{a}_1$)] Для типа A1,
$\mathbf{X}$ обладает открытой $\mathbf{K}$-орбитой (эквивалентно, замкнутой $\mathbf{G}^0$-орбитой), если и только если $\iota(\ell)=\ell$ для всех $\ell\gg 1$ (то есть, если матрица $\Omega$ из (\ref{omega}) содержит конечное число диагоналных блоков размера 2).%; in that case the open $\mathbf{K}$-orbit is $\bs{\mathcal{O}}_\mathrm{id}$.
\item[\rm ($\mbox{a}_2$)]
Для типа A2, $\mathbf{X}$ обладает открытой $\mathbf{K}$-орбитой (эквивалентно, замкнутой $\mathbf{G}^0$-орбитой),если и только если
для всех $\ell\gg 1$ элементы
$\sigma(2\ell-1),\sigma(2\ell)$ --- последовательные в  $A$, и
число $|\{k<2\ell-1:\sigma(k)\prec\sigma(2\ell-1)\}|$ --- чётное.
\item[\rm ($\mbox{a}'_{12}$)]
Для типов A1 и A2,
$\mathbf{X}$ имеет не более одной замкнутой $\mathbf{K}$-орбиты
(эквивалентно,не более одной открытой $\mathbf{G}^0$-орбиты).  $\mathbf{X}$ имеет замкнутую $\mathbf{K}$-орбиту (эквивалентно, открытую $\mathbf{G}^0$-орбиту), если и только если $\mathbf{X}$ содержит $\omega$-изотропные обобщённые флаги. Это последнее условие эквивалентно существованию инволютивного автоморфизма упорядоченных множеств $\iota_A:(A,\prec)\to(A,\prec)$ такого, что
$\iota_A\sigma(\ell)=\sigma\iota(\ell)$ для всех $\ell\gg 1$.
\item[\rm ($\mbox{a}_3$)] Для типа A3,
$\mathbf{X}$ имеет всегда беконечно много замкнутых $\mathbf{K}$-орбит (эквивалентно, бесконечно много $\mathbf{G}^0$-орбит).
%, namely $\bs{\mathcal{O}}_{(\mathrm{id},\varepsilon)}$ and $\bs{\mathfrak{O}}_{(\mathrm{id},\varepsilon)}$.
%corresponding to elements of $\mathfrak{I}_\infty(N_+,N_-)$ of the form $(\mathrm{id},\varepsilon)$.
$\mathbf{X}$ обладает открытой $\mathbf{K}$-орбитой (эквивалентно, замкнутой $\mathbf{G}^0$-орбитой),
если и только если $d:=\min\{|N_+|,|N_-|\}<\infty$ и $\mathcal{F}_\sigma$
содержит подпространства размерности $d$ и коразмерности  $d$.
\item[\rm (bcd)] Для типов B,C,D,
$\mathbf{X}$ имеет всегда бесконечно много замкнутых  $\mathbf{K}$-орбит (эквивалентно, бесконечно много открытых $\mathbf{G}^0$-орбит).
Для типов C1 и D3, $\mathbf{X}$ не имеет открытых $\mathbf{K}$-орбит (эквивалентно, не содержит замкнутых $\mathbf{G}^0$-орбит).
Для типов BD1 и C2, $\mathbf{X}$ имеет открытую $\mathbf{K}$-орбиту (соответственно, замкнутую $\mathbf{G}^0$-орбиту),
если и только если $d:=\min\{|N_+|,|N_-|\}<\infty$ и $\mathcal{F}_\sigma$ имеет подпространство размерности $d$  ( эквивалентно, содержит подпространства коразмерности $d$).
\end{itemize}
\end{corollary}

\begin{proof}
Утверждения следуют из замечаний \ref{R-open} и \ref{R-closed}, предложений \ref{P4-1.1}, \ref{P4.2-2}, \ref{P4.3} и равенств (\ref{T-proof.1}), (\ref{T-proof.2}), (\ref{T-proof.3}).
\end{proof}

%Note that, despite of the bijections of Corollary \ref{C1},
%it appears in Corollary \ref{C3} that the existence
%of open and closed orbits in $\mathbf{X}=\mathbf{G}/\mathbf{B}_{\mathcal{F}_\sigma}$
%may depend on the form of the maximal generalized flag $\mathcal{F}_\sigma$.
%The next statement follows from Corollary \ref{C3}.

\begin{corollary}
Единственный случай, в котором $\mathbf{X}$ имеет одновременно открытые и замкнутые $\mathbf{K}$-орбиты (эквивалентно, открытые и замкнутые $\mathbf{G}^0$-орбиты), это типы
A3, BD1, C2 когда $d:=\min\{|N_+|,|N_-|\}<\infty$ и $\mathcal{F}_\sigma$
содержит подпространства размерности $d$ и коразмерности $d$.
\end{corollary}

\section*{Список обозначений}

\begin{itemize}
\item[\S\ref{S.intro}:] $\mathbb{N}^*$, $|A|$, $\mathfrak{S}_n$, $\mathfrak{S}_\infty$, $(k;\ell)$

\smallskip
\item[\S\ref{S2.1}:] $\mathbf{G}(E)$, $\mathbf{G}(E,\omega)$, $\Omega$, $\omega$, $\gamma$, $\Phi$, $\phi$, $\delta$

\smallskip
\item[\S\ref{S2.2}:] $\mathcal{F}(v_1,\ldots,v_n)$, $\mathbb{O}_w$

\smallskip
\item[\S\ref{S2.3}:] $\mathcal{F}_\sigma(\underline{v})$, $\mathcal{F}_\sigma$,
$\mathbf{P}_\mathcal{F}$, $\mathbf{B}_\mathcal{F}$, $\mathbf{X}(\mathcal{F},E)$, $(\pmb{\mathbb{O}}_{\tau,\sigma})_w$, $\mathbf{X}_\omega(\mathcal{F},E)$

\smallskip
\item[\S\ref{S3.1}:] $\mathcal{F}^\perp$, $\gamma(\mathcal{F})$, $\mathfrak{I}_n$, $\mathfrak{I}'_n$, $\mathcal{O}_w$, $\mathfrak{O}_w$

\smallskip
\item[\S\ref{S3.2}:] $\delta(\mathcal{F})$, $\mathcal{F}^\dag$, $\varsigma(\phi:\mathcal{F})$, $\varsigma(\delta:\mathcal{F})$, $\varsigma(w,\varepsilon)$, $\mathfrak{I}_n(p,q)$, $\mathcal{O}_{(w,\varepsilon)}$, $\mathfrak{O}_{(w,\varepsilon)}$

\smallskip
\item[\S\ref{S3.3}:] $\mathfrak{I}^{\eta,\epsilon}_n(p,q)$, $\mathcal{O}^{\eta,\epsilon}_{(w,\varepsilon)}$, $\mathfrak{O}^{\eta,\epsilon}_{(w,\varepsilon)}$
\smallskip
\item[\S\ref{S4.1}:] $\iota$, $\mathfrak{I}_\infty(\iota)$, $\mathfrak{I}'_\infty(\iota)$, $\bs{\mathcal{O}}_w$, $\bs{\mathfrak{O}}_w$, $(A^*,\prec^*)$, $\sigma^\perp$, $\mathbf{X}^\perp$, $\mathbf{X}^\gamma$, $\pmb{\mathbb{O}}_w^\perp$, $\pmb{\mathbb{O}}_w^\gamma$
\smallskip
\item[\S\ref{S4.2}:] $\mathfrak{I}_\infty(N_+,N_-)$, $\bs{\mathcal{O}}_{(w,\varepsilon)}$, $\bs{\mathfrak{O}}_{(w,\varepsilon)}$, $\sigma^\dag$, $\mathbf{X}^\dag$, $\pmb{\mathbb{O}}_w$, $\pmb{\mathbb{O}}_w^\dag$
\smallskip
\item[\S\ref{S4.3}:] $\mathfrak{I}_\infty^{\eta,\epsilon}(N_+,N_-)$, $\bs{\mathcal{O}}_{(w,\varepsilon)}^{\eta,\epsilon}$, $\bs{\mathfrak{O}}_{(w,\varepsilon)}^{\eta,\epsilon}$
%\smallskip
%\item[\S\ref{S4.4}:] $\tau^{(n)}$
\end{itemize}

\begin{thebibliography}{gz}
\bibitem{Baranov} А. А. Baranov: Finitary simple Lie algebras.  J. Algebra \textbf{219} (1999), 299--329.
\bibitem{Berger} M. Berger: Les espaces sym\'etriques non compacts. Ann. Sci. \'Ecole Norm. Sup. \textbf{74} (1957) 85--177.
\bibitem{Brion} M. Brion and G. Helminck: On orbit closures of symmetric subgroups in flag varities. Canad. J. Math. \textbf{52} (2000), 265--292.
\bibitem{Dimitrov-Penkov} I. Dimitrov and I. Penkov: Ind-varieties of generalized flags as homogeneous spaces for classical ind-groups. Int. Math. Res. Not. \textbf{2004} (2004), 2935--2953.
\bibitem{Fresse-Penkov} L. Fresse and I. Penkov: Schubert decompositions for ind-varieties of generalized flags. Asian J. Math., to appear.
\bibitem{Gindikin-Matsuki} S. Gindikin and T. Matsuki: Stein extensions of Riemannian symmetric spaces and dualities of orbits on flag manifolds. Transform. Groups \textbf{8} (2003), 333--376.
\bibitem{FHW} G. Fels, A. Huckleberry, and J. A. Wolf: Cycle spaces of flag domains. A complex geometric viewpoint. Progr. Math., vol. 245, Birkh\"auser, Boston, MA, 2006.
\bibitem{Ignatyev-Penkov} M. V. Ignatyev and I. Penkov: Ind-varieties of generalized flags: a survey of results. Preprint (2016).
\bibitem{Ignatyev-Penkov-Wolf} M. V. Ignatyev, I. Penkov, and J. A. Wolf: Real group orbits on flag ind-varieties of $\mathrm{SL}(\infty,\mathbb{C})$. In: Lie Theory and Its Applications in Physics, pp. 111--135. Springer Proc. Math. Stat., vol. 191, Springer, Singapore, 2016.
\bibitem{Jantzen}  J. C. Jantzen: Nilpotent orbits in representation theory. Lie theory, 1--211, Progr. Math., vol. 228, Birkh\"auser Boston, Boston, MA, 2004.
\bibitem{Matsuki} T. Matsuki: The orbits of affine symetric spaces under the action of minimal parabolic
subgroups. J. Math. Soc. Japan. \textbf{31} (1979) 331--357.
%\bibitem{Matsuki2} T. Matsuki: Orbits on affine symetric spaces under the action of parabolic subgroups. Hiroshima Math. Jour. \textbf{12} (1982) 307--320.
\bibitem{Matsuki3} T. Matsuki: Closure relations for orbits on affine symmetric spaces under the action of parabolic subgroups. Intersections of associated orbits. Hiroshima Math. J. \textbf{18} (1988) 59--67.
\bibitem{MO} T. Matsuki and T. Oshima:
Embeddings of discrete series into principal series.
In: The orbit method in representation theory (Copenhagen, 1988), pp. 147--175.
Progr. Math., vol. 82, Birkh\"auser, Boston, MA, 1990.
\bibitem{Nishiyama} K. Nishiyama: Enhanced orbit embedding.
Comment. Math. Univ. St. Pauli \textbf{63} (2014), 223--232.
\bibitem{OV} A. L. Onishchik and \`E. B. Vinberg:
Lie groups and algebraic groups.
Translated from the Russian and with a preface by D. A. Leites. Springer Series in Soviet Mathematics. Springer-Verlag, Berlin, 1990.
\bibitem{Otha1} T. Ohta: The closures of nilpotent orbits in the classical symmetric pairs and their singularities. Tohoku Math. J. \textbf{43} (1991), 161--211.
\bibitem{Otha} T. Ohta: An inclusion between sets of orbits and surjectivity of the restriction map of rings of invariants.
Hokkaido Math. J. \textbf{37} (2008), 437--454.
\bibitem{Richardson-Springer} R. W. Richardson and T. A. Springer:
The Bruhat order on symmetric varieties. Geom. Dedicata \textbf{35} (1990), 389--436.
\bibitem{W1} J. A. Wolf: The action of a real semisimple Lie group on a complex flag manifold. I: Orbit structure and holomorphic arc components.  Bull. Amer. Math. Soc. \textbf{75} (1969), 1121--1237.
\bibitem{W} J. A. Wolf: Cycle spaces of infinite dimensional flag domains. Annals of Global Analysis and Geometry \textbf{50} (2016), 315--346.
\bibitem{Y} A. Yamamoto:
Orbits in the flag variety and images of the moment map for classical groups I.
Represent. Theory \textbf{1} (1997), 329--404.
\end{thebibliography}

\end{document}











